% !TeX program = xelatex
% !TeX encoding = UTF-8
\documentclass[11pt,fontset=fandol,scheme=plain]{ctexart}
\usepackage[margin=1in]{geometry}
\usepackage{amsmath,amssymb,amsthm}
\usepackage{xcolor}
\usepackage{enumitem}
% 排版补充：用于第 1.5 节原稿中的交换图。
\usepackage{tikz-cd}

\definecolor{notepurple}{RGB}{91,45,145}
\setlength{\parindent}{0pt}
\setlength{\parskip}{0.65\baselineskip}
\setlist[enumerate]{itemsep=0.3em,topsep=0.3em}

\newtheoremstyle{notes}
  {0.8\baselineskip}{0.5\baselineskip}
  {\normalfont}{}% body font, indent
  {\bfseries}{.}{0.5em}{}
\theoremstyle{notes}
\newtheorem{theorem}{Theorem}[subsection]
\newtheorem{proposition}[theorem]{Proposition}
\newtheorem{definition}[theorem]{Definition}
\newtheorem{example}[theorem]{Example}
\newtheorem{corollary}[theorem]{Corollary}
% 排版补充：第 1.6 节的两个引理与定理等环境共用编号。
\newtheorem{lemma}[theorem]{Lemma}

\newenvironment{purplenote}
  {\par\begingroup\color{notepurple}\small\noindent}
  {\par\endgroup}

% 带编号的紫色补充命题与证明：用于对应手写稿中的紫色补充结论。
\newenvironment{purpleproposition}
  {\begingroup\color{notepurple}\small\begin{proposition}}
  {\end{proposition}\endgroup}
\newenvironment{purpleproof}
  {\begingroup\color{notepurple}\small\begin{proof}}
  {\end{proof}\endgroup}

\usepackage[hidelinks,unicode]{hyperref}
% 交叉引用：每条带编号的结论在环境开头写唯一的 \label{类型:含义}；
% 正文用“定理~\ref{thm:...}”引用编号，或用 \hyperref[标签]{名称} 链接名称。
\newcommand{\unclear}[1]{\text{\color{red}[待确认：#1]}}

\title{Hungerford《代数》自学笔记}
\author{}
\date{}

\begin{document}
\maketitle

\section{群论}

% 各小节正文分别维护在 sections 目录中。
% !TeX root = ../main.tex
\subsection{半群、幺半群和群}

\begin{definition}[半群、幺半群与群]\label{def:semigroup-monoid-group}
设 $G$ 是一个非空集合。
\begin{enumerate}[label=(\arabic*)]
  \item 若在 $G$ 上定义了一个二元运算，记作 $(a,b)\mapsto ab$，且满足结合律
  \[
    a(bc)=(ab)c,\qquad \forall a,b,c\in G,
  \]
  则称 $G$ 为一个\textbf{半群}（semigroup）。

  \item 若 $G$ 是一个半群，并且存在 $e\in G$，使得
  \[
    ae=ea=a,\qquad \forall a\in G,
  \]
  则称 $G$ 为一个\textbf{幺半群}（monoid），$e$ 称为 $G$ 的\textbf{幺元}（单位元）。

  \item 若 $G$ 是一个幺半群，并且对每个 $a\in G$ 都存在 $a^{-1}\in G$，使得
  \[
    aa^{-1}=a^{-1}a=e,
  \]
  则称 $G$ 为一个\textbf{群}，$a^{-1}$ 称为 $a$ 在 $G$ 中的\textbf{逆元}。

  \item 若群 $G$ 还满足交换律
  \[
    ab=ba,\qquad \forall a,b\in G,
  \]
  则称 $G$ 为\textbf{阿贝尔群}（交换群）。
\end{enumerate}
\end{definition}

\begin{purplenote}
对于非空集合上的二元运算，常见的记法有乘法形式和加法形式。具体的半群、幺半群或群只需满足相应的运算律即可；在加法记号中，幺元通常记为 $0$，逆元记为 $-a$。
\end{purplenote}

\begin{theorem}\label{thm:basic-group-properties}
若 $G$ 是一个幺半群，则幺元 $e$ 是唯一的。若 $G$ 是一个群，则还有：
\begin{enumerate}[label=(\arabic*)]
  \item 若 $c\in G$ 且 $cc=c$，则 $c=e$；
  \item 左、右消去律成立：对任意 $a,b,c\in G$，
  \[
    ab=ac\Longrightarrow b=c,
    \qquad
    ba=ca\Longrightarrow b=c;
  \]
  \item 每个 $a\in G$ 的逆元 $a^{-1}$ 是唯一的；
  \item $(a^{-1})^{-1}=a$；
  \item $(ab)^{-1}=b^{-1}a^{-1}$；
  \item 对任意 $a,b\in G$，方程 $ax=b$ 与 $ya=b$ 在 $G$ 中分别有唯一解
  \[
    x=a^{-1}b,\qquad y=ba^{-1}.
  \]
\end{enumerate}
\end{theorem}

\begin{proof}
若 $e$ 与 $e'$ 都是幺元，则 $e=ee'=e'$，故幺元唯一。

若 $cc=c$，在等式两边左乘 $c^{-1}$，即得 $c=e$。消去律同样由在等式两边乘相应的逆元得到。

若 $b,c$ 都是 $a$ 的逆元，则
\[
  b=be=b(ac)=(ba)c=ec=c,
\]
故逆元唯一。由于 $a$ 是 $a^{-1}$ 的逆元，便有 $(a^{-1})^{-1}=a$。又因为
\[
  (ab)(b^{-1}a^{-1})=e=(b^{-1}a^{-1})(ab),
\]
所以 $(ab)^{-1}=b^{-1}a^{-1}$。最后，将 $ax=b$ 左乘 $a^{-1}$、将 $ya=b$ 右乘 $a^{-1}$，分别得到所述解；唯一性由消去律保证。
\end{proof}

\begin{purplenote}
在加法记号中，幺元记为 $0$，逆元记为 $-a$；于是 $(a+b)^{-1}=b^{-1}+a^{-1}$ 写成
\[
  -(a+b)=(-b)+(-a).
\]
只有在交换群中，才能进一步写成 $-(a+b)=-a-b$ 而不必在意次序。
\end{purplenote}

由定理~\ref{thm:basic-group-properties}，群的定义还可以适当减弱。

\begin{proposition}\label{prop:left-group-axioms}
设 $G$ 是一个半群，则 $G$ 是群当且仅当满足下列两个条件：
\begin{enumerate}[label=(\arabic*)]
  \item 存在 $e\in G$，使得 $ea=a$ 对所有 $a\in G$ 成立（左幺元）；
  \item 对每个 $a\in G$，存在 $a'\in G$，使得 $a'a=e$（相对于该左幺元的左逆元）。
\end{enumerate}
\end{proposition}

\begin{purplenote}
相应地，也可用“右幺元＋每个元素都有右逆元”刻画群。
\end{purplenote}

\begin{proof}
% 修正说明：原手写稿此处有被划去的推导；以下保留其思路并改写为完整、无循环依赖的证明。
必要性显然。下证充分性。固定 $a\in G$，取 $a'\in G$ 使 $a'a=e$；再取 $b\in G$ 使 $ba'=e$。于是
\[
  aa'=e(aa')=(ba')(aa')=b(a'a)a'=bea'=ba'=e.
\]
因此 $a'$ 也是 $a$ 的右逆元。进而
\[
  ae=a(a'a)=(aa')a=ea=a,
\]
所以 $e$ 也是右幺元。故 $e$ 是双侧幺元，且每个元素都有双侧逆元，$G$ 是群。
\end{proof}

% 修正说明：原手写稿 1.1.4 是用方程可解性刻画群；先前误录为“反向群”，现已按原稿改回。
\begin{proposition}\label{prop:solvability-group}
设 $G$ 是一个半群，则 $G$ 是群当且仅当对任意 $a,b\in G$，方程
\[
  ax=b,\qquad ya=b
\]
在 $G$ 中都有解。
\end{proposition}

\begin{proof}
必要性由定理~\ref{thm:basic-group-properties} 立即得到。下证充分性。

固定 $a\in G$。由方程 $ya=a$ 有解，可取 $e\in G$ 使 $ea=a$。任取 $b\in G$，由方程 $ax=b$ 有解，可取 $x\in G$ 使 $ax=b$，于是
\[
  eb=e(ax)=(ea)x=ax=b.
\]
因此 $e$ 是 $G$ 的左幺元。再对任意 $a\in G$，由方程 $ya=e$ 有解，可取 $a'\in G$ 使 $a'a=e$，即每个元素都有左逆元。由命题~\ref{prop:left-group-axioms}，$G$ 是群。
\end{proof}

\begin{example}[置换群]\label{ex:permutation-group}
设 $S$ 是非空集合，以 $A(S)$ 表示所有双射 $S\to S$ 构成的集合。以映射复合作为乘法，则 $A(S)$ 是一个群，称为 $S$ 上的\textbf{置换群}。若
\[
  S=\{1,\ldots,n\},
\]
则记 $A(S)=S_n$，称为 $n$ 元对称群。
\end{example}

\begin{purplenote}
对称群 $S_n$ 在有限群论中具有重要作用。
\end{purplenote}

\begin{definition}[群的直积]\label{def:group-direct-product}
设 $G,H$ 都是群，且 $e_G,e_H$ 分别为其幺元。在笛卡尔积
\[
  G\times H=\{(a,b)\mid a\in G, b\in H\}
\]
上逐分量定义
\[
  (a,b)(a',b')=(aa',bb').
\]
则 $G\times H$ 是群，称为 $G$ 与 $H$ 的\textbf{直积}。当群采用加法记号时，直积也常记作 $G\oplus H$。其幺元为 $(e_G,e_H)$，且
\[
  (a,b)^{-1}=(a^{-1},b^{-1}).
\]
\end{definition}

% 修正说明：原笔记先讨论幺半群的商，并将群作为进一步的特殊情形；先前版本误将条件直接限制为群。
\begin{theorem}[商幺半群与商群的构造]\label{thm:quotient-congruence}
设 $\sim$ 是幺半群 $G$ 上的等价关系，并且与乘法相容，即
\[
  a_1\sim a_2,\qquad b_1\sim b_2
  \quad\Longrightarrow\quad
  a_1b_1\sim a_2b_2.
\]
在等价类集合 $G/{\sim}$ 上定义
\[
  \overline a\,\overline b:=\overline{ab}.
\]
则该运算良定义，并使 $G/{\sim}$ 成为幺半群，其幺元为 $\overline e$。

进一步，若 $G$ 是群，则 $G/{\sim}$ 也是群，且
\[
  (\overline a)^{-1}=\overline{a^{-1}}.
\]
若 $G$ 是阿贝尔群，则 $G/{\sim}$ 也是阿贝尔群。
\end{theorem}

\begin{purplenote}
满足上述相容条件的等价关系称为幺半群 $G$ 上的一个\textbf{同余关系}。当 $G$ 是群时，相应的商幺半群就是商群。例如，对固定的正整数 $m$，整数上的模 $m$ 同余给出商群
\[
  \mathbb Z_m.
\]
\end{purplenote}

\begin{example}\label{ex:rationals-mod-integers}
% 修正说明：原手写稿 1.1.8 的底集是全体有理数，而非整数；商群应为 \mathbb Q/\mathbb Z。
在全体有理数 $\mathbb Q$ 上定义等价关系
\[
  a\sim b\quad\Longleftrightarrow\quad a-b\in\mathbb Z.
\]
则等价类之和由
\[
  \overline a+\overline b=\overline{a+b}
\]
定义，所得商群为 $\mathbb Q/\mathbb Z$。
\end{example}

\begin{purplenote}
在半群中，由结合律可无歧义地记
\[
  a_1a_2\cdots a_n=\prod_{i=1}^n a_i,
  \qquad
  a^m=\underbrace{aa\cdots a}_{m\text{ 个}}.
\]
\end{purplenote}

% !TeX root = ../main.tex
\subsection{同态与子群}

\begin{definition}[同态]\label{def:homomorphism}
设 $G,H$ 是两个半群。若映射 $f:G\to H$ 保持运算，即
\[
  f(ab)=f(a)f(b),\qquad \forall a,b\in G,
\]
则称 $f$ 为一个\textbf{同态}（homomorphism）。单射同态称为\textbf{单同态}（monomorphism），满射同态称为\textbf{满同态}（epimorphism），双射同态称为\textbf{同构}（isomorphism）。若 $G$ 与 $H$ 之间存在同构，则记作 $G\cong H$。

同态 $f:G\to G$ 称为 $G$ 的\textbf{自同态}（endomorphism）。

同构 $f:G\to G$ 称为 $G$ 的\textbf{自同构}（automorphism）。
\end{definition}

\begin{purplenote}
同态的本质是保持代数运算；例如采用加法记号时，同态条件写作
\[
  f(a+b)=f(a)+f(b).
\]
\end{purplenote}

\begin{proposition}\label{prop:hom-basic}
\begin{enumerate}[label=(\arabic*)]
  \item 若 $f:G\to H$ 与 $g:H\to K$ 是半群同态，则 $g\circ f:G\to K$ 也是半群同态。
  \item 单同态的复合仍是单同态，满同态的复合仍是满同态，同构的复合仍是同构，自同构的复合仍是自同构。
  \item 若 $f:G\to H$ 是群同态，则
  \[
    f(e_G)=e_H,\qquad f(a^{-1})=f(a)^{-1}\quad(a\in G).
  \]
\end{enumerate}
\end{proposition}

\begin{purplenote}
% 修正说明：群同态必保持幺元；一般幺半群同态若定义中不要求保幺，则未必保持幺元。并规范了习题引用的空格。
第 (3) 条使用了 $G,H$ 是群这一条件；对一般幺半群同态，若未额外要求保持幺元，则 $f(e_G)=e_H$ 未必成立。反例见 Chapter~1, Exercise~2.1。
\end{purplenote}

\begin{definition}[核、像与原像]\label{def:kernel-image-preimage}
设 $f:G\to H$ 是群同态。定义
\[
  \ker f=\{a\in G\mid f(a)=e_H\}.
\]
若 $A\subseteq G$，定义 $A$ 的像
\[
  f(A)=\{f(a)\mid a\in A\},
\]
特别地，$f(G)$ 记为 $\operatorname{im}f$。若 $B\subseteq H$，定义 $B$ 的原像
\[
  f^{-1}(B)=\{a\in G\mid f(a)\in B\}.
\]
\end{definition}

\begin{theorem}\label{thm:mono-kernel}
设 $f:G\to H$ 是群同态。
\begin{enumerate}[label=(\arabic*)]
  \item $f$ 是单同态当且仅当 $\ker f=\{e_G\}$；
  \item $f$ 是同构当且仅当存在群同态 $f^{-1}:H\to G$，使
  \[
    f\circ f^{-1}=1_H,\qquad f^{-1}\circ f=1_G.
  \]
\end{enumerate}
\end{theorem}

\begin{definition}[对运算封闭]\label{def:closed-under-operation}
设 $G$ 是半群，$H\subseteq G$。若对任意 $a,b\in H$ 都有 $ab\in H$，则称 $H$ 对 $G$ 的运算\textbf{封闭}。
\end{definition}

\begin{definition}[子群]\label{def:subgroup}
设 $G$ 是群，$H$ 是 $G$ 的非空子集。若 $H$ 在 $G$ 的二元运算下本身构成群，则称 $H$ 是 $G$ 的\textbf{子群}，记作 $H\leq G$。
\end{definition}

\begin{purplenote}
% 补充说明：按 Hungerford《Algebra》第一章第 2 节（印刷页码 32）的约定区分 trivial subgroup 与 proper subgroup；后者同时排除幺元子群和全群。
群 $G$ 本身以及只含幺元的 $\{e_G\}=\langle e_G\rangle$ 都是 $G$ 的子群。其中 $\{e_G\}$ 称为\textbf{平凡子群}（trivial subgroup）；非平凡子群是指不等于 $\{e_G\}$ 的子群。

特别注意 Hungerford 的术语约定：\textbf{proper subgroup} 指同时满足
\[
  H\leq G,\qquad H\neq\{e_G\},\qquad H\neq G
\]
的子群，即\textbf{非平凡真子群}，也可写为
\[
  \{e_G\}\subsetneq H\subsetneq G.
\]
其他一些教材将 proper subgroup 仅用于表示 $H\leq G$ 且 $H\neq G$，不排除平凡子群。为避免混淆，本笔记在对应 Hungerford 的 proper subgroup 时明确写作“非平凡真子群”；若仅写“真子群”，则只表示 $H\neq G$。记号 $H\leq G$ 本身不排除 $H=\{e_G\}$ 或 $H=G$。

若 $H\leq G$，则 $H$ 的幺元就是 $e_G$，且 $a\in H$ 在 $H$ 中的逆元就是它在 $G$ 中的逆元。
\end{purplenote}

\begin{theorem}\label{thm:image-preimage-subgroups}
设 $f:G\to H$ 是群同态，$A\leq G$，$B\leq H$，则
\[
  \ker f\leq G,\qquad f^{-1}(B)\leq G,\qquad f(A)\leq H.
\]
\end{theorem}

\begin{theorem}[子群判别法]\label{thm:subgroup-test}
设 $H$ 是群 $G$ 的非空子集，则
\[
  H\leq G
  \quad\Longleftrightarrow\quad
  ab^{-1}\in H\quad\text{对所有 }a,b\in H.
\]
\end{theorem}

\begin{corollary}\label{cor:intersection-subgroups}
任意一族子群的交仍是子群。即若 $H_i\leq G$（$i\in I$），则
\[
  \bigcap_{i\in I}H_i\leq G.
\]
\end{corollary}

\begin{definition}[生成子群]\label{def:generated-subgroup}
设 $G$ 是群，$X\subseteq G$。所有包含 $X$ 的G的子群之交称为由 $X$ \textbf{生成的子群}，记作 $\langle X\rangle$。若 $X=\{a_1,\ldots,a_n\}$，也写作 $\langle a_1,\ldots,a_n\rangle$；若 $X=\{a\}$，则称 $\langle a\rangle$ 为由 $a$ 生成的\textbf{循环子群}。

% 修正说明：原手写稿 1.2.10 还定义了生成元、有限生成群和循环群；先前转录遗漏了这些内容。
$X$ 中的元素称为 $\langle X\rangle$ 的\textbf{生成元}。若存在有限多个元素 $a_1,\ldots,a_n\in G$，使得
\[
  G=\langle a_1,\ldots,a_n\rangle,
\]
则称 $G$ 为\textbf{有限生成群}。若存在 $a\in G$ 使得
\[
  G=\langle a\rangle,
\]
则称 $G$ 为由 $a$ 生成的\textbf{循环群}。
\end{definition}

\begin{purplenote}
% 修正说明：补回原手写稿中关于不同生成集及子群之并所生成子群的记号。
不同的生成集可能生成同一个群，即使 $X\neq Y$，也可能有
\[
  \langle X\rangle=\langle Y\rangle.
\]
例如，$\langle x\rangle=\langle x^{-1}\rangle$。一般而言，若 $H,K\leq G$，则 $H\cup K$ 未必是 $G$ 的子群；由 $H\cup K$ 生成的子群记作
\[
  H\vee K:=\langle H\cup K\rangle.
\]
\end{purplenote}

\begin{theorem}\label{thm:generated-subgroup-words}
设 $G$ 是群，$X$ 是 $G$ 的非空子集，则
\[
  \langle X\rangle
  % 修正说明：原转录漏写了集合条件符号 \mid 的反斜杠。
  =\bigl\{a_1^{n_1}\cdots a_t^{n_t}\mid
  t\geq1,\ a_i\in X,\ n_i\in\mathbb Z\bigr\}.
\]
特别地，对任意 $a\in G$，
\[
  \langle a\rangle=\{a^n\mid n\in\mathbb Z\}.
\]
\end{theorem}

% 修正说明：原手写稿 1.2.12 分为两项，并把若干具体自同构群同构作为命题的一部分；先前版本误将它们全部移入紫色备注。
\begin{proposition}[自同构群]\label{prop:automorphism-group}
设 $G$ 是一个群，以 $\operatorname{Aut}G$ 表示 $G$ 的所有自同构，则：
\begin{enumerate}[label=(\arabic*)]
  \item $\operatorname{Aut}G$ 在映射复合下构成一个群；
  \item 有
  \[
    \operatorname{Aut}(\mathbb Z)\cong \mathbb Z_2,
    \qquad
    \operatorname{Aut}(\mathbb Z_6)\cong \mathbb Z_2,
  \]
  \[
    \operatorname{Aut}(\mathbb Z_8)
    \cong \mathbb Z_2\oplus\mathbb Z_2,
  \]
  且对任意素数 $p$，
  \[
    \operatorname{Aut}(\mathbb Z_p)
    \cong \mathbb Z_{p-1}.
  \]
\end{enumerate}
\end{proposition}

\begin{purplenote}
% 修正说明：将手写稿中尚未写完的猜想补充为一般结论，并说明自同构的具体形式。
% 补充说明：详细解释单位群 (\mathbb Z_n)^\times 的定义及其与自同构群的对应关系。
% 补充说明：明确区分 \mathbb Z_n 的加法群结构与单位群上的乘法结构。
这里 $\mathbb Z_n$ 表示模 $n$ 的剩余类集合。单独将它视为群时，本笔记默认指其\textbf{加法群}
\[
  (\mathbb Z_n,+).
\]
与此不同，下面的单位群采用模 $n$ 乘法作为二元运算。

对任意正整数 $n$，称 $\overline k\in\mathbb Z_n$ 为一个\textbf{可逆元素}（或\textbf{单位}），若存在 $\overline l\in\mathbb Z_n$，使得
\[
  \overline k\,\overline l=\overline 1.
\]
这等价于存在整数 $l$ 使 $kl\equiv1\pmod n$，也等价于
\[
  \gcd(k,n)=1.
\]
所有可逆元素在模 $n$ 乘法下构成一个群，称为 $\mathbb Z_n$ 的\textbf{单位群}，即
\[
  \bigl((\mathbb Z_n)^{\times},\cdot\bigr),
  \qquad
  (\mathbb Z_n)^{\times}
  =\{\overline k\in\mathbb Z_n\mid \gcd(k,n)=1\}.
\]
其幺元是 $\overline1$，而 $\overline k$ 的逆元是满足 $kl\equiv1\pmod n$ 的 $\overline l$；该群共有 $\varphi(n)$ 个元素。

另一方面，加法循环群 $(\mathbb Z_n,+)$ 的每个自同构都由 $\overline1$ 的像唯一决定。对 $\overline k\in(\mathbb Z_n)^{\times}$，定义
\[
  \varphi_k(\overline a)=\overline{ka}\qquad(\overline a\in\mathbb Z_n),
\]
则 $\varphi_k$ 是加法群 $(\mathbb Z_n,+)$ 的自同构，并且每个自同构都以这种方式唯一得到。由于
\[
  \varphi_k\circ\varphi_l=\varphi_{kl},
\]
故
\[
  \operatorname{Aut}(\mathbb Z_n,+)
  \cong\bigl((\mathbb Z_n)^{\times},\cdot\bigr).
\]

% 补充说明：群的阶为 \varphi(n) 并不意味着该群一定是阶为 \varphi(n) 的循环群。
需要特别注意，虽然
\[
  \bigl|\operatorname{Aut}(\mathbb Z_n,+)\bigr|=\varphi(n),
\]
但 $(\mathbb Z_n)^{\times}$ 不一定是循环群，因而一般不能断言
\[
  \operatorname{Aut}(\mathbb Z_n,+)\cong\mathbb Z_{\varphi(n)}.
\]
% 修正说明：删除 \varphi 前误录的多余反斜杠。
例如，$\varphi(8)=4$，但
\[
  \operatorname{Aut}(\mathbb Z_8,+)
  \cong(\mathbb Z_8)^{\times}
  \cong\mathbb Z_2\oplus\mathbb Z_2
  \not\cong\mathbb Z_4.
\]
\end{purplenote}

% !TeX root = ../main.tex
\subsection{循环群}

\begin{theorem}\label{thm:subgroups-of-integers}
% 修正说明：原手写稿 1.3.1 区分零子群与非零子群，并说明非零子群含有无限多个元素；先前误写成“每个非零子群都有有限指数”。
加法群 $\mathbb Z$ 的每个子群 $H$ 都是循环群。具体地，
\[
  H=\langle0\rangle,
\]
或者
\[
  H=\langle m\rangle=m\mathbb Z,
\]
其中 $m$ 是 $H$ 中最小的正整数。若 $H\neq\langle0\rangle$，则 $H$ 含有无限多个元素。
\end{theorem}

\begin{theorem}\label{thm:cyclic-classification}
任意无限循环群都与加法群 $\mathbb Z$ 同构；任意阶为 $m$ 的有限循环群都与 $\mathbb Z_m$ 同构。
\end{theorem}

\begin{definition}[元素的阶]\label{def:element-order}
设 $G$ 是群，$a\in G$。循环子群 $\langle a\rangle$ 的阶称为元素 $a$ 的\textbf{阶}，记作 $|a|$。
\end{definition}

\begin{theorem}\label{thm:order-powers}
设 $G$ 是群，$a\in G$。

若 $|a|=\infty$，则
\begin{enumerate}[label=(\arabic*)]
  \item $a^k=e$ 当且仅当 $k=0$；
  \item $a^r=a^s$ 当且仅当 $r=s$。
\end{enumerate}

若 $|a|=m<\infty$，则
\begin{enumerate}[label=(\arabic*),resume]
  \item $m$ 是满足 $a^m=e$ 的最小正整数；
  \item $a^k=e$ 当且仅当 $m\mid k$；
  \item $a^r=a^s$ 当且仅当 $r\equiv s\pmod m$；
  \item $\langle a\rangle=\{e,a,a^2,\ldots,a^{m-1}\}$，且这些元素两两不同；
  \item 对任意 $k\in\mathbb Z$，
  % 修正说明：手写稿给出了 k\mid m 时的特例；这里写成对任意整数 k 成立的完整公式。
  \[
    |a^k|=\frac{m}{\gcd(m,k)}.
  \]
\end{enumerate}
\end{theorem}


\begin{purplenote}
当 $k\mid m$ 时，上式简化为 $|a^k|=m/k$（取 $k>0$）。
\end{purplenote}

\begin{theorem}\label{thm:cyclic-images-subgroups}
% 修正说明：原手写稿 1.3.5 同时包含“同态像”和“子群”两部分；先前转录遗漏了前一部分。
循环群的每一个同态像都是循环群；循环群的每一个子群也都是循环群。更具体地，若 $G$ 是循环群且 $f:G\to H$ 是群同态，则 $\operatorname{im}f$ 是循环群；若 $K\leq G$，则 $K$ 是循环群。
\end{theorem}

\begin{theorem}[循环群的生成元]\label{thm:cyclic-generators}
设 $G=\langle a\rangle$。
\begin{enumerate}[label=(\arabic*)]
  \item 若 $G$ 为无限循环群，则它的生成元只有 $a$ 与 $a^{-1}$；
  \item 若 $|G|=m<\infty$，则 $a^k$ 是 $G$ 的生成元当且仅当 $\gcd(m,k)=1$。
\end{enumerate}
\end{theorem}

% 第 8 张照片：以下均为循环群一节的补充结论，故置于 1.3 节末尾。
% 修正说明：照片中的补充命题主体为黑笔，现改用普通 proposition；只有原稿紫笔注释继续使用紫色环境。
\begin{proposition}\label{prop:abelian-lcm-order}
设 $G$ 是阿贝尔群，$a,b\in G$ 的阶分别为有限数 $m,n$。则 $G$ 中存在元素 $c$，使得
\[
  |c|=\operatorname{lcm}(m,n).
\]
\end{proposition}

\begin{proof}
将 $\operatorname{lcm}(m,n)$ 分解为互异素数幂的乘积
\[
  \operatorname{lcm}(m,n)=\prod_{i=1}^r p_i^{\gamma_i}.
\]
对每个 $p_i$，它在 $m,n$ 中的指数至少有一个等于 $\gamma_i$。若 $p_i^{\gamma_i}$ 是整除 $m$ 的最高 $p_i$-幂，则取
\[
  x_i=a^{m/p_i^{\gamma_i}};
\]
否则取 $x_i=b^{n/p_i^{\gamma_i}}$。由定理~\ref{thm:order-powers}，$|x_i|=p_i^{\gamma_i}$。由于 $G$ 为阿贝尔群，$x_1,\ldots,x_r$ 两两交换，且它们的阶两两互素，故
\[
  c=x_1x_2\cdots x_r
\]
的阶为 $\prod_i p_i^{\gamma_i}=\operatorname{lcm}(m,n)$。
\end{proof}

\begin{purplenote}
% 修正说明：存在性结论是正确的；仅“总可直接取 c=ab”不正确。当 gcd(m,n)=1 时可以取 c=ab。
上述 $c$ 一般不能直接取为 $ab$；若 $\gcd(m,n)=1$，则确有 $|ab|=mn$。
\end{purplenote}

\begin{proposition}\label{prop:prufer-subgroups}
\textbf{Prüfer $p$-群的子群。}
设 $p$ 为素数，定义
\[
  \mathbb Z(p^\infty)
  =\left\{\overline{\frac{a}{b}}\in\mathbb Q/\mathbb Z
  \ \middle|\ a,b\in\mathbb Z,\ b=p^i,\ i\geq0\right\}.
\]
设 $H\leq\mathbb Z(p^\infty)$，则：
\begin{enumerate}[label=(\arabic*)]
  \item 对任意 $x\in\mathbb Z(p^\infty)$，存在 $n\geq0$，使得
  \[
    |x|=p^n.
  \]

  \item 若 $H$ 中至少有一个元素的阶为 $p^k$，并且 $H$ 中没有元素的阶大于 $p^k$，则
  \[
    H=\left\langle\overline{\frac{1}{p^k}}\right\rangle
    \cong\mathbb Z_{p^k}.
  \]

  \item 若 $H$ 中元素的阶没有上界，则
  \[
    H=\mathbb Z(p^\infty).
  \]

  % 修正说明：原稿写“非平凡子群”，但整个 Prüfer p-群本身也是非平凡子群；这里严谨地改为“非平凡真子群”。
  \item $\mathbb Z(p^\infty)$ 的非平凡真子群恰为
  \[
    C_n:=\left\langle\overline{\frac{1}{p^n}}\right\rangle
    \qquad(n=1,2,\ldots).
  \]
  若记 $C_0=\langle\overline0\rangle$，则
  \[
    C_0<C_1<C_2<\cdots<C_n<\cdots<\mathbb Z(p^\infty).
  \]

  % 修正说明：补回原手写稿 1.3.8 的第 (5) 条。
  \item 设 $G$ 是阿贝尔群，且 $x_1,x_2,\ldots\in G$ 满足
  \[
    |x_1|=p,\qquad px_{n+1}=x_n\quad(n\geq1).
  \]
  则
  \[
    \langle x_1,x_2,\ldots\rangle\cong\mathbb Z(p^\infty).
  \]
\end{enumerate}
\end{proposition}

\begin{purplenote}
由上述结构可知，$\mathbb Z(p^\infty)$ 含有无限多个有限阶元素；更精确地，阶为 $p^n$ 的元素恰有 $\varphi(p^n)$ 个。
\end{purplenote}

\begin{proposition}\label{prop:finitely-many-subgroups}
若一个群只有有限多个子群，则它必为有限群。
\end{proposition}

\begin{purpleproof}
若某元素生成无限循环子群，则该循环子群本身已有无限多个子群，矛盾。因此每个循环子群都有限。群是其所有循环子群之并，而循环子群只有有限多个，故该群有限。
\end{purpleproof}

\begin{proposition}\label{prop:torsion-subgroup}
若 $G$ 是阿贝尔群，则所有有限阶元素组成的集合
\[
  T=\{a\in G\mid |a|<\infty\}
\]
是 $G$ 的子群，称为 $G$ 的挠子群。
\end{proposition}

\begin{proof}
若 $|a|=m$、$|b|=n$，取 $N=\operatorname{lcm}(m,n)$。因 $G$ 交换，
\[
  (ab^{-1})^N=a^N b^{-N}=e,
\]
故 $ab^{-1}\in T$。由\hyperref[thm:subgroup-test]{子群判别法}，$T\leq G$。
\end{proof}

\begin{purplenote}
阿贝尔条件不可省略：一般群中，两个有限阶元素的乘积未必仍有有限阶。
\end{purplenote}

\begin{proposition}\label{prop:infinite-cyclic-subgroups}
% 修正说明：重新核查后，原手写稿中的“当且仅当”是正确的；先前删除逆命题属于误修正。
设 $G$ 是无限群，则 $G$ 是循环群当且仅当 $G$ 的每一个非平凡子群都与 $G$ 同构。
\end{proposition}

\begin{proof}
若 $G$ 是无限循环群，则其每个非平凡子群仍是无限循环群，因而与 $G$ 同构。

反之，任取 $a\in G\setminus\{e\}$，则 $\langle a\rangle$ 是 $G$ 的非平凡子群。由假设 $\langle a\rangle\cong G$，而 $\langle a\rangle$ 是循环群，故 $G$ 也是循环群。
\end{proof}

\begin{purplenote}
正向证明使用了无限循环群的每个非平凡子群仍为无限循环群。实际判断一个无限群是否循环时，若能找到一个与原群不同构的非平凡子群，即可断定原群不是循环群。
\end{purplenote}

% !TeX root = ../main.tex
\subsection{陪集和计数技巧}

\begin{definition}[模子群的同余]\label{def:congruence-mod-subgroup}
设 $G$ 是群，$H\leq G$。对 $a,b\in G$：
\begin{align*}
  a&\equiv_r b\pmod H &&\Longleftrightarrow&& ab^{-1}\in H,\\
  a&\equiv_l b\pmod H &&\Longleftrightarrow&& a^{-1}b\in H.
\end{align*}
分别称 $a,b$ 模 $H$ 右同余、左同余。
\end{definition}

\begin{purplenote}
% 修正说明：补回原手写稿中非阿贝尔情形的记号区别、加法形式以及整数模 m 同余这一特殊情形。
若 $G$ 不是阿贝尔群，模 $H$ 的右同余和左同余一般并不相同，因此必须分别记作
\[
  a\equiv_r b\pmod H,
  \qquad
  a\equiv_l b\pmod H.
\]

若 $G$ 是阿贝尔群，则
\[
  ab^{-1}\in H\quad\Longleftrightarrow\quad a^{-1}b\in H,
\]
所以左、右同余相同，此时可以省略下标 $l,r$。

% 修正说明：一般加法记号下左同余应写为 (-a)+b；只有在阿贝尔群中才能进一步交换为 b-a。
若仅把群运算改用加法记号，而不预设交换律，则右同余条件写作
\[
  a+(-b)\in H,
\]
左同余条件写作
\[
  (-a)+b\in H.
\]
前者通常简写为 $a-b\in H$；后者一般不能改写成 $b-a\in H$。在当前 $G$ 为阿贝尔群的假设下，才有
\[
  (-a)+b=b-a,
\]
并且 $a-b=-(b-a)$。由于 $H$ 对取负元封闭，两种同余条件等价。

通常的整数模 $m$ 同余正是这一概念的特殊情形：取加法群 $G=\mathbb Z$ 及其子群
\[
  H=m\mathbb Z=\langle m\rangle,
\]
则
\[
  a\equiv b\pmod m
  \quad\Longleftrightarrow\quad
  m\mid(a-b)
  \quad\Longleftrightarrow\quad
  a-b\in m\mathbb Z.
\]
\end{purplenote}

\begin{theorem}\label{thm:cosets-equivalence}
若$H$是群$G$的子群则模 $H$ 的左、右同余都是 $G$ 上的等价关系。元素 $a$ 的右同余类和左同余类分别为
\[
  Ha=\{ha\mid h\in H\},
  \qquad
  aH=\{ah\mid h\in H\}.
\]
并且对任意 $a\in G$，
\[
  |Ha|=|H|=|aH|.
\]
\end{theorem}

\begin{definition}[陪集]\label{def:coset}
集合 $Ha$ 称为 $H$ 在 $G$ 中的\textbf{右陪集}，$aH$ 称为 $H$ 在 $G$ 中的\textbf{左陪集}。
\end{definition}

\begin{purplenote}
% 修正说明：补回第 7 张照片中 1.4.3 后的紫色备注，明确左右陪集及其加法记号。
一般不能保证 $Ha=aH$，即左、右陪集可能不同，但并非对所有 $a$ 都有 $Ha\neq aH$。
若将群运算写成加法，则右陪集和左陪集分别为
\[
  H+a=\{h+a\mid h\in H\},
  \qquad
  a+H=\{a+h\mid h\in H\}.
\]
加法记号本身并不意味着交换律；若 $G$ 是阿贝尔群，则 $H+a=a+H$。
\end{purplenote}

\begin{corollary}\label{cor:coset-partition}
% 修正说明：按第 10 张照片恢复 1.4.4 的四条结构及陪集集合记号。
设 $G$ 是群，$H\leq G$，则：
\begin{enumerate}[label=(\arabic*)]
  \item
  \[
    G=\bigcup_{a\in G}Ha=\bigcup_{a\in G}aH,
  \]
  即 $G$ 是 $H$ 的全体右陪集之并，也是全体左陪集之并；
  \item 任意两个右陪集要么不相交，要么相同；任意两个左陪集也要么不相交，要么相同；
  \item 对任意 $a,b\in G$，
  \[
    Ha=Hb\Longleftrightarrow ab^{-1}\in H,
    \qquad
    aH=bH\Longleftrightarrow a^{-1}b\in H;
  \]
  \item 若记
  \[
    \mathcal R=\{Ha\mid a\in G\},
    \qquad
    \mathcal L=\{aH\mid a\in G\},
  \]
  则 $|\mathcal R|=|\mathcal L|$，即右陪集与左陪集的个数相同。
\end{enumerate}
\end{corollary}

\begin{purplenote}
% 修正说明：补回原稿第 (4) 条的紫色证明提示，并解释其良定义性及双射性。
第 (4) 条的证明是建立双射
\[
  \mathcal R\longrightarrow\mathcal L,
  \qquad Ha\longmapsto a^{-1}H.
\]
这是对陪集中的每个元素取逆得到的，因为 $(Ha)^{-1}=a^{-1}H$；再取一次逆即回到原陪集，因此该映射良定义且为双射。
\end{purplenote}

\begin{definition}[指数]\label{def:subgroup-index}
子群 $H\leq G$ 的左（等价地，右）陪集个数称为 $H$ 在 $G$ 中的\textbf{指数}，记作 $[G:H]$。
\end{definition}

\begin{theorem}[指数的乘法公式]\label{thm:index-multiplication}
若 $K\leq H\leq G$，则
\[
  [G:K]=[G:H]\,[H:K].
\]
特别地，只要其中任意两个指数有限，第三个也有限。
\end{theorem}

\begin{corollary}[Lagrange 定理]\label{cor:lagrange}
若 $H\leq G$，则
\[
  |G|=[G:H]|H|.
\]
特别地，当 $G$ 有限时，$|H|\mid |G|$，并且对每个 $a\in G$ 都有 $|a|\mid|G|$。
\end{corollary}

\begin{purplenote}
% 恢复说明：第 10 张照片中 1.4.8 前的紫色计数技巧引言；集合乘积定义适用于任意子集，而不仅是子群。
下面介绍一些计数技巧。设 $G$ 是群，$H,K\subseteq G$，定义集合乘积
\[
  HK=\{hk\mid h\in H,\ k\in K\}.
\]
陪集是集合乘积的特殊情形：右陪集 $Ha=H\{a\}$，左陪集 $aH=\{a\}H$。
需要注意，即使 $H,K\leq G$，集合 $HK$ 也不一定是 $G$ 的子群。
\end{purplenote}

\begin{theorem}[子群乘积的大小]\label{prop:HK-cardinality}
设 $H,K$ 是群 $G$ 的有限子群，则
\[
  |HK|=\frac{|H|\,|K|}{|H\cap K|}.
\]
\end{theorem}

\begin{proof}
% 恢复说明：补回第 9 张照片中 1.4.8 的陪集分解证明，并明确陪集代表元的使用。
令 $C=H\cap K$则$C$也是$G$的子群，取 $C$ 在 $K$ 中的不同右陪集代表元 $k_1,\ldots,k_t$，则
\[
  K=\bigcup_{i=1}^t Ck_i,
  \qquad t=[K:C]=\frac{|K|}{|C|}.
\]
因为 $C\leq H$，所以
\[
  HK=\bigcup_{i=1}^t HCk_i=\bigcup_{i=1}^t Hk_i.
\]
若 $Hk_i\cap Hk_j\neq\varnothing$，则存在 $h_i,h_j\in H$ 使 $h_i k_i=h_j k_j$，于是
\[
  k_i k_j^{-1}=h_i^{-1}h_j\in H\cap K=C.
\]
从而 $Ck_i=Ck_j$，故 $i=j$。因此上述 $Hk_i$ 两两不交，且每个都有 $|H|$ 个元素，故
\[
  |HK|=t|H|=\frac{|H|\,|K|}{|H\cap K|}.
\]
\end{proof}

\begin{proposition}[子群中的相对指数]\label{prop:index-intersection}
% 恢复说明：第 10 张照片中的 1.4.9 是下述较一般的相对指数不等式；先前误用原稿 1.4.10 的结论替换了它。
设 $G$ 是群，$H,K\leq G$，则
\[
  [H:H\cap K]\leq[G:K].
\]
进一步，若 $[G:K]$ 有限，则
\[
  [H:H\cap K]=[G:K]\quad\Longleftrightarrow\quad G=KH.
\]
\end{proposition}

\begin{proof}
% 修正说明：按第 9 张照片使用右陪集集合，避免用商群记号暗示子群正规；保留等号判据所需的有限性条件。
令 $C=H\cap K$，记 $\mathcal A=\{Ch\mid h\in H\}$、$\mathcal B=\{Kg\mid g\in G\}$。考虑映射
\[
  \phi:\mathcal A\longrightarrow\mathcal B,
  \qquad Ch\longmapsto Kh.
\]
若 $Ch_1=Ch_2$，则 $h_1h_2^{-1}\in C\subseteq K$，故 $Kh_1=Kh_2$，所以 $\phi$ 良定义。
反之，若 $Kh_1=Kh_2$，则 $h_1h_2^{-1}\in K$；又因 $h_1,h_2\in H$，故 $h_1h_2^{-1}\in C$，从而 $Ch_1=Ch_2$。因此 $\phi$ 为单射，得到
\[
  [H:H\cap K]\leq[G:K].
\]
当 $[G:K]$ 有限时，两边相等当且仅当此单射 $\phi$ 为满射。
而 $\phi$ 满射等价于每个 $g\in G$ 都能找到 $h\in H$ 使 $Kg=Kh$，即 $g\in Kh$，也就等价于 $G=KH$。
\end{proof}

\begin{proposition}[有限指数子群的交]\label{prop:finite-index-intersection}
% 恢复说明：按第 11 张照片恢复 1.4.10 的有限指数假设、不等式及等号判据；先前新增的有限性逆命题与下界移至命题后的独立补充。
设 $G$ 是群，$H,K\leq G$，且 $[G:H]$ 与 $[G:K]$ 均有限。则 $[G:H\cap K]$ 有限，且
\[
  [G:H\cap K]\leq[G:H][G:K].
\]
进一步，
\[
  [G:H\cap K]=[G:H][G:K]\quad\Longleftrightarrow\quad G=HK.
\]
\end{proposition}

\begin{purplenote}
% 恢复说明：原稿 1.4.10 的紫色证明提示是“利用 1.4.6 和 1.4.9 即得”；以下展开该推导。
由定理~\ref{thm:index-multiplication} 和命题~\ref{prop:index-intersection}，
\[
  [G:H\cap K]=[G:H][H:H\cap K]
  \leq[G:H][G:K]<\infty.
\]
因为 $[G:H]$ 是正的有限整数，等号成立当且仅当 $[H:H\cap K]=[G:K]$，即 $G=KH$。
又由 $(KH)^{-1}=HK$ 及 $G^{-1}=G$，可知 $G=KH$ 与 $G=HK$ 等价。
\end{purplenote}

% 转录补充（非原稿）：下面是先前额外加入的正确推论，并非原稿 1.4.10；不另编号，也不冒充原紫笔批注。
\par\textbf{转录补充（非原稿）。}
先前版本增加的有限性逆命题与下界也成立：
\[
  [G:H\cap K]<\infty
  \quad\Longleftrightarrow\quad
  [G:H]<\infty\ \text{且}\ [G:K]<\infty.
\]
事实上，若左边有限，由 $H\cap K\leq H,K$ 及\hyperref[thm:index-multiplication]{指数乘法公式}即知 $[G:H]$、$[G:K]$ 都有限，且
\[
  \max\{[G:H],[G:K]\}\leq[G:H\cap K].
\]
反向由命题~\ref{prop:finite-index-intersection} 得到。这里是补充推论，不替代原命题。

\begin{proposition}\label{prop:coset-of-intersection}
任意一族子群 $\{H_i\mid i\in I\}$ 满足
\[
  \left(\bigcap_{i\in I}H_i\right)a
  =\bigcap_{i\in I}(H_i a),
  \qquad \forall a\in G.
\]
\end{proposition}

\begin{proposition}\label{prop:prime-order-characterizations}
对有限群 $G$，下列条件等价：
\begin{enumerate}[label=(\arabic*)]
  \item $|G|$ 是素数；
  \item $G\neq\{e\}$，且 $G$ 没有非平凡真子群；
  \item 存在素数 $p$ 使 $G\cong\mathbb Z_p$。
\end{enumerate}
\end{proposition}

\begin{theorem}\label{thm:HK-subgroup}
设 $H,K\leq G$，则
\[
  HK\leq G\quad\Longleftrightarrow\quad HK=KH.
\]
\end{theorem}

\begin{proposition}\label{prop:p-subgroup-product-not-subgroup}
设 $G$是群且$|G|=p^k m$，其中 $p$ 为素数且 $\gcd(p,m)=1$。若 $H\leq G$ 的阶为 $p^k$，$K\leq G$ 的阶为 $p^d$，其中 $0<d\leq k$，且 $K\nsubseteq H$，则 $HK$ 不是 $G$ 的子群。
\end{proposition}

\begin{proof}
% 修正说明：补足 K\nsubseteq H 推出 H\cap K 为 K 的真子群，以及阶的 p-指数矛盾。
若 $HK\leq G$，由命题~\ref{prop:HK-cardinality}，
\[
  |HK|=\frac{p^{k+d}}{|H\cap K|}.
\]
因为 $K\nsubseteq H$，所以 $H\cap K$ 是 $K$ 的真子群，从而 $|H\cap K|\leq p^{d-1}$。于是 $p^{k+1}\mid |HK|$。但 $|HK|\mid |G|=p^k m$ 且 $p\nmid m$，矛盾。
\end{proof}

% 新增：第 1.5 节单独维护，延续原笔记编号与双色排版。
% !TeX root = ../main.tex
\subsection{正规子群、商群和同态}

% 来源：微信图片_20260910201726_17_215.jpg，1.5.1--1.5.6。
% 排版约定：黑笔正文用普通定理与证明环境；原紫笔批注用 purplenote 或 purpleproof。
% 记号说明：\trianglelefteq 表示正规子群，不要求为真子群；\leq 同样允许相等。
\begin{theorem}\label{thm:normal-equivalences}
设 $N$ 是群 $G$ 的子群，则下列条件等价：
\begin{enumerate}[label=(\arabic*)]
  \item 模 $N$ 的左、右同余是 $G$ 上的同一个等价关系，即对任意 $a,b\in G$，
  \[
    a\equiv_l b\pmod N\quad\Longleftrightarrow\quad a\equiv_r b\pmod N;
  \]
  \item $N$ 在 $G$ 中的每一个左陪集都是 $N$ 在 $G$ 中的某个右陪集；
  \item 对任意 $a\in G$，$aN=Na$；
  \item 对任意 $a\in G$，$aNa^{-1}\subseteq N$，其中
  \[
    aNa^{-1}=\{ana^{-1}\mid n\in N\};
  \]
  \item 对任意 $a\in G$，$aNa^{-1}=N$。
\end{enumerate}
\end{theorem}

\begin{purplenote}
第 (4) 条是判别正规子群时常用的条件。
\end{purplenote}

\begin{definition}[正规子群]\label{def:normal-subgroup}
设 $N\leq G$。若 $N$ 满足定理~\ref{thm:normal-equivalences} 中的任意一个条件，则称 $N$ 为 $G$ 的\textbf{正规子群}，记作 $N\trianglelefteq G$。
\end{definition}

\begin{purplenote}
阿贝尔群的任意子群都是其正规子群。
\end{purplenote}

\begin{proposition}\label{prop:normal-intersection-restriction}
设 $N\trianglelefteq G$，并且 $N_i\trianglelefteq G$（$i\in I$），则：
\begin{enumerate}[label=(\arabic*)]
  \item $\displaystyle\bigcap_{i\in I}N_i\trianglelefteq G$；
  \item 若 $N\leq H\leq G$，则 $N\trianglelefteq H$。
\end{enumerate}
\end{proposition}

\begin{purplenote}
需要注意，正规性一般不具有传递性：即使 $N\leq H\leq G$，且 $N\trianglelefteq H$、$H\trianglelefteq G$，也不一定有 $N\trianglelefteq G$。
\end{purplenote}

\begin{theorem}\label{thm:normal-product}
设 $G$ 是群，$K\leq G$，$N\trianglelefteq G$，则：
\begin{enumerate}[label=(\arabic*)]
  \item $N\cap K\trianglelefteq K$；
  \item $N\trianglelefteq N\vee K$；
  \item $NK=N\vee K=KN$；
  \item 若还有 $K\trianglelefteq G$ 且 $K\cap N=\{e\}$，则对任意 $k\in K$、$n\in N$，都有 $nk=kn$。
\end{enumerate}
\end{theorem}

\begin{proof}
% 来源：微信图片_20260910201738_18_215.jpg，上半页的 1.5.4 证明移至本定理后。
(1) 首先 $N\cap K\leq K$。对任意 $a\in K$、$t\in N\cap K$，由 $N\trianglelefteq G$ 得 $ata^{-1}\in N$；又因 $a,t\in K$，有 $ata^{-1}\in K$。故 $ata^{-1}\in N\cap K$，所以 $N\cap K\trianglelefteq K$。

(2) 由 $N\leq N\vee K\leq G$，利用命题~\ref{prop:normal-intersection-restriction} 的第 (2) 条即得。

% 修正说明：原稿第 (3) 条的“两两交换”应理解为利用正规性重写乘积，不是断言 n 与 k 交换；重写后 N 中的元素可能改变。
(3) 显然 $NK\subseteq N\vee K$。由生成子群的描述，$N\vee K$ 的每个元素都可写成有限交错乘积
\[
  n_1k_1\cdots n_rk_r,\qquad n_i\in N,\ k_i\in K.
\]
由于 $N\trianglelefteq G$，对任意 $k\in K$、$n\in N$，可写成
\[
  kn=(knk^{-1})k\in NK.
\]
反复利用这一重写规则，并在各子群内合并乘积，可将上述交错乘积写成 $nk$ 的形式，故 $N\vee K=NK$。
同理，利用 $nk=k(k^{-1}nk)$ 可得 $N\vee K=KN$。

(4) 因为 $K\trianglelefteq G$，所以 $nkn^{-1}\in K$，从而 $nkn^{-1}k^{-1}\in K$。
又由 $N\trianglelefteq G$，有 $kn^{-1}k^{-1}\in N$，故
\[
  nkn^{-1}k^{-1}\in N\cap K=\{e\}.
\]
因此 $nk=kn$。
\end{proof}

\begin{theorem}\label{thm:quotient-group-construction}
设 $N\trianglelefteq G$，以 $G/N$ 表示 $N$ 在 $G$ 中的所有左陪集（也就是所有右陪集）组成的集合。在 $G/N$ 上定义
\[
  (aN)(bN)=(ab)N.
\]
则该运算使 $G/N$ 成为群，且 $|G/N|=[G:N]$。
\end{theorem}

\begin{definition}[商群]\label{def:quotient-group}
设 $N\trianglelefteq G$，则定理~\ref{thm:quotient-group-construction} 中定义的群 $G/N$ 称为\textbf{由 $N$ 决定的 $G$ 的商群}。
\end{definition}

% 来源：微信图片_20260910201741_19_215.jpg 顶部，接续 1.5.6 的紫色备注。
\begin{purplenote}
\begin{enumerate}[label=(\arabic*)]
  \item 若 $G$ 的运算采用加法记号，则商群中的运算相应写为
  \[
    (a+N)+(b+N)=(a+b)+N.
  \]
  \item 对任意整数 $m>1$，$\mathbb Z_m$ 就是由加法群 $\mathbb Z$ 对子群 $\langle m\rangle$ 取商得到的群。
\end{enumerate}
以下讨论正规子群、商群和同态之间的关系。
\end{purplenote}

\begin{theorem}\label{thm:kernel-and-projection}
设 $f:G\to H$ 是群同态，则 $\ker f\trianglelefteq G$。反之，若 $N\trianglelefteq G$，则映射
% 修正说明：原稿 1.5.7 及其紫色备注中将标准投影的定义域误写为 N；应为 G，否则一般不满射到 G/N。
\[
  \pi:G\longrightarrow G/N,\qquad \pi(a)=aN
\]
是满同态，且 $\ker\pi=N$。
\end{theorem}

\begin{proof}
% 来源：照片 18 中的 1.5.7 证明。
已知 $\ker f\leq G$。对任意 $a\in G$、$t\in\ker f$，
\[
  f(ata^{-1})=f(a)f(t)f(a)^{-1}=f(a)e_Hf(a)^{-1}=e_H,
\]
故 $ata^{-1}\in\ker f$，从而 $\ker f\trianglelefteq G$。
映射 $\pi$ 显然为满射，且
\[
  \pi(ab)=abN=(aN)(bN)=\pi(a)\pi(b).
\]
此外，$\pi(a)=N$ 当且仅当 $a\in N$，故 $\ker\pi=N$。
\end{proof}

\begin{purplenote}
上述映射 $\pi:G\to G/N$ 称为\textbf{标准满同态}或\textbf{标准投影}。以后若无特别说明，提到 $G\to G/N$（$N\trianglelefteq G$）时均指这一映射。
\end{purplenote}

\begin{theorem}\label{thm:factor-through-quotient}
设 $f:G\to H$ 是群同态，$N\trianglelefteq G$ 且 $N\leq\ker f$。则存在唯一的群同态
\[
  \overline f:G/N\longrightarrow H,
  \qquad \overline f(aN)=f(a)\quad(a\in G).
\]
并且
\[
  \operatorname{im}\overline f=\operatorname{im}f,
  \qquad \ker\overline f=(\ker f)/N.
\]
$\overline f$ 是同构，当且仅当 $f$ 是满同态且 $N=\ker f$。
\end{theorem}

\begin{proof}
% 来源：照片 18 中的 1.5.8 证明；补明唯一性，统一 \overline f 与 f 的记号。
由 $N\trianglelefteq G$、$N\leq\ker f$，有 $N\trianglelefteq\ker f$。
若 $aN=bN$，由正规性可得 $ab^{-1}\in N\subseteq\ker f$，所以 $f(a)=f(b)$，故 $\overline f$ 良定义。
又有
\[
  \overline f\bigl((aN)(bN)\bigr)=f(ab)=f(a)f(b)
  =\overline f(aN)\,\overline f(bN),
\]
因此 $\overline f$ 是群同态。每个陪集都有形式 $aN$，故指定 $\overline f(aN)=f(a)$ 已唯一确定该映射；像的相等也立即可得。

对任意 $cN\in G/N$，
\[
  cN\in\ker\overline f
  \quad\Longleftrightarrow\quad f(c)=e_H
  \quad\Longleftrightarrow\quad c\in\ker f.
\]
所以 $\ker\overline f=(\ker f)/N$。由群同态的单射判别及像的相等，即得最后的同构判据。
\end{proof}

\begin{purplenote}
定理~\ref{thm:factor-through-quotient} 也就是说，存在唯一的群同态 $\overline f$，使下图交换，即 $\overline f\circ\pi=f$：
\[
\begin{tikzcd}[row sep=large,column sep=large]
  G \arrow[r,"f"] \arrow[d,"\pi" swap] & H \\
  G/N \arrow[ur,"\overline f" swap] &
\end{tikzcd}
\]
\end{purplenote}

\begin{corollary}[第一同构定理]\label{cor:first-isomorphism}
设 $f:G\to H$ 是群同态，则 $f$ 诱导一个群同构
\[
  G/\ker f\cong\operatorname{im}f.
\]
\end{corollary}

\begin{corollary}\label{cor:induced-quotient-hom}
设 $f:G\to H$ 是群同态，$N\trianglelefteq G$、$M\trianglelefteq H$，且 $f(N)\leq M$。则 $f$ 诱导群同态
\[
  \overline f:G/N\longrightarrow H/M,
  \qquad \overline f(aN)=f(a)M.
\]
$\overline f$ 是群同构，当且仅当
\[
  \operatorname{im}f\vee M=H,
  \qquad f^{-1}(M)\subseteq N.
\]
特别地，若 $f$ 是满同态，且 $f(N)=M$、$\ker f\subseteq N$，则 $\overline f$ 是群同构。
\end{corollary}

\begin{proof}
% 来源：照片 18 底部的 1.5.10 证明。原稿的 \ker\pi\circ f 统一加括号为 \ker(\pi\circ f)。
令 $\pi:H\to H/M$ 为标准投影，则 $\pi\circ f:G\to H/M$ 是群同态，且由 $f(N)\leq M$ 得 $N\subseteq\ker(\pi\circ f)$。
利用定理~\ref{thm:factor-through-quotient}，得到唯一的群同态 $\overline f$，满足
\[
  \overline f(aN)=(\pi\circ f)(a)=f(a)M.
\]
由\hyperref[thm:factor-through-quotient]{同一定理}，$\overline f$ 是同构当且仅当 $\pi\circ f$ 满射且 $N=\ker(\pi\circ f)$。

% 润色说明：展开照片 18 中标为“需多想”的两个等价条件；未改变原稿 1.5.10 的假设或结论。
由于 $M\trianglelefteq H$，有 $\operatorname{im}f\vee M=(\operatorname{im}f)M$，所以
\[
  \pi\circ f\text{ 满射}
  \quad\Longleftrightarrow\quad (\operatorname{im}f)M=H
  \quad\Longleftrightarrow\quad \operatorname{im}f\vee M=H.
\]
又因为 $\ker(\pi\circ f)=f^{-1}(M)$，且已知 $N\subseteq f^{-1}(M)$，故核等于 $N$ 当且仅当 $f^{-1}(M)\subseteq N$，从而得到所述判据。
\end{proof}

\begin{purplenote}
特别情形的理由是：$f$ 满射立即给出 $\operatorname{im}f\vee M=H$；若 $f(N)=M$ 且 $\ker f\subseteq N$，则对任意 $a\in f^{-1}(M)$，可取 $n\in N$ 使 $f(a)=f(n)$。于是 $an^{-1}\in\ker f\subseteq N$，故 $a\in N$，即 $f^{-1}(M)\subseteq N$。
\end{purplenote}

% 来源：微信图片_20260910201745_20_215.jpg，原稿 1.5.11--1.5.16。
\begin{corollary}[第二同构定理]\label{cor:second-isomorphism}
设 $G$ 是群，$N\trianglelefteq G$、$K\leq G$，则
\[
  K/(N\cap K)\cong NK/N.
\]
\end{corollary}

\begin{purpleproof}
由定理~\ref{thm:normal-product}，有 $NK=N\vee K$、$N\cap K\trianglelefteq K$，且 $N\trianglelefteq NK$。
% 修正说明：原紫笔证明所称“恒等映射”实际是包含映射 K -> N\vee K；定义域、陪域不同，不称为恒等映射。
考虑包含映射
\[
  \psi:K\longrightarrow N\vee K,\qquad\psi(k)=k.
\]
则 $\psi$ 是群同态，并且
\[
  \psi(N\cap K)\leq N,\qquad
  \operatorname{im}\psi\vee N=N\vee K,\qquad
  \psi^{-1}(N)=N\cap K.
\]
由推论~\ref{cor:induced-quotient-hom} 即得所述同构。
\end{purpleproof}

\begin{corollary}[第三同构定理]\label{cor:third-isomorphism}
设 $G$ 是群，$H\trianglelefteq G$、$K\trianglelefteq G$，且 $K\leq H$，则
\[
  H/K\trianglelefteq G/K,
  \qquad (G/K)/(H/K)\cong G/H.
\]
\end{corollary}

\begin{proof}
% 来源：微信图片_20260910201750_21_215.jpg；按原稿三步组织，略去划去的推导。
(1) 由 $K\trianglelefteq G$ 及 $K\leq H$，得 $K\trianglelefteq H$。对 $a\in H$，陪集 $aK$ 的所有元素仍属于 $H$，因而将 $H/K=\{aK\mid a\in H\}$ 视为 $G/K$ 的子集是自然的。它对乘法与取逆封闭，所以 $H/K\leq G/K$。

(2) 对任意 $aK\in H/K$、$bK\in G/K$，
\[
  (bK)(aK)(bK)^{-1}=(bab^{-1})K.
\]
因为 $H\trianglelefteq G$，所以 $bab^{-1}\in H$，故上述陪集属于 $H/K$。因此 $H/K\trianglelefteq G/K$。

% 修正说明：原稿第 (3) 步用恒等映射诱导商映射；这里明确先取标准投影 G -> G/H，再通过 G/K 分解，避免把 G/K -> G/H 称为恒等映射。
(3) 考虑标准投影 $q:G\to G/H$。由于 $K\subseteq H=\ker q$，由定理~\ref{thm:factor-through-quotient} 得到群同态
\[
  \overline q:G/K\longrightarrow G/H,
  \qquad\overline q(bK)=bH.
\]
它是满同态，且
\[
  \ker\overline q=\{aK\mid a\in H\}=H/K.
\]
再由\hyperref[cor:first-isomorphism]{第一同构定理}，得到
\[
  (G/K)/(H/K)\cong G/H.
\]
\end{proof}

\begin{purplenote}
第 (3) 步的思路是先从标准投影得到 $\overline q:G/K\to G/H$，再对它的核 $H/K$ 取商，得到同构
\[
  (G/K)/(H/K)\longrightarrow G/H,
  \qquad (aK)(H/K)\longmapsto aH.
\]
\end{purplenote}

\begin{theorem}[子群对应定理]\label{thm:subgroup-correspondence}
设 $f:G\to H$ 是群的满同态，记
\[
  S_f(G)=\{K\leq G\mid\ker f\leq K\},
  \qquad S(H)=\{L\mid L\leq H\}.
\]
则映射
\[
  \psi:S_f(G)\longrightarrow S(H),\qquad\psi(K)=f(K)
\]
是一一对应，且这一对应将 $S_f(G)$ 中 $G$ 的正规子群与 $H$ 的正规子群一一对应起来。
\end{theorem}

\begin{purplenote}
证明一一对应可分别验证满射与单射，也可验证逆映射是 $L\mapsto f^{-1}(L)$。
正规子群之间的对应，只需进一步证明
\[
  K\trianglelefteq G\Longrightarrow f(K)\trianglelefteq H,
  \qquad
  L\trianglelefteq H\Longrightarrow f^{-1}(L)\trianglelefteq G.
\]
这里第一个蕴含使用了 $f$ 的满射性。
\end{purplenote}

\begin{corollary}\label{cor:quotient-subgroup-correspondence}
设 $G$ 是群，$N\trianglelefteq G$，则 $G/N$ 的任意子群都可唯一写为 $K/N$，其中 $N\leq K\leq G$。进一步，
\[
  K/N\trianglelefteq G/N\quad\Longleftrightarrow\quad K\trianglelefteq G.
\]
\end{corollary}

\begin{purplenote}
在定理~\ref{thm:subgroup-correspondence} 中取 $f$ 为标准满同态 $\pi:G\to G/N$，利用 $\ker\pi=N$ 即得。

下面是一些习题结论。
\end{purplenote}

% 来源：照片 20 中原稿 1.5.15 前的未编号结论。
% 修改说明：按用户要求改为黑色编号命题（现 1.5.15）；原 1.5.15--1.5.21 顺延为 1.5.16--1.5.22，结论不变。
\begin{proposition}[指数为二的子群]\label{prop:index-two-normal}
若 $N\leq G$ 且 $[G:N]=2$，则 $N\trianglelefteq G$。
\end{proposition}

\begin{proposition}\label{prop:normal-congruence}
设 $G$ 是群，$N\leq G$，则 $N\trianglelefteq G$ 当且仅当模 $N$ 的右同余是 $G$ 上的一个同余关系。
\end{proposition}

\begin{proposition}\label{prop:congruence-identity-class}
设 $G$ 是群，$\sim$ 是 $G$ 上的等价关系，记
\[
  N=\{a\in G\mid a\sim e\}.
\]
则 $\sim$ 是 $G$ 上的同余关系，当且仅当 $N\trianglelefteq G$ 且 $\sim$ 就是模 $N$ 的同余关系。
\end{proposition}

% 来源：微信图片_20260910201755_22_215.jpg，原稿 1.5.17--1.5.21（现 1.5.18--1.5.22）。
\begin{example}\label{ex:normality-not-transitive}
% 转录补充：原稿引用教材印刷页码 25 的 D_4^*；为使笔记可独立阅读，补明其为正方形对称群及相关变换的含义。
设 $D_4^*$ 为教材中的正方形对称群，它是一个 $8$ 阶群，运算为映射复合。以 $I$ 表示恒等变换，$R$ 表示绕正方形中心旋转 $90^\circ$ 的变换，$T_x,T_y$ 分别表示关于水平轴和竖直轴的反射。取
\[
  H=\{T_y,I\},\qquad K=\{I,R^2,T_x,T_y\},
\]
则
\[
  H\trianglelefteq K,\qquad K\trianglelefteq D_4^*,
  \qquad H\not\trianglelefteq D_4^*.
\]
\end{example}

\begin{proposition}\label{prop:subgroups-of-cyclic-normal}
设 $G$ 是群，$H$ 是 $G$ 的循环子群，且 $H\trianglelefteq G$。则 $H$ 的任意子群都是 $G$ 的正规子群。
\end{proposition}

\begin{theorem}\label{thm:finitely-generated-extension}
设 $G$ 是群，$H\trianglelefteq G$。若 $H$ 和 $G/H$ 都是有限生成群，则 $G$ 也是有限生成群。
\end{theorem}

\begin{proposition}\label{prop:quotient-direct-product}
设 $G_1,G_2$ 是群，$N_1\trianglelefteq G_1$、$N_2\trianglelefteq G_2$，则
\[
  N_1\times N_2\trianglelefteq G_1\times G_2,
\]
且
\[
  (G_1\times G_2)/(N_1\times N_2)
  \cong(G_1/N_1)\times(G_2/N_2).
\]
\end{proposition}

\begin{theorem}\label{thm:preimage-image-subgroup}
设 $f:G\to H$ 是群同态，$N=\ker f$，且 $K\leq G$，则
\[
  f^{-1}\bigl(f(K)\bigr)=KN.
\]
进一步，
\[
  f^{-1}\bigl(f(K)\bigr)=K\quad\Longleftrightarrow\quad N\leq K.
\]
\end{theorem}

% 新增：第 1.6 节独立维护，保留原稿编号及黑、紫双色区分。
% !TeX root = ../main.tex
\subsection{对称群、偶置换与二面体群}

% 来源：微信图片_20260914150308_24_215.jpg，1.6.1--1.6.6。
% 排版约定：原黑笔正文和证明用普通环境，原紫笔批注用 purplenote；单独的证明移至对应结论后。
% 转录补充：明确映射复合的运算次序，便于核对本节的循环乘积公式。
本节记 $I_n=\{1,2,\ldots,n\}$，以 $S_n$ 表示 $I_n$ 上的对称群。置换的乘积采用映射复合，即
\[
  (\sigma\tau)(k)=\sigma\bigl(\tau(k)\bigr),
\]
也就是先作用右边的置换。恒等置换记作 $(1)$。

\begin{definition}[循环与对换]\label{def:cycles-transpositions}
% 排版修正：将取值条件与循环记号分段，避免行内公式越过页边界。
设 $1\leq r\leq n$，$i_1,\ldots,i_r\in I_n$ 互不相同。

记 $(i_1\,i_2\,\cdots\,i_r)$ 为置换
\[
  i_1\longmapsto i_2,\quad i_2\longmapsto i_3,\quad\ldots,\quad
  i_{r-1}\longmapsto i_r,\quad i_r\longmapsto i_1,
\]
而 $I_n$ 中其余的元素都映到自身。称这一置换为长度为 $r$ 的\textbf{循环}，简称 $r$-循环。一个 $2$-循环称为\textbf{对换}。
\end{definition}

\begin{purplenote}
由定义立即得到
\[
  (i_1\,i_2\,\cdots\,i_r)
  =(i_2\,i_3\,\cdots\,i_r\,i_1)=\cdots,
\]
且
\[
  (i_1\,i_2\,\cdots\,i_r)^{-1}=(i_r\,i_{r-1}\,\cdots\,i_1).
\]
一般的循环（更一般地，置换）不一定可交换，即不一定有 $\sigma\tau=\tau\sigma$，但在一些特殊情形下二者可交换。
\end{purplenote}

\begin{definition}[不相交的置换]\label{def:disjoint-permutations}
设 $\sigma_1,\ldots,\sigma_r\in S_n$。若对任意 $1\leq i\leq r$ 及 $k\in I_n$，都有
\[
  \sigma_i(k)\neq k\quad\Longrightarrow\quad
  \sigma_j(k)=k\qquad(j\neq i),
\]
则称 $\sigma_1,\ldots,\sigma_r$ 是\textbf{不相交的置换}。
\end{definition}

\begin{purplenote}
换言之，$\sigma_1,\ldots,\sigma_r$ 不相交，当且仅当 $I_n$ 中任意元素至多被其中一个置换移动，即至多有一个 $\sigma_i$ 将它映为不同于自身的元素。因此，当 $\sigma$ 与 $\tau$ 不相交时，有
\[
  \sigma\tau=\tau\sigma.
\]
\end{purplenote}

\begin{theorem}\label{thm:disjoint-cycle-decomposition}
% 修正说明：原稿叙述漏写“不相交”，证明页已补出；唯一分解必须限定为两两不相交的循环。
$S_n$ 中任意非恒等置换都可分解为两两不相交的循环的乘积，其中每个循环的长度至少为 $2$。在忽略循环因子的排列次序以及同一循环的循环移位写法的意义下，这一分解是唯一的。
\end{theorem}

\begin{proof}
% 来源：微信图片_20260914150318_26_215.jpg 顶部的 1.6.3 证明。
对任意 $k\in I_n$，由于 $I_n$ 有限，且 $\sigma$ 是双射，必有正整数 $m$ 使 $\sigma^m(k)=k$。于是集合
\[
  B(k)=\{\sigma^j(k)\mid j\in\mathbb Z\}
\]
有限；任意两个这样的集合要么相同，要么不相交。因此它们将 $I_n$ 分成有限个互不相交的集合 $B_1,\ldots,B_t$。

对每个 $B_i$ 定义
\[
  \sigma_i(x)=
  \begin{cases}
    \sigma(x),&x\in B_i,\\
    x,&x\notin B_i.
  \end{cases}
\]
若 $x\in B_i$，且 $m$ 是使 $\sigma^m(x)=x$ 的最小正整数，则
\[
  \sigma_i=\bigl(x\,\sigma(x)\,\cdots\,\sigma^{m-1}(x)\bigr).
\]
这些循环两两不相交，且 $\sigma=\sigma_1\cdots\sigma_t$。
% 补充说明：略去单点集合对应的恒等循环，以满足定理中循环长度至少为 2 的要求。
略去 $|B_i|=1$ 时的恒等循环，即得所需分解。

为证唯一性，设 $\sigma=\tau_1\cdots\tau_u$ 是另一组两两不相交、长度至少为 $2$ 的循环的乘积。若 $\tau_j(x)\neq x$，则 $\sigma(x)=\tau_j(x)$，进一步有
\[
  \sigma^k(x)=\tau_j^k(x)\qquad(k\in\mathbb Z).
\]
因此 $\tau_j$ 所移动的元素恰好构成某个 $B_i$，且其循环次序也由 $\sigma$ 确定。这就证明了唯一性。
\end{proof}

\begin{corollary}\label{cor:permutation-order}
对任意 $\sigma\in S_n$，$\sigma$ 的阶等于其不相交循环分解中各循环的阶（也就是长度）的最小公倍数。即若
\[
  \sigma=\sigma_1\cdots\sigma_r,
\]
其中各循环两两不相交，则
\[
  |\sigma|=\operatorname{lcm}\bigl(|\sigma_1|,\ldots,|\sigma_r|\bigr).
\]
% 补充说明：明确恒等置换这一边界情形，不必使用空集的最小公倍数记号。
恒等置换的阶为 $1$。
\end{corollary}

\begin{corollary}\label{cor:transposition-decomposition}
任意 $\sigma\in S_n$ 都可分解为对换的乘积，其中的对换不一定互不相交。
% 补充说明：允许空乘积表示恒等置换，以包括 n=1 的情形。
这里允许用空乘积表示恒等置换。
\end{corollary}

\begin{definition}[偶置换与奇置换]\label{def:permutation-parity}
对任意 $\sigma\in S_n$，若 $\sigma$ 可分解为偶数个对换的乘积，则称 $\sigma$ 为\textbf{偶置换}；若可分解为奇数个对换的乘积，则称其为\textbf{奇置换}。
\end{definition}

\begin{purplenote}
这一定义的良定义性由定理~\ref{thm:parity-well-defined} 保证。
\end{purplenote}

% 来源：微信图片_20260914150312_25_215.jpg，1.6.7--1.6.13。
\begin{theorem}\label{thm:parity-well-defined}
任意 $\sigma\in S_n$ 不可能既是偶置换又是奇置换。
\end{theorem}

\begin{proof}
% 来源：照片 26 中部的 1.6.7 证明；展开对换使乘积变号及奇偶性唯一的推导。
设
\[
  \sigma=\begin{pmatrix}
    1&2&\cdots&n\\
    i_1&i_2&\cdots&i_n
  \end{pmatrix},
  \qquad
  \Delta(\sigma)=\prod_{1\leq k<j\leq n}(i_k-i_j).
\]
因 $i_1,\ldots,i_n$ 两两不同，$\Delta(\sigma)\neq0$。若 $\tau$ 是一个对换，则左乘 $\tau$ 恰好交换上述 $i_1,\ldots,i_n$ 中的两个值，故
\[
  \Delta(\tau\sigma)=-\Delta(\sigma).
\]
所以若 $\sigma$ 分解为 $r$ 个对换的乘积，则
\[
  \Delta(\sigma)=(-1)^r\Delta((1)).
\]
由于 $\Delta((1))\neq0$，不同分解中对换个数的奇偶性必相同。$n=1$ 时只有恒等置换，结论也成立。
\end{proof}

\begin{theorem}\label{thm:alternating-group}
对任意 $n\geq2$，记 $A_n$ 为 $S_n$ 中所有偶置换组成的集合，则
\[
  A_n\trianglelefteq S_n,\qquad
  [S_n:A_n]=2,\qquad |A_n|=\frac{|S_n|}{2}=\frac{n!}{2}.
\]
进一步，$A_n$ 是 $S_n$ 中唯一的指数为 $2$ 的子群。
% 转录补充：补明 A_n 的名称，不另增编号。
$A_n$ 称为 $n$ 元\textbf{交错群}。
\end{theorem}

\begin{purplenote}
% 来源说明：原稿将唯一性的习题出处记作“教材习题 8”，此处保留原稿提示。
唯一性的证明见教材习题 $8$（原稿所记）。
\end{purplenote}

\begin{definition}[简单群]\label{def:simple-group}
设 $G$ 是群。若 $G$ 没有非平凡真正规子群，即 $G$ 的正规子群只能是 $\{e\}$ 或 $G$ 本身，则称 $G$ 为\textbf{简单群}。
\end{definition}

% 约定说明：遵循本节原稿及 Hungerford Definition 6.9 / Theorem 6.10 的约定，不额外要求简单群非平凡；否则下面 n\ne4 的说法需调整。
\par\textbf{术语约定。}这里的定义不排除平凡群。有些文献在简单群的定义中还要求 $G\neq\{e\}$，阅读时需注意这一区别。

\begin{purplenote}
% 修正说明：原稿“简单阿贝尔群只能同构于 Z_p”需限定非平凡群，才能与本节允许平凡简单群的约定一致。
阿贝尔群的任意子群都是正规子群，而无限群必有非平凡真子群，故非平凡的简单阿贝尔群一定有限，进一步只能同构于 $\mathbb Z_p$，其中 $p$ 为素数。
非阿贝尔简单群也有许多例子，例如下面的交错群。
\end{purplenote}

\begin{theorem}\label{thm:alternating-simple}
% 约定说明：保留原稿 n\ne4 的结论；n=2 时 A_2 为平凡群，在上面的简单群约定下仍满足结论。
设 $n\geq2$。若 $n\neq4$，则 $A_n$ 是简单群。
\end{theorem}

\begin{purplenote}
定理~\ref{thm:alternating-simple} 的证明涉及以下两个重要引理。
\end{purplenote}

\begin{lemma}\label{lem:alternating-generators}
设 $n\geq3$，$r,s$ 是 $I_n$ 中两个不同的元素，则 $A_n$ 由下列 $3$-循环生成：
\[
  A_n=\bigl\langle(r\,s\,k)\mid 1\leq k\leq n,\ k\neq r,s\bigr\rangle.
\]
\end{lemma}

\begin{proof}
% 来源：照片 26 下半页的 1.6.11 证明；所有乘积均按从右向左作用核对。
$n=3$ 时，$A_3$ 由恒等置换及两个 $3$-循环组成，结论可直接验证。以下设 $n\geq4$。

任意 $\sigma\in A_n$ 都可写成偶数个对换的乘积。将相邻的对换两两配对：相同的对换相消，其余各对可写成 $(a\,b)(c\,d)$ 或 $(a\,b)(a\,c)$ 的形式。注意到
\[
  (a\,b)(c\,d)=(a\,c\,b)(c\,d\,a),
  \qquad (a\,b)(a\,c)=(a\,c\,b),
\]
所以每个偶置换都是 $3$-循环的乘积。

令 $L=\langle(r\,s\,k)\mid k\neq r,s\rangle$。只需说明每个 $3$-循环都属于 $L$。
按是否含有 $r,s$，并允许循环移位及取逆，只需考虑
\[
  (r\,s\,a),\quad(r\,a\,b),\quad(s\,a\,b),\quad(a\,b\,c),
\]
其中 $a,b,c$ 互不相同且都不等于 $r,s$；最后一种情形仅在 $n\geq5$ 时出现。由
\begin{align*}
  (r\,a\,s)&=(r\,s\,a)^2,\\
  (r\,a\,b)&=(r\,s\,b)(r\,s\,a)^2,\\
  (s\,a\,b)&=(r\,s\,b)^2(r\,s\,a),\\
  (a\,b\,c)&=(r\,s\,a)^2(r\,s\,c)(r\,s\,b)^2(r\,s\,a),
\end{align*}
可知这些循环都属于 $L$。另一方面，$3$-循环是偶置换，所以 $L\leq A_n$，故 $L=A_n$。
\end{proof}

\begin{lemma}\label{lem:normal-containing-three-cycle}
设 $n\geq3$。若 $N\trianglelefteq A_n$，且 $N$ 包含一个 $3$-循环，则 $N=A_n$。
\end{lemma}

\begin{proof}
% 来源：微信图片_20260914150322_27_215.jpg 顶部的 1.6.12 证明。
设 $(r\,s\,c)\in N$。若 $k\notin\{r,s,c\}$，则 $(r\,s)(c\,k)\in A_n$，并且
\begin{align*}
  (r\,s\,k)
  &=(r\,s)(c\,k)(r\,s\,c)^2(c\,k)(r\,s)\\
  &=\bigl((r\,s)(c\,k)\bigr)(r\,s\,c)^2
    \bigl((r\,s)(c\,k)\bigr)^{-1}\in N.
\end{align*}
这里使用了 $N\trianglelefteq A_n$。连同已知的 $(r\,s\,c)\in N$，可知 $N$ 包含引理~\ref{lem:alternating-generators} 中的全部生成元，故 $N=A_n$。这也包括 $n=3$ 时只有一个所需生成元的情形。
\end{proof}

\begin{purplenote}
定理~\ref{thm:alternating-simple} 的完整证明详见教材。主要策略是取 $A_n$ 的一个非平凡正规子群 $N$，分类讨论其元素的循环分解，归结为引理~\ref{lem:normal-containing-three-cycle}。这一分类论证用于 $n\geq5$；小的 $n$ 单独处理。
\end{purplenote}

\begin{definition}[二面体群]\label{def:dihedral-group}
% 补充说明：明确 n\geq3，与后面的阶及关系式定理保持一致。
设 $n\geq3$，在 $S_n$ 中取
\[
  a=(1\,2\,\cdots\,n),\qquad
  b=\begin{pmatrix}
    1&2&3&\cdots&n\\
    1&n&n-1&\cdots&2
  \end{pmatrix}.
\]
也可写成不相交对换的乘积
\[
  b=\prod_{2\leq i<(n+2)/2}(i\,\,n+2-i).
\]
子群 $D_n=\langle a,b\rangle\leq S_n$ 称为次数为 $n$ 的\textbf{二面体群}。
\end{definition}

% 来源：微信图片_20260914150326_28_215.jpg，1.6.14--1.6.17。
\begin{theorem}\label{thm:dihedral-relations}
对任意 $n\geq3$，二面体群 $D_n$ 的阶为 $2n$，且其生成元 $a,b$ 满足：
\begin{enumerate}[label=(\arabic*)]
  \item $a^n=(1)$，$b^2=(1)$，且 $a^k\neq(1)$ 对所有 $0<k<n$ 成立；
  \item $ba=a^{-1}b$。
\end{enumerate}
反过来，若群 $G=\langle a,b\rangle$ 的生成元 $a,b$ 对某个 $n\geq3$ 满足以上两条关系（将 $(1)$ 换成 $G$ 的幺元 $e$），则 $G\cong D_n$。
\end{theorem}

\begin{proof}
% 来源：照片 27 中部的 1.6.14 证明；按原稿的标准形思路补足 2n 个元素互异及同构保持乘法的验证。
定义~\ref{def:dihedral-group} 中的置换 $a,b$ 满足上述关系，且 $a$ 的阶为 $n$。由 $ba=a^{-1}b$ 可得
\[
  ba^i=a^{-i}b\qquad(i\in\mathbb Z).
\]
因此每个由 $a,b$ 生成的元素都可写成 $a^i b^j$，其中 $0\leq i<n$、$j=0,1$。

事实上，这 $2n$ 个标准形互不相同。首先 $b\notin\langle a\rangle$：否则 $a,b$ 可交换，结合 $ba=a^{-1}b$ 得 $a=a^{-1}$，与 $|a|=n\geq3$ 矛盾。
同一类中的 $a^i$（或 $a^i b$）两两不同；而 $a^i=a^k b$ 会推出 $b\in\langle a\rangle$，也不可能。故 $|D_n|=2n$。

反过来，若 $G=\langle a,b\rangle$ 满足上述关系，同样的论证给出
\[
  G=\{a^i b^j\mid 0\leq i<n,\ j=0,1\},
\]
且每个元素的这种写法唯一。为区分两群的生成元，将 $D_n$ 中的生成元记为 $\alpha,\beta$。映射
\[
  G\longrightarrow D_n,\qquad a^i b^j\longmapsto\alpha^i\beta^j
\]
是双射。两群中的乘法都满足
\[
  (a^i b^j)(a^k b^\ell)=a^{\,i+(-1)^j k}b^{\,j+\ell},
\]
其中 $a$ 的指数按模 $n$ 化简，$b$ 的指数按模 $2$ 化简。因此该双射保持乘法，是群同构。
\end{proof}

\begin{purplenote}
本质上，二面体群 $D_n$ 与正 $n$ 边形的全体对称在复合下组成的群同构。

以下是一些练习题结论。
\end{purplenote}

\begin{proposition}\label{prop:transposition-generating-sets}
设 $n\geq2$，则：
\begin{enumerate}[label=(\arabic*)]
  \item $S_n=\langle(1\,2),(1\,3),\ldots,(1\,n)\rangle$；
  \item $S_n=\langle(1\,2),(2\,3),\ldots,(n-1\,\,n)\rangle$。
\end{enumerate}
\end{proposition}

\begin{proposition}\label{prop:conjugation-of-cycle}
若 $\sigma=(i_1\,i_2\,\cdots\,i_r)\in S_n$，$\tau\in S_n$，则
\[
  \tau\sigma\tau^{-1}
  =\bigl(\tau(i_1)\,\tau(i_2)\,\cdots\,\tau(i_r)\bigr).
\]
\end{proposition}

\begin{proposition}\label{prop:two-generators-symmetric}
% 补充说明：明确生成式的适用范围；第 (2) 条取 n\geq3，使 (2\,3\,\cdots\,n) 的写法不含歧义。
\begin{enumerate}[label=(\arabic*)]
  \item 设 $n\geq2$，$\sigma_1=(1\,2)$，$\tau=(1\,2\,\cdots\,n)$，则
  \[
    S_n=\langle\sigma_1,\tau\rangle.
  \]
  \item 设 $n\geq3$，则
  \[
    S_n=\langle(1\,2),(2\,3\,\cdots\,n)\rangle.
  \]
\end{enumerate}
\end{proposition}
% 新增：第 1.7 节独立维护，按原稿整理范畴、积、余积与自由对象。
% !TeX root = ../main.tex
\subsection{范畴论：积、余积与自由对象}

% 来源：微信图片_20261008195012_31_215.jpg，1.7.1--1.7.3。
% 排版约定：黑笔定义、定理用普通环境；原稿中的紫笔解释和证明提示保留为紫色。
\begin{definition}[范畴]\label{def:category}
一个\textbf{范畴} $\mathcal C$ 由以下数据组成：
\begin{enumerate}[label=(\arabic*)]
  \item 一类对象，记作 $A,B,C,\ldots$；
  \item 对每对对象 $A,B$，一个态射集合 $\operatorname{Hom}_{\mathcal C}(A,B)$，其元素称为从 $A$ 到 $B$ 的\textbf{态射}，记作 $f:A\to B$；
  \item 对任意对象 $A,B,C$，一个复合运算
  \[
    \operatorname{Hom}_{\mathcal C}(B,C)\times
    \operatorname{Hom}_{\mathcal C}(A,B)
    \longrightarrow \operatorname{Hom}_{\mathcal C}(A,C),
    \qquad (g,f)\longmapsto g\circ f.
  \]
\end{enumerate}
这些数据满足：
\begin{enumerate}[label=(\arabic*),resume]
  \item 复合满足结合律：若 $f:A\to B$、$g:B\to C$、$h:C\to D$，则
  \[
    h\circ(g\circ f)=(h\circ g)\circ f;
  \]
  \item 对每个对象 $A$，存在恒等态射 $1_A:A\to A$，且对任意态射 $f:A\to B$，
  \[
    1_B\circ f=f=f\circ1_A.
  \]
\end{enumerate}
\end{definition}

\begin{purplenote}
范畴中的对象可以是集合、群等代数对象，也可以是其他数学对象；态射是对象之间指定的映射，未必是通常意义下的集合函数，也不必有逆态射。结合律与恒等态射的条件分别要求复合可结合，以及恒等态射在复合中不起作用。
\end{purplenote}

\begin{definition}[等价对象（同构对象）]\label{def:equivalent-objects}
设 $A,B$ 是范畴 $\mathcal C$ 中的对象。若存在态射 $f:A\to B$ 和 $g:B\to A$，使得
\[
  g\circ f=1_A,\qquad f\circ g=1_B,
\]
则称 $f$、$g$ 互为\textbf{逆态射}，并称 $A$ 与 $B$ \textbf{等价}（或\textbf{同构}）。
\end{definition}

\begin{purplenote}
两个等价态射的复合仍是等价态射。
\end{purplenote}

\begin{definition}[积]\label{def:category-product}
设 $\{A_i\mid i\in I\}$ 是范畴 $\mathcal C$ 中的一族对象。若对象 $P$ 连同一族态射 $\pi_i:P\to A_i$ 满足下列泛性质，则称 $(P,\{\pi_i\}_{i\in I})$ 为这族对象的\textbf{积}：对任意对象 $B$ 及任意一族态射 $\psi_i:B\to A_i$，存在唯一态射 $\psi:B\to P$，使得
\[
  \pi_i\circ\psi=\psi_i\qquad(i\in I).
\]
\end{definition}

\begin{purplenote}
积对象 $P$ 通常记作 $\prod_{i\in I}A_i$。不是每个范畴中的每一族对象都必定有积。
\end{purplenote}

\begin{theorem}\label{thm:products-unique}
若 $(P,\{\pi_i\}_{i\in I})$ 和 $(Q,\{\rho_i\}_{i\in I})$ 都是范畴 $\mathcal C$ 中对象族 $\{A_i\mid i\in I\}$ 的积，则 $P$ 与 $Q$ 等价。
\end{theorem}

\begin{purpleproof}
由 $P$、$Q$ 的积泛性质，存在唯一态射 $u:P\to Q$ 和 $v:Q\to P$，使得
\[
  \rho_i\circ u=\pi_i,\qquad \pi_i\circ v=\rho_i\qquad(i\in I).
\]
于是 $\pi_i\circ(v\circ u)=\pi_i$。恒等态射 $1_P$ 也满足 $\pi_i\circ1_P=\pi_i$；由积泛性质的唯一性，$v\circ u=1_P$。同理，$u\circ v=1_Q$，故 $P$ 与 $Q$ 等价。
\end{purpleproof}

\begin{definition}[余积]\label{def:category-coproduct}
设 $\{A_i\mid i\in I\}$ 是范畴 $\mathcal C$ 中的一族对象。若对象 $S$ 连同一族态射 $\tau_i:A_i\to S$ 满足下列泛性质，则称 $(S,\{\tau_i\}_{i\in I})$ 为这族对象的\textbf{余积}：对任意对象 $B$ 及任意一族态射 $\psi_i:A_i\to B$，存在唯一态射 $\psi:S\to B$，使得
\[
  \psi\circ\tau_i=\psi_i\qquad(i\in I).
\]
\end{definition}

\begin{theorem}\label{thm:coproducts-unique}
若 $(S,\{\tau_i\}_{i\in I})$ 和 $(S',\{\tau'_i\}_{i\in I})$ 都是范畴 $\mathcal C$ 中对象族 $\{A_i\mid i\in I\}$ 的余积，则 $S$ 与 $S'$ 等价。
\end{theorem}

\begin{purpleproof}
证明与定理~\ref{thm:products-unique} 对偶。由余积的泛性质，得到相互逆的唯一态射 $u:S\to S'$、$v:S'\to S$，满足 $u\circ\tau_i=\tau'_i$、$v\circ\tau'_i=\tau_i$。
\end{purpleproof}

% 来源：微信图片_20261008195026_32_215.jpg，1.7.7--1.7.9。
\begin{definition}[具体范畴]\label{def:concrete-category}
\textbf{具体范畴}是一个范畴 $\mathcal C$ 连同一个赋值 $\sigma$，它给每个对象 $A$ 指定一个集合 $\sigma(A)$（称为 $A$ 的\textbf{底集}），并满足：
\begin{enumerate}[label=(\arabic*)]
  \item $\mathcal C$ 中每个态射 $f:A\to B$ 都是底集间的函数 $\sigma(f):\sigma(A)\to\sigma(B)$；
  \item $A$ 的恒等态射 $1_A$ 在底集上是恒等函数：$\sigma(1_A)=1_{\sigma(A)}$；
  \item 范畴中的态射复合与底集函数的复合相容：$\sigma(g\circ f)=\sigma(g)\circ\sigma(f)$。
\end{enumerate}
\end{definition}

\begin{purplenote}
% 补充说明：用群范畴中的具体函数解释“底集函数未必是范畴态射”。
例如，群范畴中每个群的底集就是它的元素集合，每个态射在底集上是群同态所对应的函数。但底集之间的任意函数未必保持群结构：把加法群 $\mathbb Z$ 映到自身，函数 $f(n)=n+1$ 是底集间的函数，却不是群同态，因为 $f(0)=1\neq0$，而群同态必须把单位元 $0$ 映到 $0$。所以群范畴中的态射只是底集函数中的一部分。通常为简便，会把对象 $A$ 与它的底集 $\sigma(A)$ 用同一符号表示；在集合范畴中，态射就是任意集合函数，但在群范畴中还须保持群运算。
\end{purplenote}

\begin{purplenote}
% 转录补充：应用户要求，加入常见具体范畴的例子，并与仅由对象、态射及复合组成的一般范畴对照；此备注不另编号。
\textbf{常见具体范畴的例子。}
\begin{enumerate}[label=(\arabic*)]
  \item 集合范畴 $\mathbf{Set}$：对象是集合，态射是任意集合函数；
  \item 群范畴 $\mathbf{Grp}$：对象是群，态射是群同态；阿贝尔群范畴 $\mathbf{Ab}$ 是其子范畴；
  \item 偏序集范畴 $\mathbf{Pos}$：对象是偏序集，态射是保序映射。
\end{enumerate}
这些范畴都能把每个对象看成一个底集，并把每个态射看成底集之间的函数。一般范畴的定义本身不要求有这样的底集或函数表示；具体范畴则额外指定这种表示，并要求它与恒等态射、复合相容。换句话说，具体范畴保留了“底层元素”和“底层函数”可供讨论。

本节沿用教材列出的三项条件。常见的更严格约定还要求底集赋值在每个态射集合上是单射：若两个态射诱导出同一个底集函数，则它们本来就是同一个态射。前三项条件保证底集赋值与恒等态射、复合相容；加上这项单射条件，才是通常所说的忠实底集赋值。具体范畴可以理解为：一般范畴结构之外，还能一致地谈论对象的底层元素和态射的底层函数。
\end{purplenote}

\begin{definition}[自由对象]\label{def:free-object}
% 记号说明：i 与 f 是底集之间的集合映射；其唯一延拓 \overline f 才是范畴中的态射。
设 $\mathcal C$ 是具体范畴，$F$ 是 $\mathcal C$ 中的对象，$X$ 是非空集合，且 $i:X\to\sigma(F)$ 是一个集合映射。若对任意对象 $A\in\mathcal C$ 及任意集合映射 $f:X\to\sigma(A)$，都存在唯一态射 $\overline f:F\to A$，使得
\[
  \sigma(\overline f)\circ i=f,
\]
则称 $F$ 是\textbf{由集合 $X$ 自由生成的对象}，或称 $F$ 是 $X$ 上的\textbf{自由对象}。
\end{definition}

\begin{purplenote}
常把底集映射 $\sigma(\overline f)$ 简写为 $\overline f$，于是条件写作 $\overline f\circ i=f$。自由对象的泛性质表示：从 $F$ 出发的态射由它在指定生成元 $i(X)$ 上的取值唯一确定，而且任意给定这些生成元的像都能唯一延拓为态射。
\end{purplenote}

\begin{theorem}\label{thm:free-objects-equivalent}
设 $\mathcal C$ 是具体范畴，$F,F'$ 是 $\mathcal C$ 中的对象，分别由映射 $i:X\to\sigma(F)$ 和 $i':X'\to\sigma(F')$ 自由生成。若 $|X|=|X'|$，则 $F$ 与 $F'$ 等价。
\end{theorem}

\begin{purplenote}
% 转录补充：根据原稿中的泛性质提示，展开以双射 X→X' 诱导对象同构的证明。
取双射 $b:X\to X'$。由自由对象的泛性质，映射 $i'\circ b:X\to\sigma(F')$ 唯一延拓为态射 $u:F\to F'$；同理，$i\circ b^{-1}$ 唯一延拓为态射 $v:F'\to F$。复合 $v\circ u$ 在 $i(X)$ 上与 $1_F$ 一致，故由唯一性 $v\circ u=1_F$；同理 $u\circ v=1_{F'}$。这是定理~\ref{thm:free-objects-equivalent} 的证明思路；$X=X'$ 是其特殊情形。
\end{purplenote}

% 来源：微信图片_20261008195034_33_215.jpg，1.7.10--1.7.12。
\begin{definition}[泛对象与余泛对象]\label{def:initial-terminal}
设 $\mathcal C$ 是范畴。
\begin{enumerate}[label=(\arabic*)]
  \item 若对象 $I\in\mathcal C$ 满足对每个对象 $A\in\mathcal C$，都存在唯一态射 $I\to A$，则称 $I$ 为\textbf{泛对象}（或\textbf{初始对象}）。
  \item 若对象 $T\in\mathcal C$ 满足对每个对象 $A\in\mathcal C$，都存在唯一态射 $A\to T$，则称 $T$ 为\textbf{余泛对象}（或\textbf{终对象}）。
\end{enumerate}
\end{definition}

\begin{theorem}\label{thm:initial-terminal-unique}
范畴 $\mathcal C$ 中任意两个泛对象等价，任意两个余泛对象也等价。
\end{theorem}

\begin{theorem}\label{thm:free-group-generators}
在群范畴中，设 $F$ 是非空集合 $X$ 上的自由对象，$i:X\to\sigma(F)$ 是相应的集合映射（F的存在性将在第~\ref{sec:free-groups}~节中证明），则 $i(X)$ 是 $F$ 的生成集，即
\[
  F=\langle i(X)\rangle.
\]
\end{theorem}

\begin{purpleproof}
% 转录补充：展开原稿 1.7.12 的证明提示，用自由对象的唯一延拓性质证明 i(X) 生成 F。
令 $U=\langle i(X)\rangle\leq F$，并令 $j:U\hookrightarrow F$ 为包含同态。由于 $i(X)\subseteq U$，可将 $i$ 看成集合映射 $i_U:X\to U$。由 $F$ 的自由性质，存在群同态 $\varphi:F\to U$ 使 $\varphi\circ i=i_U$。于是 $j\circ\varphi:F\to F$ 满足 $(j\circ\varphi)\circ i=i$。恒等同态 $1_F$ 也满足此条件，故由自由性给出的唯一性，$j\circ\varphi=1_F$。所以 $j$ 是满射，从而 $U=F$。
\end{purpleproof}

% 新增：第 1.8 节独立维护，整理群的直积、弱直积与内直积。
% !TeX root = ../main.tex
\subsection{直积和直和}

% 来源：微信图片_20261008195558_34_215.jpg 至微信图片_20261008195616_39_215.jpg，1.8.1--1.8.15。
% 排版约定：黑笔定义、定理和命题用普通环境；原紫笔说明及证明提示保留为紫色。
\begin{purplenote}
本节讨论如何由已知的群构造新群，并在已知各组成部分结构的情形下描述所得群的结构。
\end{purplenote}

\begin{definition}[直积与完全直积]\label{def:external-direct-product}
设 $\{G_i\mid i\in I\}$ 是一族群。在集合
\[
  \prod_{i\in I}G_i
  =\{(g_i)_{i\in I}\mid g_i\in G_i\text{ 对每个 }i\in I\}
\]
上定义二元运算：对 $f=(f(i))_{i\in I}$、$g=(g(i))_{i\in I}$，令
\[
  (fg)(i)=f(i)g(i)\qquad(i\in I).
\]
此运算称为逐坐标运算。所得集合连同此运算称为这族群的\textbf{完全直积}。当 $I=\{1,2,\ldots,n\}$ 时，通常记作
\[
  G_1\times G_2\times\cdots\times G_n.
\]
\end{definition}

\begin{purplenote}
若各 $G_i$ 采用加法记号，则逐坐标运算写作
\[
  (f+g)(i)=f(i)+g(i)\qquad(i\in I).
\]
当 $I=\{1,2,\ldots,n\}$ 且各 $G_i$ 是阿贝尔群时，有限直积也常记作
\[
  G_1\oplus G_2\oplus\cdots\oplus G_n;
\]
有限指标下，直和与直积相同。
\end{purplenote}

\begin{theorem}\label{thm:direct-product-group}
设 $\{G_i\mid i\in I\}$ 是一族群，则：
\begin{enumerate}[label=(\arabic*)]
  \item 完全直积 $\prod_{i\in I}G_i$ 在逐坐标运算下也是一个群；
  \item 对任意 $k\in I$，映射
  \[
    \pi_k:\prod_{i\in I}G_i\longrightarrow G_k,
    \qquad (g_i)_{i\in I}\longmapsto g_k
  \]
  是满同态。
\end{enumerate}
\end{theorem}

\begin{purplenote}
$\pi_k$ 称为第 $k$ 个坐标投影，称为直积的标准投影。
\end{purplenote}

\begin{theorem}\label{thm:direct-product-universal-property}
设 $\{G_i\mid i\in I\}$ 是一族群，$H$ 是群，且对每个 $i\in I$ 给定群同态 $\varphi_i:H\to G_i$。则存在唯一群同态
\[
  \varphi:H\longrightarrow\prod_{i\in I}G_i
\]
使得
\[
  \pi_i\circ\varphi=\varphi_i\qquad(i\in I).
\]
即$(\prod_{i\in I}G_i, \{\pi_i\}_{i\in I})$是群范畴下$\{G_i\mid i\in I\}$的一个积
\end{theorem}

\begin{purplenote}
由范畴论中积在等价的意义下唯一知群范畴中一族给定群的完全直积在同构意义下是唯一的。若各 $G_i$ 都是阿贝尔群，则其完全直积仍是阿贝尔群；因此它不仅是群范畴中的积，也是阿贝尔群范畴中的积。
\end{purplenote}

\begin{definition}[外弱直积与外直和]\label{def:external-weak-direct-product}
设 $\{G_i\mid i\in I\}$ 是一族群。完全直积中满足“除有限多个 $i$ 外，$g_i=e_i$”的元素所成的子集记作
\[
  \prod\nolimits_{i\in I}^{w}G_i
  =\{(g_i)_{i\in I}\in\prod_{i\in I}G_i
      \mid g_i=e_i\text{ 除有限多个 }i\text{ 外}
    \},
\]
称为这族群的\textbf{外弱直积}。若各 $G_i$ 是加法记号的阿贝尔群，则外弱直积通常称为\textbf{外直和}，记作 $\sum_{i\in I}G_i$。
\end{definition}

\begin{purplenote}
当指标集 $I$ 有限时，外弱直积就是完全直积。
\end{purplenote}

\begin{theorem}\label{thm:weak-product-embeddings}
设 $\{G_i\mid i\in I\}$ 是一族群，则：
\begin{enumerate}[label=(\arabic*)]
  \item 外弱直积 $\prod\nolimits_{i\in I}^{w}G_i$ 是完全直积 $\prod_{i\in I}G_i$ 的一个正规子群；
  \item 对每个 $k\in I$，映射
  \[
    \tau_k:G_k\longrightarrow\prod\nolimits_{i\in I}^{w}G_i,
    \qquad a\longmapsto(g_i)_{i\in I},
  \]
  其中 $g_k=a$ 且 $g_i=e_i$（$i\neq k$），是单同态；
  \item 对每个 $k\in I$，$\tau_k(G_k)$ 是外弱直积$\prod\nolimits_{i\in I}^{w}G_i$的一个正规子群。
\end{enumerate}
\end{theorem}

\begin{purplenote}
上述 $\tau_k$ 称为标准嵌入。
\end{purplenote}

\begin{theorem}\label{thm:weak-direct-sum-universal-property}
设 $\{A_i\mid i\in I\}$ 是一族加法记号的阿贝尔群，$B$ 也是阿贝尔群，且对每个 $i\in I$ 给定群同态 $\psi_i:A_i\to B$。则存在唯一群同态
\[
  \psi:\sum_{i\in I}A_i\longrightarrow B
\]
使得
\[
  \psi\circ\tau_i=\psi_i\qquad(i\in I),
\]
其中 $\tau_i:A_i\to\sum_{i\in I}A_i$ 是标准嵌入。因此，外直和$\sum_{i\in I}A_i$是阿贝尔群范畴中的余积。
\end{theorem}

\begin{purpleproof}
每个 $x=(a_i)_{i\in I}\in\sum_{i\in I}A_i$ 只有有限多个非零坐标，且可写成有限和 $x=\sum_i\tau_i(a_i)$。定义
\[
  \psi(x)=\sum_i\psi_i(a_i),
\]
右侧也是有限和。此定义给出群同态，并满足 $\psi\circ\tau_i=\psi_i$。若另有群同态满足同一条件，则它在每个 $\tau_i(a_i)$ 上的取值都被确定，因而在有限和 $x$ 上的取值也被确定，故该同态唯一。
\end{purpleproof}

\begin{purplenote}
这里的阿贝尔条件不能省略：非阿贝尔群的外弱直积一般不是群范畴中的余积。按手写笔记中的思路，取
\[
  A_1=A_2=B=\mathrm{GL}_2(\mathbb R),
  \qquad \psi_1=\psi_2=\operatorname{id}_{\mathrm{GL}_2(\mathbb R)}.
\]
这里 $\mathrm{GL}_2(\mathbb R)$ 表示所有可逆 $2\times2$ 实矩阵在矩阵乘法下构成的群。
若存在同态 $\psi:A_1\times A_2\to B$ 延拓这两个同态，则在直积中
\[
  (g_1,e_2)(e_1,g_2)=(e_1,g_2)(g_1,e_2),
\]
故它们的像也必须交换，即
\[
  g_1g_2=\psi_1(g_1)\psi_2(g_2)
  =\psi_2(g_2)\psi_1(g_1)=g_2g_1.
\]
但 $\mathrm{GL}_2(\mathbb R)$ 中并非所有矩阵都彼此交换，例如
\[
  \begin{pmatrix}1&1\\0&1\end{pmatrix}
  \begin{pmatrix}1&0\\1&1\end{pmatrix}
  \neq
  \begin{pmatrix}1&0\\1&1\end{pmatrix}
  \begin{pmatrix}1&1\\0&1\end{pmatrix}.
\]
因此这样的延拓 $\psi$ 不存在，这说明外弱直积一般不满足群范畴中余积的泛性质。

下面讨论一个群何时是其子群族的内弱直积。
\end{purplenote}

\begin{theorem}\label{thm:internal-weak-direct-product}
设 $\{N_i\mid i\in I\}$ 是群 $G$ 的一族正规子群。若
\begin{enumerate}[label=(\roman*)]
  \item $G=\left\langle\bigcup_{i\in I}N_i\right\rangle$；
  \item 对每个 $k\in I$，
  \[
    N_k\cap\left\langle\bigcup_{i\neq k}N_i\right\rangle=\{e\},
  \]
\end{enumerate}
则 $G$ 同构于 $\{N_i\mid i\in I\}$ 的外弱直积：
\[
  G\cong\prod\nolimits_{i\in I}^{w}N_i.
\]
\end{theorem}

\begin{purpleproof}
若 $i\neq j$，则条件 (ii) 给出 $N_i\cap N_j=\{e\}$。又因 $N_i,N_j\trianglelefteq G$，由定理~\ref{thm:normal-product} 的第 (4) 条，$N_i$ 与 $N_j$ 的元素彼此交换。

定义
\[
  \Phi:\prod\nolimits_{i\in I}^{w}N_i\longrightarrow G,
  \qquad (n_i)_{i\in I}\longmapsto\prod_{i\in I}n_i,
\]
其中右侧只有有限多个非单位元因子，且因各因子彼此交换，乘积与排列次序无关。条件 (i) 保证 $\Phi$ 满射；逐坐标相乘及因子间的交换性保证 $\Phi$ 是群同态。

若 $\Phi((n_i))=e$，则对每个 $k$，
\[
  n_k=\left(\prod_{i\neq k}n_i\right)^{-1}
  \in N_k\cap\left\langle\bigcup_{i\neq k}N_i\right\rangle=\{e\}.
\]
所以所有 $n_k=e$，$\Phi$ 是单同态，因而为同构。
\end{purpleproof}

\begin{corollary}\label{cor:finite-internal-direct-product}
设 $N_1,\ldots,N_r\trianglelefteq G$，且
\[
  G=N_1N_2\cdots N_r.
\]
若对每个 $1\leq k\leq r$，
\[
  N_k\cap N_1\cdots N_{k-1}N_{k+1}\cdots N_r=\{e\},
\]
则
\[
  G\cong N_1\times N_2\times\cdots\times N_r.
\]
\end{corollary}

\begin{definition}[内弱直积与内直和]\label{def:internal-weak-direct-product}
设 $\{N_i\mid i\in I\}$ 是群 $G$ 的一族正规子群。若
\[
  G=\left\langle\bigcup_{i\in I}N_i\right\rangle,
  \qquad
  N_k\cap\left\langle\bigcup_{i\neq k}N_i\right\rangle=\{e\}
  \quad(k\in I),
\]
则称 $G$ 是 $\{N_i\mid i\in I\}$ 的\textbf{内弱直积}。若 $G$ 是加法记号的阿贝尔群，也称为 $G$ 的\textbf{内直和}。
\end{definition}

\begin{theorem}\label{thm:unique-internal-weak-product-decomposition}
设 $\{N_i\mid i\in I\}$ 是群 $G$ 的一族正规子群。则 $G$ 是 $\{N_i\mid i\in I\}$ 的内弱直积，当且仅当每个非单位元$x\in G$ 都能唯一写成有限个非单位元因子的乘积
\[
  x=n_{i_1}n_{i_2}\cdots n_{i_r},
  \qquad n_{i_j}\in N_{i_j}\setminus\{e\},
\]
其中指标 $i_1,\ldots,i_r\in I$ 两两不同。这里的唯一性是指各坐标上的因子唯一；因子彼此交换，排列次序不计。
\end{theorem}

\begin{purplenote}
通常将内弱直积记作 $G=\prod\nolimits_{i\in I}^{w}N_i$。严格地说，$G$ 与外弱直积是同构，而不一定是作为集合相等；此记号是为了方便，表示两者在所讨论的直积结构下对应。
\end{purplenote}

\begin{theorem}\label{thm:weak-product-homomorphism}
设 $\{f_i:G_i\to H_i\mid i\in I\}$ 是一族群同态。映射
\[
  f\stackrel{\text{def}}{=}\prod\nolimits_{i\in I}^{w}f_i:
  \prod\nolimits_{i\in I}^{w}G_i\longrightarrow\prod\nolimits_{i\in I}^{w}H_i,
  \qquad (g_i)_{i\in I}\longmapsto(f_i(g_i))_{i\in I}
\]
是群同态，并且
\[
  \ker f=\prod\nolimits_{i\in I}^{w}\ker f_i,
  \qquad
  \operatorname{im}f=\prod\nolimits_{i\in I}^{w}\operatorname{im}f_i.
\]
进一步，$f$ 是单同态（满同态），当且仅当每个 $f_i$ 都是单同态（满同态）。
\end{theorem}

\begin{corollary}\label{cor:quotient-direct-products}
设 $\{G_i\mid i\in I\}$ 是一族群，且 $N_i\trianglelefteq G_i$（$i\in I$）。则：
\begin{enumerate}[label=(\arabic*)]
  \item $\prod_{i\in I}N_i\trianglelefteq\prod_{i\in I}G_i$，且
  \[
    \left(\prod_{i\in I}G_i\right)\Big/\left(\prod_{i\in I}N_i\right)
    \cong\prod_{i\in I}(G_i/N_i);
  \]
  \item $\prod\nolimits_{i\in I}^{w}N_i\trianglelefteq\prod\nolimits_{i\in I}^{w}G_i$，且
  \[
    \left(\prod\nolimits_{i\in I}^{w}G_i\right)\Big/\left(\prod\nolimits_{i\in I}^{w}N_i\right)
    \cong\prod\nolimits_{i\in I}^{w}(G_i/N_i).
  \]
\end{enumerate}
\end{corollary}

\begin{purpleproof}
(1) 对每个 $i\in I$，令 $\pi_i:G_i\to G_i/N_i$ 为标准投影。考虑同态
\[
  \pi=\prod_{i\in I}\pi_i:
  \prod_{i\in I}G_i\longrightarrow\prod_{i\in I}(G_i/N_i),
  \qquad (g_i)_{i\in I}\longmapsto(g_iN_i)_{i\in I}.
\]
由定理~\ref{thm:weak-product-homomorphism} 的坐标映射结论（对完全直积同理），$\pi$ 是满同态，且
\[
  \ker\pi=\prod_{i\in I}\ker\pi_i=\prod_{i\in I}N_i,
  \qquad
  \operatorname{im}\pi=\prod_{i\in I}\operatorname{im}\pi_i
  =\prod_{i\in I}(G_i/N_i).
\]
由\hyperref[cor:first-isomorphism]{第一同构定理}，得到第 (1) 条的同构。

(2) 将 $\pi$ 限制在外弱直积上，得到
\[
  \pi:\prod\nolimits_{i\in I}^{w}G_i
  \longrightarrow\prod\nolimits_{i\in I}^{w}(G_i/N_i),
  \qquad (g_i)_{i\in I}\longmapsto(g_iN_i)_{i\in I}.
\]
由定理~\ref{thm:weak-product-homomorphism}，该映射是满同态，且
\[
  \ker\pi=\prod\nolimits_{i\in I}^{w}\ker\pi_i
  =\prod\nolimits_{i\in I}^{w}N_i.
\]
再由\hyperref[cor:first-isomorphism]{第一同构定理}，得到第 (2) 条的同构。
\end{purpleproof}

\begin{proposition}\label{prop:quotient-of-internal-direct-product}
若群 $G$ 是其正规子群 $H$ 与 $K$ 的内直积，则
\[
  G/H\cong K,
  \qquad
  G/K\cong H.
\]
\end{proposition}

\begin{proposition}\label{prop:weak-product-internal-sum}
设 $\{G_i\mid i\in I\}$ 是一族群，且 $\tau_i:G_i\to\prod\nolimits_{i\in I}^{w}G_i$ 是标准嵌入。则外弱直积 $\prod\nolimits_{i\in I}^{w}G_i$ 是其正规子群族 $\{\tau_i(G_i)\mid i\in I\}$ 的内弱直积。
\end{proposition}

\begin{proposition}\label{prop:internal-weak-product-criterion}
设 $\{N_i\mid i\in I\}$ 是群 $G$ 的一族子群。则 $G$ 是这族子群的内弱直积，当且仅当：
\begin{enumerate}[label=(\arabic*)]
  \item 对任意 $i\neq j$、$a_i\in N_i$、$a_j\in N_j$，有 $a_ia_j=a_ja_i$；
  \item 每个非单位元$x\in G$ 都能唯一写成有限个非单位元因子的乘积
  \[
    x=a_{i_1}a_{i_2}\cdots a_{i_r},
    \qquad a_{i_j}\in N_{i_j}\setminus\{e\},
  \]
  其中指标 $i_1,\ldots,i_r\in I$ 两两不同；唯一性按各坐标上的因子理解，不计排列次序。
\end{enumerate}
\end{proposition}

\begin{purplenote}
注意与定理~\ref{thm:unique-internal-weak-product-decomposition} 对照：此处没有预先假设各 $N_i$ 正规；由两两交换及唯一分解条件，可以推出每个 $N_i\trianglelefteq G$。
\end{purplenote}

% 新增：第 1.9 节独立维护，整理自由群、群表示与自由积。
% !TeX root = ../main.tex
\subsection{自由群、群表示与自由积}\label{sec:free-groups}

% 来源：微信图片_20261111163357_47_215.jpg 至微信图片_20261111163413_53_215.jpg，1.9.1--1.9.12。
% 排版约定：黑笔定义、定理、命题和证明用普通环境；原紫笔说明保留为紫色。
\begin{purplenote}
本节将说明：给定一个集合，可以构造以它为自由生成集的群；给定一族群，也可以构造它们的自由积即群范畴中自由对象(自由群)是存在的。同时将推出一种利用生成元和生成元之间的关系来描述群的方法。
\end{purplenote}

\begin{definition}[自由群的既约字模型]\label{def:free-group-reduced-words}
设 $X$ 是非空集合。取一个与 $X$ 不相交的集合
\[
  X^{-1}=\{x^{-1}\mid x\in X\},
\]
并令 $x\mapsto x^{-1}$ 是 $X$ 与 $X^{-1}$ 之间的双射(这里包含了$|X|=|X^{-1}|$的条件)；再取一个不属于 $X\cup X^{-1}$ 的元素为单位字母记为 $1$。称由 $X\cup X^{-1}\cup\{1\}$ 中元素组成、且除有限项外都等于 $1$ 的序列
\[
  (a_1,a_2,\ldots)
\]
为一个\textbf{字}。若它满足以下条件，则称为\textbf{既约字}：
\begin{enumerate}[label=(\arabic*)]
  \item 对每个 $x\in X$，字中 $x$ 与 $x^{-1}$ 不相邻出现；
  \item 若 $a_k=1$，则对每个 $i\geq k$ 都有 $a_i=1$。
\end{enumerate}
特别地，$(1,1,\ldots)$ 称为空字，仍记作 $1$。可以发现空字是既约的。
\end{definition}

\begin{purplenote}
这里的 $X^{-1}$ 只是与 $X$ 不相交的形式副本；具体取成什么集合并不重要。后文主要使用字的形式，不依赖其底层集合实现。
\end{purplenote}

\begin{purplenote}
将所有既约字组成的集合记作 $F(X)$。由字的定义，每个非空既约字都可唯一写成有限序列
\[
  (x_1^{\varepsilon_1},\ldots,x_n^{\varepsilon_n},1,1,\ldots),
  \qquad n\in\mathbb N^+,\quad x_i\in X,\quad \varepsilon_i\in\{1,-1\}.
\]
我们定义两个这样的非空既约字相等，当且仅当它们的长度相同，且对应位置上的 $x_i$ 与指数 $\varepsilon_i$ 分别相同。因此映射
\[
  \tau:X\longrightarrow F(X),\qquad x\longmapsto(x,1,1,\ldots)
\]
是单射。后文常把 $X$ 与其像 $\tau(X)$ 视为同一集合，并简写为 $X\subseteq F(X)$。对任意整数 $m\geq1$， $m$ 个 $x^{-1}$ 的连续出现可简写为$x^{-m}$。
\end{purplenote}

\begin{purplenote}
下面具体说明 $F(X)$ 上定义乘法二元运算的“拼接和消去”规则。空字 $1$ 作为单位元，即对每个 $w\in F(X)$，有 $1w=w1=w$。设两个非空既约字为
\[
  w=(x_1^{\varepsilon_1},\ldots,x_m^{\varepsilon_m}),\qquad
  v=(y_1^{\delta_1},\ldots,y_n^{\delta_n}).
\]
令 $k$ 为可以消去的相邻互逆字母对的最大个数，即满足
\[
  x_{m-j}^{\varepsilon_{m-j}}=y_{j+1}^{-\delta_{j+1}}
  \qquad(j=0,1,\ldots,k-1)
\]
的最大整数 $0\leq k\leq\min(m,n)$。定义 $wv$ 为消去这 $k$ 对字母后剩下的字：
\[
  wv=
  \begin{cases}
    (x_1^{\varepsilon_1},\ldots,x_{m-k}^{\varepsilon_{m-k}},
      y_{k+1}^{\delta_{k+1}},\ldots,y_n^{\delta_n}),
      & k<m,\ k<n,\\[2pt]
    (y_{m+1}^{\delta_{m+1}},\ldots,y_n^{\delta_n}),
      & k=m<n,\\[2pt]
    (x_1^{\varepsilon_1},\ldots,x_{m-n}^{\varepsilon_{m-n}}),
      & k=n<m,\\[2pt]
    1, & k=m=n.
  \end{cases}
\]
由于各情形分别对应 $k<m$ 且 $k<n$、$k=m<n$、$k=n<m$、$k=m=n$，第一种情形中的两段都非空。由于 $w$ 与 $v$ 各自已经既约，拼接后可能出现的互逆相邻字母只会出现在两字的连接处，因此按此规则消去后所得仍是既约字。
\end{purplenote}

\begin{theorem}\label{thm:free-group-existence}
设 $X$ 是非空集合，$F(X)$ 为 $X$ 上所有既约字构成的集合。由以上讨论定义$F(X)$上的二元运算: 将两个既约字连接后，把相邻的互逆字母 $xx^{-1}$ 或 $x^{-1}x$ 逐次消去，所得既约字定义为它们的乘积。则 $F(X)$ 在此运算下是群，且$F(X)=\langle X\rangle$。称 $F(X)$ 为 $X$ 上的\textbf{自由群}。
\end{theorem}

\begin{proof}
空字 $1$ 是单位元；对既约字
\[
  w=(x_1^{\varepsilon_1},\ldots,x_n^{\varepsilon_n}),
\]
其逆字为
\[
  w^{-1}=(x_n^{-\varepsilon_n},\ldots,x_1^{-\varepsilon_1}),
\]
由乘法的消去规则，$ww^{-1}=w^{-1}w=1$。乘积仍是既约字，故运算封闭。下面验证结合律。

对每个 $x\in X$ 及 $\varepsilon\in\{1,-1\}$，定义左乘映射 $L_{x^{\varepsilon}}:F(X)\to F(X)$，令 $L_{x^{\varepsilon}}(1)=(x^{\varepsilon})$；对非空字 $w=(x_1^{\delta_1},\ldots,x_n^{\delta_n})$，令
\[
  L_{x^{\varepsilon}}(w)=
  \begin{cases}
    (x_2^{\delta_2},\ldots,x_n^{\delta_n}),
      & x^{\varepsilon}=x_1^{-\delta_1},\\[2pt]
    (x^{\varepsilon},x_1^{\delta_1},\ldots,x_n^{\delta_n}),
      & x^{\varepsilon}\neq x_1^{-\delta_1}.
  \end{cases}
\]
第一种情形若 $n=1$，右侧为空字。由定义可知
\[
  L_{x^{\varepsilon}}\circ L_{x^{-\varepsilon}}
  =\operatorname{id}_{F(X)}
  =L_{x^{-\varepsilon}}\circ L_{x^{\varepsilon}},
\]
所以每个 $L_{x^{\varepsilon}}$ 都是 $F(X)$ 上的置换。记 $\operatorname{Sym}(F(X))$ 为集合 $F(X)$ 上所有置换（即双射）在复合运算下构成的群，也就是笔记中记作 $A(F)$ 的群。令
\[
  F_0=\langle L_{x^{\varepsilon}}\mid x\in X,\ \varepsilon\in\{1,-1\}\rangle
  \leq\operatorname{Sym}(F(X)).
\]
定义映射 $\Phi:F(X)\to F_0$：
\[
  \Phi(1)=\operatorname{id}_{F(X)},\qquad
  \Phi(w)=L_{x_1^{\varepsilon_1}}\circ\cdots\circ L_{x_n^{\varepsilon_n}}
\]
（当 $w=(x_1^{\varepsilon_1},\ldots,x_n^{\varepsilon_n})$）。拼接并消去互逆字母对，对应于置换复合中消去 $L_{x^{\varepsilon}}$ 与 $L_{x^{-\varepsilon}}$，故
\[
  \Phi(wv)=\Phi(w)\circ\Phi(v)\qquad(w,v\in F(X)).
\]
显然$\Phi$ 是满射。又因 $\Phi(w)(1)=w$，故对$w_1 \neq w_2$ 有$\Phi(w_1) \neq \Phi(w_2)$ 因此$\Phi$ 是单射，故 $\Phi$ 为双射。又因为$F_0$是群，对任意 $u,v,w\in F(X)$，
\[
  \Phi((uv)w)
  =(\Phi(u)\circ\Phi(v))\circ\Phi(w)
  =\Phi(u)\circ(\Phi(v)\circ\Phi(w))
  =\Phi(u(vw)).
\]
由 $\Phi$ 的单射性，$(uv)w=u(vw)$，故乘法满足结合律。最后，每个既约字都是 $X$ 中单字母及其逆的乘积，所以 $F(X)=\langle X\rangle$。
\end{proof}

\begin{purplenote}
上述消去过程给出唯一的既约字；乘积就是连接后逐次消去相邻互逆字母所得的字。单位元为空字，逆字由原字中字母倒序并逐个取逆得到。若 $|X|\geq2$，则 $F(X)$ 非阿贝尔；此外，$F(X)$ 中每个非单位元都有无限阶。
\end{purplenote}

\begin{theorem}[自由群的泛性质]\label{thm:free-group-universal-property}
设 $X$ 是非空集合，$F(X)$ 是 $X$ 上的自由群，且 $\tau:X\to F(X)$ 是标准嵌入将 $x$ 映为对应的单字母字。对任意群 $G$ 及任意集合映射 $f:X\to G$，存在唯一群同态
\[
  \overline f:F(X)\longrightarrow G
\]
使得 $\overline f\circ\tau=f$。因此，在群范畴中，$F(X)$ 是集合 $X$ 上的一个自由对象。
\end{theorem}

\begin{proof}
对既约字 $w=(x_1^{\varepsilon_1},\ldots,x_n^{\varepsilon_n})$，其中 $x_i\in X$、$\varepsilon_i\in\{1,-1\}$，定义
\[
  \overline f(w)=f(x_1)^{\varepsilon_1}\cdots f(x_n)^{\varepsilon_n},
  \qquad \overline f(1)=e_G.
\]
每个元素都有唯一的既约字表示，故 $\overline f$ 定义良好。由于 $G$ 是群，右侧的乘积总有定义。消去相邻互逆字母不改变右侧的群元素，因此 $\overline f$ 与字的乘法相容，是群同态，且满足 $\overline f\circ\tau=f$。

下面证明唯一性。若 $g:F(X)\to G$ 是群同态且 $g\circ\tau=f$，则对任意既约字 $w=(x_1^{\varepsilon_1},\ldots,x_n^{\varepsilon_n})$，
\begin{align*}
  g(w)
  &=g(x_1^{\varepsilon_1})\cdots g(x_n^{\varepsilon_n})\\
  &=g(\tau(x_1))^{\varepsilon_1}\cdots g(\tau(x_n))^{\varepsilon_n}\\
  &=f(x_1)^{\varepsilon_1}\cdots f(x_n)^{\varepsilon_n}
   =\overline f(w).
\end{align*}
所以 $g=\overline f$，延拓唯一。
\end{proof}

\begin{corollary}\label{cor:every-group-quotient-of-free}
任意群 $G$ 都是某个自由群在群同态下的像。
\end{corollary}

\begin{purpleproof}
取 $G$ 的一个生成集 $X$，即 $G=\langle X\rangle$，并令 $f:X\hookrightarrow G$ 为包含映射。由定理~\ref{thm:free-group-universal-property}，存在唯一群同态 $\overline f:F(X)\to G$ 使得 $\overline f\circ\tau=f$。因为 $\overline f$ 的像包含 $X$，而 $X$ 生成 $G$，所以 $\overline f$ 是满同态，故 $G=\operatorname{im}\overline f$。

需要注意，即使 $G=\langle X\rangle$，自由群 $F(X)$ 中的 $\tau(x)^{-1}$ 与 $G$ 中的 $f(x)^{-1}$ 是不同群中的元素，不应直接视为同一个元素；同态所保证的是
\[
  \overline f\bigl(\tau(x)^{-1}\bigr)=f(x)^{-1}\qquad(x\in X).
\]
\end{purpleproof}

\begin{definition}[群的表示]\label{def:group-presentation}
设 $X$ 是非空集合，$Y$ 是自由群 $F(X)$ 中一些既约字组成的集合。令 $N$ 为 $F(X)$ 中包含 $Y$ 的最小正规子群，即由 $Y$ 正规生成的子群。若群 $G$ 同构于 $F(X)/N$，则称 $G$ 是由生成元 $X$ 和关系 $Y$ 定义的群，或称 $G$ 有表示
\[
  G\cong\langle X\mid Y\rangle.
\]
此时也称 $(X\mid Y)$ 是 $G$ 的一个\textbf{群表示}。
\end{definition}

\begin{purplenote}
群的表示通常不唯一。严格地说，若 $\pi:F(X)\to F(X)/N$ 是自然投影，则关系 $w\in Y$ 的含义是 $\pi(w)=N$，即它在商群中成为单位元。由于 $F(X)=\langle X\rangle$，故商群 $F(X)/N$ 由 $X$ 在标准投影下的像生成。
\end{purplenote}

\begin{theorem}[范·戴克定理]\label{thm:van-dyck}
设 $N$ 是 $F(X)$ 中由 $Y$ 正规生成的子群，并将 $G=F(X)/N$ 记作 $G=\langle X\mid Y\rangle$。若群 $H$ 中给定元素族 $\{h_x\mid x\in X\}$ 满足关系 $Y$，即将 $x$ 映为 $h_x$ 所诱导的同态(自由对象 $F(X)$ 的泛性质所得的) $\overline f:F(X)\to H$ 满足 $\overline f(y)=e_H$（$y\in Y$），则存在唯一群同态
\[
  \varphi:G\longrightarrow H
\]
使得 $\varphi(xN)=h_x$（$x\in X$）。若 $\{h_x\mid x\in X\}$ 生成 $H$，则 $\varphi$ 为满同态。
\end{theorem}

\begin{proof}
由假设，$\overline f(y)=e_H$ 对每个 $y\in Y$ 成立。因此 $Y\subseteq\ker\overline f$。由于 $\ker\overline f$ 是正规子群，而 $N$ 是包含 $Y$ 的最小正规子群，故
\[
  N\subseteq\ker\overline f.
\]
定义
\[
  \varphi:F(X)/N\longrightarrow H,\qquad
  \varphi(wN)=\overline f(w)\quad(w\in F(X)).
\]
若 $wN=w'N$，则 $w^{-1}w'\in N\subseteq\ker\overline f$，从而 $\overline f(w)=\overline f(w')$；故 $\varphi$ 定义良好。由 $\overline f$ 是群同态可知 $\varphi$ 也是群同态，并且
\[
  \varphi(xN)=\overline f(x)=h_x\qquad(x\in X).
\]

商群 $F(X)/N$ 由 $\{xN\mid x\in X\}$ 生成，所以任何满足上述条件的群同态都必须与 $\varphi$ 相同，故延拓唯一。若 $\{h_x\mid x\in X\}$ 生成 $H$，则 $\operatorname{im}\varphi$ 含有这组生成元，因而 $\varphi$ 是满同态。
\end{proof}

\begin{purplenote}
这里一般得到的是满同态（当指定元素生成 $H$ 时），而不一定是同构；$H$ 中还可能满足表示中未列出的其他关系。
\end{purplenote}

\begin{example}\label{ex:quaternion-presentation}
四元数群 $Q_8$ 有表示
\[
  Q_8\cong\langle a,b\mid a^4=e,\ a^2=b^2,\ bab^{-1}=a^{-1}\rangle.
\]
\begin{proof}
记 $F=F(\{a,b\})$，令 $N$ 为 $F$ 中由 $a^4$、$a^2b^{-2}$、$abab^{-1}$ 正规生成的子群，并记 $P=F/N$。这三个关系分别等价于显示的 $a^4=e$、$a^2=b^2$、$bab^{-1}=a^{-1}$。在 $P$ 中，
\[
  b^2=a^2,\qquad ba=a^{-1}b.
\]
此外，$a^{-1}=a^3$，且 $b^4=(b^2)^2=a^4=e$，故 $b^{-1}=b^3$。因此任意元素都可先写成由 $a,b$ 构成的字。利用 $ba=a^{-1}b$（从 $bab^{-1}=a^{-1}$ 两边右乘 $b$ 得到），可将字中的 $b$ 逐个移到 $a$ 的右侧，写成 $a^m b^n$ 的形式；再用 $a^4=e$ 和 $b^2=a^2$ 化简，得到 $a^i$ 或 $a^ib$，其中 $0\leq i<4$。故 $|P|\leq8$。

令 $f:\{a,b\}\to Q_8$ 满足 $f(a)=i$、$f(b)=j$。由自由群的泛性质，$f$ 延拓为群同态 $\overline f:F\to Q_8$。由于 $i^4=e$、$i^2=j^2=-1$ 且 $jij^{-1}=-i$，有
\[
  \overline f(a^4)=e,\qquad
  \overline f(a^2b^{-2})=i^2j^{-2}=e,\qquad
  \overline f(abab^{-1})=i(jij^{-1})=e.
\]
所以 $\ker\overline f$ 是包含这三个关系字的正规子群，从而 $N\subseteq\ker\overline f$。定义
\[
  \varphi:P\longrightarrow Q_8,\qquad
  \varphi(wN)=\overline f(w)\quad(w\in F).
\]
若 $wN=w'N$，则 $w^{-1}w'\in N\subseteq\ker\overline f$，故 $\overline f(w)=\overline f(w')$，因此 $\varphi$ 定义良好；由定义它是群同态，且 $\varphi(aN)=i$、$\varphi(bN)=j$。因 $i,j$ 生成 $Q_8$，$\varphi$ 是满同态。又 $|P|\leq8=|Q_8|$，故 $\varphi$ 为同构。
\end{proof}
\end{example}

\begin{purplenote}
在群表示中，常把诸如 $x^{-1}x^{-1}y^{-1}y^{-1}y^{-1}$ 的字简写为幂的形式$x^{-2}y^{-3}$，而不逐个写出 $x^{-1}$、$y^{-1}$。
\end{purplenote}

\begin{purplenote}
下面讨论群范畴中的余积，即自由积。它的构造与自由群类似：用各因子群中的非单位元作字母，再合并相邻且属于同一因子群的字母，并消去乘积为单位元的情形。
\end{purplenote}

\begin{definition}[自由积]\label{def:free-product}
设 $\{G_i\mid i\in I\}$ 是一族两两不交的群，各群的单位元分别记为 $e_i$。另取一个不属于 $\bigcup_{i\in I}G_i$ 的符号 $1$，并记字母集合
\[
  \mathcal A=\left(\bigcup_{i\in I}G_i\right)\cup\{1\}.
\]
一个\textbf{字}是一个序列
\[
  (a_1,a_2,\ldots),\qquad a_k\in\mathcal A,
\]
且除有限项外都等于 $1$。全由 $1$ 组成的序列称为空字，仍记为 $1$。
若一个字满足以下三个条件，则称为\textbf{既约字}：
\begin{enumerate}[label=(\arabic*)]
  \item 若 $a_k=1$，则对每个 $j\geq k$ 都有 $a_j=1$；
  \item 若 $a_k\in G_i$，则 $a_k\neq e_i$；
  \item 若相邻两项 $a_k\in G_i$、$a_{k+1}\in G_j$，则 $i\neq j$(即相邻两项不属于同一个群)。
\end{enumerate}
所以既约字的有限有效部分不含任何因子群单位元 $e_i$，末尾可能出现的 $1$ 只是补齐项。

约化一个普通字时，先删去所有等于标记 $1$ 或某个因子群单位元 $e_i$ 的项，再从左到右依次读取剩余字母，并维护一个既约字。读到 $a\in G_i\setminus\{e_i\}$ 时，若当前字为空，或其末字母属于 $G_j$ 且 $j\neq i$，就把 $a$ 接到末尾；若末字母 $b$ 也属于 $G_i$，就在 $G_i$ 内计算 $ba$，并以 $ba$ 替换末字母 $b$；若 $ba=e_i$，则删去末字母 $b$。处理完后在所得有限序列末尾补上无限多个 $1$。约化过程是确定的，故所得既约字唯一。

两个既约字 $u,v$ 的乘积定义为先去掉它们末尾补齐的 $1$，连接其有限有效部分，再按上述规则约化所得的既约字。空字 $1$ 是单位元；若 $u=(a_1,\ldots,a_n,1,1,\ldots)$ 且 $a_n\neq1$，则其逆字为
\[
  u^{-1}=(a_n^{-1},\ldots,a_1^{-1},1,1,\ldots),
\]
其中每个逆元在相应的因子群中取；空字的逆仍为空字。所有既约字在此乘法下构成群，称为这族群的\textbf{自由积}，记作
\[
  \prod_{i\in I}{}^{*}G_i\quad\text{（通常也记作 }\mathop{*}_{i\in I}G_i\text{）}.
\]
对每个 $k\in I$，定义标准嵌入
\[
  \tau_k:G_k\longrightarrow\prod_{i\in I}{}^{*}G_i
\]
为 $\tau_k(e_k)=1$，并将 $a\in G_k\setminus\{e_k\}$ 映为 $(a,1,1,\ldots)$。易知$\tau_k$ 是群单同态，以后视$G_k$和$\tau_k(G_k)$为一个东西。
\end{definition}

\begin{theorem}[自由积的泛性质]\label{thm:free-product-universal-property}
设 $\{G_i\mid i\in I\}$ 是一族两两不交的群，$H$ 是群，且对每个 $i\in I$ 给定群同态 $\psi_i:G_i\to H$。则存在唯一群同态
\[
  \psi:\prod_{i\in I}{}^{*}G_i\longrightarrow H
\]
使得 $\psi\circ\tau_i=\psi_i$（$i\in I$）。因此，$\prod_{i\in I}{}^{*}G_i$ 是群范畴中这族群的余积。
\end{theorem}

\begin{proof}
对既约字 $(a_1,\ldots,a_n)$，令 $a_k\in G_{i_k}$，定义
\[
  \psi(a_1,\ldots,a_n)=\psi_{i_1}(a_1)\cdots\psi_{i_n}(a_n),
  \qquad \psi(1)=e_H.
\]
在某个因子群内合并相邻字母或删去单位元，不改变右侧的值，故这一定义与合并、消去过程相容并给出群同态。它在每个标准嵌入上的限制都是 $\psi_i$。自由积由各因子的像生成，因此这样的同态唯一。
\end{proof}

\begin{purplenote}
以下列出一些群表示的习题结论；其中的验证可参照例~\ref{ex:quaternion-presentation} 的方法。
\end{purplenote}

\begin{proposition}\label{prop:free-group-basic-facts}
设 $X$ 是非空集合，$F(X)$ 是 $X$ 上的自由群。则：
\begin{enumerate}[label=(\arabic*)]
  \item $F(X)$ 中非单位元的元素的阶为无限；
  \item 若 $X=\{x\}$，则 $F(X)$ 是无限循环群，故 $F(X)\cong\mathbb Z$；
  \item 若 $Y\subseteq X$，则 $F(X)$ 除以 $Y$ 的正规闭包所得的商群是由 $X\setminus Y$ 自由生成的自由群（当 $X\setminus Y=\varnothing$ 时，该商群为平凡群）。
\end{enumerate}
\end{proposition}

\begin{proposition}\label{prop:group-presentations-examples}
下列群表示给出相应的群：
\begin{enumerate}[label=(\arabic*)]
  \item 六阶循环群有表示
  \[
    \mathbb Z_6\cong\langle a,b\mid a^2=e,\ b^3=e,\ a^{-1}b^{-1}ab=e\rangle;
  \]
  \item $n$ 阶二面体群（$n\geq3$）有表示
  \[
    D_n\cong\langle a,b\mid a^n=e,\ b^2=e,\ abab=e\rangle;
  \]
  \item 由一个生成元 $b$ 及关系 $b^m=e$（$m\in\mathbb N^+$）定义的群同构于 $\mathbb Z_m$；
  \item 由生成元 $a,b$ 及关系 $a^2=b^3
  =e$ 定义的群是无限非阿贝尔群。
\end{enumerate}
\end{proposition}

\begin{proposition}\label{prop:free-product-facts}
设 $G,H,K$ 是群，则：
\begin{enumerate}[label=(\arabic*)]
  \item 自由积满足交换律和结合律：$G*H\cong H*G, (G*H)*K\cong G*H*K\cong G*(H*K)$；
  \item 若 $N$ 是 $G*H$ 中包含标准嵌入像 $G$ 的最小正规子群，则 $(G*H)/N\cong H$；
  \item 若 $G,H$ 都非平凡，则 $G*H$ 是无限群，且其中心为平凡群$\langle e\rangle$；
  \item 由生成元 $a,b$ 及关系 $a^2=e$、$b^3=e$ 定义的群同构于 $\mathbb Z_2*\mathbb Z_3$；
  \item 若 $f:G_1\to G_2$、$g:H_1\to H_2$ 是群同态，则存在唯一群同态
  \[
    h:G_1*H_1\longrightarrow G_2*H_2
  \]
  使得 $h$ 在 $G_1$ 上的限制等于 $G_1\xrightarrow{f}G_2\hookrightarrow G_2*H_2$，在 $H_1$ 上的限制等于 $H_1\xrightarrow{g}H_2\hookrightarrow G_2*H_2$；这里的 $\hookrightarrow$ 表示相应自由积中的标准嵌入（故为单同态）。
\end{enumerate}
\end{proposition}

\begin{proof}
(1) 由\hyperref[def:free-product]{自由积的定义}可知交换律是显然的。关于结合律，只需证明 $(G*H)*K\cong G*(H*K)$；从既约字的形式表示即可证明。

(2) 考察 $N$ 中元素的具体形式。首先 $G\leq N$；其次，对任意 $h\in H$、$g\in G$，有 $hgh^{-1}\in N$。由此可猜想3，若
\[
  x=g_1h_1g_2h_2\cdots g_nh_n\in N,\qquad g_i\in G,\quad h_i\in H,
\]
则 $h_1h_2\cdots h_n=e_H$。因此，商群中的陪集由 $H$ 中相应字母的乘积决定，从而 $(G*H)/N\cong H$。

(4) 记 $\mathbb Z_2=\langle 0_2, 1_2\rangle$、$\mathbb Z_3=\langle 0_3, 1_3, 2_3\rangle$。考虑映射
\[
  f:X=\{a,b\}\longrightarrow\mathbb Z_2*\mathbb Z_3,\qquad f(a)=1_2,\quad f(b)=1_3.
\]
由自由群的泛性质，存在群同态 $\overline f:F(X)\to\mathbb Z_2*\mathbb Z_3$，使得 $\overline f\circ\tau=f$。又因 $\mathbb Z_2*\mathbb Z_3=\langle 1_2,1_3\rangle$，故 $\overline f$ 是满同态。令 $N$ 为由 $a^2,b^3$ 正规生成的子群。由于 $\overline f(a^2)=\overline f(b^3)=1$，有 $N\subseteq\ker\overline f$，故存在满同态
\[
  \widetilde f:F(X)/N\longrightarrow\mathbb Z_2*\mathbb Z_3,\qquad \widetilde f(wN)=\overline f(w).
\]
任取 $w\in F(X)$，将其写成 $a,b$ 的幂交替出现的形式。若 $w\in\ker\overline f$，则它在 $\mathbb Z_2*\mathbb Z_3$ 中约化为空字。约化过程中将 $\bar{f}(a)$ 的幂按模 $2$ 化简、将 $\bar{f}(b)$ 的幂按模 $3$ 化简，并合并相邻的同因子字母, 故易知 $w\in N$。于是 $\ker\overline f\subseteq N$，从而 $N=\ker\overline f$。由定理~\ref{thm:factor-through-quotient}，$\widetilde f$ 是同构。
\end{proof}


\section{群的结构}

% 新增：第 2.1 节独立维护，保留自由阿贝尔群原稿编号及黑、紫双色区分。
% !TeX root = ../main.tex
\subsection{自由阿贝尔群}

% 来源目录：../../NotePhotos/BaiduSyncdisk/algebra/2.1/（相对于本科目目录）。
% 共 7 张照片，文件名编号 32--38；编号不是原稿页码。
% 保留原稿 2.1.1--2.1.19 的编号；另页证明移至对应结论之后。
% 黑笔正文用普通环境，原紫笔说明用紫色；明确标出的“整理注”不属于原稿。

% 来源：微信图片_20261114000824_32_26.jpg，2.1.1--2.1.2。
\begin{proposition}[乘法记号与加法记号的对照]\label{prop:abelian-additive-notation}
以下是一般群与采用加法记号的阿贝尔群的符号对照表：
\[
  \begin{array}{c|c}
    \text{乘法记号} & \text{加法记号} \\ \hline
    ab & a+b \\
    a^{-1} & -a \\
    e & 0 \\
    a^n & na \\
    ab^{-1} & a-b \\
    HK & H+K \\
    aH & a+H \\
    G\times H & G\oplus H \\
    H\vee K & H+K \\
    \prod\nolimits_{i\in I}^{w}G_i & \sum_{i\in I}G_i \\
    \text{weak direct product} & \text{direct sum}
  \end{array}
\]
对任意用加法记号表示运算的群（不一定是阿贝尔群），有
\[
  (m+n)a=ma+na\qquad(a\in G,\ m,n\in\mathbb Z).
\]
若 $G$ 是阿贝尔群，则还有
\[
  m(a+b)=ma+mb\qquad(a,b\in G,\ m\in\mathbb Z).
\]
\end{proposition}

\begin{purplenote}
若 $G$ 是阿贝尔群，$H,K\leq G$，则 $H+K$ 是一个子群；用乘法记号说，这是因为 $HK=KH$。这里 $H\vee K$ 表示由 $H\cup K$ 生成的子群。
\end{purplenote}

\begin{definition}[阿贝尔群的基]\label{def:abelian-group-basis}
设 $G$ 是用加法记号表示运算的群，$X$ 是 $G$ 的一个非空子集，则
\[
  \langle X\rangle
  =\{n_1x_1+\cdots+n_kx_k\mid
      k\in\mathbb N^+,\ n_i\in\mathbb Z,\ x_i\in X\}.
\]
特别地，若 $X=\{x_0\}$，则
\[
  \langle X\rangle=\{nx_0\mid n\in\mathbb Z\}.
\]
现设 $F$ 是阿贝尔群，$X$ 是 $F$ 的一个非空子集。若 $F=\langle X\rangle$，并且对任意两两不同的 $x_1,\ldots,x_k\in X$ 及任意 $n_1,\ldots,n_k\in\mathbb Z$，都有
\[
  n_1x_1+\cdots+n_kx_k=0
  \quad\Longrightarrow\quad n_1=\cdots=n_k=0,
\]
则称 $X$ 是 $F$ 的一个\textbf{基}。
\end{definition}

\begin{purplenote}
后面需注意阿贝尔群中的基与向量空间中的基的区别与联系。
\end{purplenote}

% 来源：微信图片_20261114000831_33_26.jpg，2.1.3--2.1.8 的陈述。
\begin{theorem}\label{thm:free-abelian-characterizations}
以下关于阿贝尔群 $F$ 的描述是等价的：
\begin{enumerate}[label=(\arabic*)]
  \item $F$ 有一个非空的基；
  \item $F$ 是一族无限阶循环子群的内直和；
  \item $F$ 同构于一族整数群 $\mathbb Z$ 的外直和；
  \item 存在非空集合 $X$ 和映射 $\tau:X\to F$，满足：对任意给定的阿贝尔群 $G$ 及任意集合映射 $f:X\to G$，存在唯一群同态
  \[
    \overline f:F\longrightarrow G
  \]
  使得 $\overline f\circ\tau=f$。换句话说，$F$ 是阿贝尔群范畴中的一个自由对象。
\end{enumerate}
这里第 (2)、(3) 条的指标集均取非空集。
\end{theorem}

\begin{purplenote}
第 (4) 条中的 $X$ 不一定是 $F$ 的子集。
\end{purplenote}

% 来源：微信图片_20261114000835_34_26.jpg，2.1.3 的证明。
\begin{proof}
$(1)\Rightarrow(2)$：设 $X=\{x_i\mid i\in I\}$ 是 $F$ 的一个非空基，记 $N_i=\langle x_i\rangle$。由于 $F$ 是阿贝尔群，各 $N_i$ 都是 $F$ 的正规子群，且
\[
  F=\left\langle\bigcup_{i\in I}N_i\right\rangle.
\]
每个 $x_i$ 都有无限阶，否则有某个非零整数 $n$ 满足 $nx_i=0$，与\hyperref[def:abelian-group-basis]{基的定义}矛盾。若
\[
  z\in N_i\cap\left\langle\bigcup_{j\ne i}N_j\right\rangle,
\]
则存在有限多个不同于 $i$ 的指标 $j_1,\ldots,j_k$，使得
\[
  z=nx_i=n_1x_{j_1}+\cdots+n_kx_{j_k}.
\]
由基的线性无关性，各系数都为零，故 $z=0$。因此 $F$ 是这族无限阶循环子群的内直和。

$(2)\Rightarrow(3)$：无限阶循环群同构于整数群 $\mathbb Z$，再用定理~\ref{thm:internal-weak-direct-product} 即得。

$(3)\Rightarrow(1)$：先考察 $\sum_{i\in I}\mathbb Z_i$，其中每个 $\mathbb Z_i=\mathbb Z$。令 $u_i$ 的第 $i$ 个坐标为 $1$，其余坐标为 $0$，并记 $Y=\{u_i\mid i\in I\}$。易知
\[
  \sum_{i\in I}\mathbb Z_i=\langle Y\rangle,
\]
且 $Y$ 线性无关，故为一个非空基。通过给定的同构，把这个基移到 $F$ 中即可。

$(1)\Rightarrow(4)$：设 $X$ 是 $F$ 的基，取 $\tau:X\rightarrow F$ 为嵌入映射。对给定的 $G$ 和 $f$，定义
\[
  \overline f(n_1x_1+\cdots+n_kx_k)
  =n_1f(x_1)+\cdots+n_kf(x_k).
\]
把相同基元素的系数合并后，每个元素的有限线性表示唯一，故此定义良好。易验证 $\overline f\circ\tau=f$，且 $\overline f$ 是群同态。任何满足这一条件的同态在 $X$ 上都与 $f$ 相同，而 $X$ 生成 $F$，故同态唯一。

$(4)\Rightarrow(3)$：设 $X=\{x_i\mid i\in I\}$，取 $A=\sum_{i\in I}\mathbb Z_i$(其中$\mathbb Z_i \cong \mathbb Z, i\in I$)，及上面构造的基 $Y=\{u_i\mid i\in I\}$。由 $(1)\Rightarrow(4)$，$A$ 是 $Y$ 上的自由对象；由假设，$F$ 是 $X$ 上的自由对象。因为 $X$ 与 $Y$ 一一对应，由定理~\ref{thm:free-objects-equivalent} 得
\[
  F\cong A=\sum_{i\in I}\mathbb Z_i.
\]
\end{proof}

\begin{definition}[自由阿贝尔群]\label{def:free-abelian-group}
若阿贝尔群 $F$ 满足定理~\ref{thm:free-abelian-characterizations} 中的等价条件，则称 $F$ 为\textbf{自由阿贝尔群}；相对于第 (4) 条中的集合 $X$，也称为 $X$ 上的自由阿贝尔群。
\end{definition}

\begin{purplenote}
对于任意给定的非空集合 $X$，定理~\ref{thm:free-abelian-characterizations} 的证明展示了如何构造“以 $X$ 为基”的自由阿贝尔群：直接考虑 $\sum_{i\in I}\mathbb Z$，把 $X$ 看成指标集 $I$，并把 $x_i$ 与相应的标准基元素 $u_i$ 对应起来。
\end{purplenote}

\begin{theorem}\label{thm:free-abelian-rank}
若 $F$ 是一个自由阿贝尔群，则 $F$ 的任意两个基有相同的基数。因此，给定自由阿贝尔群的基的基数是固定的，称为该自由阿贝尔群的\textbf{秩}（rank）。
\end{theorem}

% 来源：微信图片_20261114000835_34_26.jpg，2.1.5 的证明概要。
\begin{proof}[证明概要（有限秩情形）]
设 $F$ 有一个基 $X$，$|X|=n<\infty$，则 $F\cong\sum_{i=1}^{n}\mathbb Z$。定义
\[
  2F=\{2x\mid x\in F\}.
\]
易知 $2F\trianglelefteq F$，且在上述同构下(同一个同构)对应于 $\sum_{i=1}^{n}2\mathbb Z$。由推论~\ref{cor:quotient-direct-products}，
\[
  F/2F\cong\sum_{i=1}^{n}(\mathbb Z/2\mathbb Z),
  \qquad |F/2F|=2^n.
\]
若 $Y$ 是 $F$ 的另一个基，对于某个任意给定的正整数$r$若有$|Y|\geq r$则由上面类似分析可得 $|F/2F| \geq 2^r$，因此有 $2^n \geq 2^r$，故 $r\leq n$。于是 $Y$ 为有限集，记 $|Y|=m$，再由同一商群的计算得到 $2^m=2^n$，所以 $m=n$。
\end{proof}

\begin{purplenote}
原稿注明：$X$ 为无限集时的证明见教材。
\end{purplenote}

\begin{proposition}\label{prop:free-abelian-isomorphism-rank}
设 $F_1$ 是以 $X_1$ 为基的自由阿贝尔群，$F_2$ 是以 $X_2$ 为基的自由阿贝尔群。则
\[
  F_1\cong F_2
  \quad\Longleftrightarrow\quad
  \operatorname{rank}F_1=\operatorname{rank}F_2
  \quad\Longleftrightarrow\quad |X_1|=|X_2|.
\]
\end{proposition}

\begin{theorem}\label{thm:abelian-quotient-of-free-abelian}
设 $G$ 是阿贝尔群，$X$ 是 $G$ 的一个生成集，则 $G$ 是某个秩为 $|X|$ 的自由阿贝尔群在某个群同态下的像。
\end{theorem}

\begin{purplenote}
以下介绍研究有限生成阿贝尔群结构时有用的工具。
\end{purplenote}

\begin{lemma}\label{lem:free-abelian-basis-shear}
若 $\{x_1,\ldots,x_n\}$ 是自由阿贝尔群 $F$ 的一个基，$a\in\mathbb Z$，则对所有 $i\ne j$，
\[
  \{x_1,\ldots,x_{j-1},\ x_j+ax_i,\ x_{j+1},\ldots,x_n\}
\]
也是 $F$ 的一个基。
\end{lemma}

% 来源：微信图片_20261114000839_35_26.jpg，2.1.9--2.1.12 的陈述及紫笔说明。
\begin{theorem}\label{thm:subgroup-finite-rank-free-abelian}
设 $F$ 是秩为 $n<\infty$ 的自由阿贝尔群，$G$ 是 $F$ 的一个非零子群。则存在 $F$ 的一个基 $\{x_1,\ldots,x_n\}$、一个正整数 $r$（$1\leq r\leq n$）及满足
\[
  d_1\mid d_2\mid\cdots\mid d_r
\]
的正整数 $d_1,\ldots,d_r$，使得 $G$ 是以 $\{d_1x_1,\ldots,d_rx_r\}$ 为基的自由阿贝尔群。
\end{theorem}

% 来源：微信图片_20261114000843_36_26.jpg，2.1.9 的完整证明。
% 整理：统一归纳步骤中 G∩H 的基的末项指标，允许 G∩H=0。
\begin{proof}
对 $n$ 作归纳。当 $n=1$ 时，结论由无限循环群的非零子群仍为无限循环群得到。设秩小于 $n$ 时结论成立，下面考虑秩为 $n$ 的情形。

记
\[
  S=\left\{s\in\mathbb Z\ \middle|\
    \begin{array}{l}
      \text{存在 }F\text{ 的一个基 }\{x_1,\ldots,x_n\}\text{ 及 }g\in G,\\
      g=sx_1+k_2x_2+\cdots+k_nx_n,\quad k_i\in\mathbb Z
    \end{array}\right\}.
\]
由于 $G\ne0$，且交换基元素位置仍是基及改变基元素符号仍是基，故 $S$ 中有正整数。取 $d_1$ 为 $S$ 中最小的正整数，选取相应的基和元素
\[
  g_0=d_1x_1+k_2x_2+\cdots+k_nx_n\in G.
\]
先证明 $d_1\mid k_i$（$2\leq i\leq n$）。否则，对某个 $i$ 可写成
\[
  k_i=q_id_1+r_i,\qquad 0<r_i<d_1.
\]
于是
\[
  g_0=d_1(x_1+q_ix_i)
      +\sum_{\substack{2\leq j\leq n\\j\ne i}}k_jx_j+r_ix_i.
\]
由引理~\ref{lem:free-abelian-basis-shear}，以 $x_1+q_ix_i$ 替换 $x_1$ 后仍得到 $F$ 的基。再把 $x_i$ 移到第一位，便有 $r_i\in S$，与 $d_1$ 的最小性矛盾。

因此可写 $k_i=d_1q_i$。令
\[
  v=x_1+q_2x_2+\cdots+q_nx_n,
\]
则 $d_1v=g_0\in G$。反复使用引理~\ref{lem:free-abelian-basis-shear}，可知 $\{v,x_2,\ldots,x_n\}$ 仍为 $F$ 的基。记 $H=\langle x_2,\ldots,x_n\rangle$，则
\[
  F=\langle v\rangle\oplus H.
\]
任取 $g\in G$，唯一写成
\[
  g=n_1v+n_2x_2+\cdots+n_nx_n.
\]
将 $n_1$ 除以 $d_1$，写作 $n_1=qd_1+r$，其中 $0\leq r<d_1$。因为 $d_1v\in G$，所以
\[
  g-qd_1v=rv+n_2x_2+\cdots+n_nx_n\in G.
\]
若 $r>0$，则 $r\in S$，矛盾。因此 $d_1\mid n_1$，从而
\[
  G=\langle d_1v\rangle\oplus(G\cap H).
\]

若 $G\cap H=0$，取 $r=1$ 即得结论。否则，$H$ 的秩为 $n-1$，由归纳假设，存在 $H$ 的一个基 $\{y_2,\ldots,y_n\}$ 及某个 $2\leq r\leq n$，使得
\[
  \{d_2y_2,\ldots,d_ry_r\}
\]
是 $G\cap H$ 的基，且 $d_2\mid\cdots\mid d_r$。于是 $\{v,y_2,\ldots,y_n\}$ 是 $F$ 的基，$\{d_1v,d_2y_2,\ldots,d_ry_r\}$ 是 $G$ 的基。对
\[
  d_1v+d_2y_2+\cdots+d_ry_r\in G
\]
使用前面相同的带余除法和换基论证，得到 $d_1\mid d_2$，故所需的整除链成立。
\end{proof}

\begin{purplenote}
这里归纳法最核心的技巧是 $d_1$ 的选取。
\end{purplenote}

\begin{corollary}\label{cor:subgroups-finitely-generated-abelian}
若阿贝尔群 $G$ 由 $n$ 个元素生成，则 $G$ 的任意子群都可由某个 $m\leq n$ 个元素生成。
\end{corollary}

\begin{purpleproof}
考虑由这 $n$ 个元素对应的形式生成元所构造的自由阿贝尔群 $F$，由泛性质得到满同态 $\pi:F\to G$。若 $H\leq G$，则 $\pi^{-1}(H)\leq F$。由定理~\ref{thm:subgroup-finite-rank-free-abelian}，非零的 $\pi^{-1}(H)$ 有一个至多含 $n$ 个元素的基；这些基元素在 $\pi$ 下的像生成 $H$。零子群的情形显然。
\end{purpleproof}

下面是一些本节的习题。

\begin{proposition}\label{prop:multiples-and-direct-sums}
设 $G$ 及各 $G_i$ 都是阿贝尔群，$m\in\mathbb Z$，则：
\begin{enumerate}[label=(\arabic*)]
  \item $mG=\{mu\mid u\in G\}$ 是 $G$ 的正规子群；
  \item 若 $G\cong\sum_{i\in I}G_i$，则
  \[
    mG\cong\sum_{i\in I}mG_i,
    \qquad
    G/mG\cong\sum_{i\in I}(G_i/mG_i).
  \]
\end{enumerate}
\end{proposition}

\begin{example}[线性独立与基]\label{ex:abelian-independence-and-bases}
设 $F$ 是阿贝尔群，$X\subseteq F$。称 $X$ \textbf{线性独立}，当且仅当对任意有限个两两不同的 $x_1,\ldots,x_k\in X$ 及任意整数 $n_1,\ldots,n_k$，都有
\[
  n_1x_1+\cdots+n_kx_k=0
  \quad\Longleftrightarrow\quad n_1=\cdots=n_k=0.
\]
\begin{enumerate}[label=(\arabic*)]
  \item $X$ 线性独立，当且仅当 $\langle X\rangle$ 中任意非零元素均可唯一写成
  \[
    n_1x_1+\cdots+n_kx_k,
  \]
  其中 $n_i\in\mathbb Z\setminus\{0\}$，各 $x_i\in X$ 两两不同（表示的唯一性不计各项的排列次序）。
  \item 若 $F$ 是秩为 $n$ 的自由阿贝尔群，则含 $n$ 个元素的线性独立子集 $X$ 不一定是 $F$ 的基。
  \item 若 $F$ 是自由阿贝尔群，则 $F$ 的线性独立子集 $X$ 不一定能扩充为 $F$ 的一个基。
  \item 若 $F$ 是自由阿贝尔群，则 $F$ 的生成集不一定包含 $F$ 的一个基；但若 $F$ 由 $n$ 个元素生成，则其秩 $m\leq n$。
\end{enumerate}
\end{example}

\begin{purplenote}

(2)、(3) 的反例为 $F=\mathbb Z$、$X=\{ 2 \}$。

(4) 的反例是 $\mathbb Z=\langle2,3\rangle$：生成集 $\{2,3\}$ 中不包含 $\mathbb Z$ 的基。

对于 (4) 最后的秩不等式，原稿还给出矩阵说明。若 $\{e_1,\ldots,e_m\}$ 是 $F$ 的基，$g_1,\ldots,g_n$ 生成 $F$，则可取整数矩阵 $A\in M_{m\times n}(\mathbb Z)$、$B\in M_{n\times m}(\mathbb Z)$，使得
\[
  \begin{pmatrix}e_1\\\vdots\\e_m\end{pmatrix}
  =A\begin{pmatrix}g_1\\\vdots\\g_n\end{pmatrix},
  \qquad
  \begin{pmatrix}g_1\\\vdots\\g_n\end{pmatrix}
  =B\begin{pmatrix}e_1\\\vdots\\e_m\end{pmatrix}.
\]
代入并用基的线性无关性，得到 $AB=I_m$。把整数矩阵视为有理数域上的矩阵，便有
\[
  m=\operatorname{rank}(AB)
  \leq\min\{\operatorname{rank}A,\operatorname{rank}B\}
  \leq\min\{m,n\},
\]
所以 $m\leq n$。这里需区分阿贝尔群的秩与矩阵的秩。
\end{purplenote}

% 来源：微信图片_20261114000846_37_26.jpg，2.1.13--2.1.19。
\begin{proposition}\label{prop:free-abelian-commutator-presentation}
设 $X=\{a_i\mid i\in I\}$ 是一个集合，则 $X$ 上的自由阿贝尔群同构于由生成元 $X$ 和关系
\[
  \{a_ia_ja_i^{-1}a_j^{-1}=e\mid i,j\in I\}
\]
决定的群。
\end{proposition}

% 来源：微信图片_20261114000850_38_26.jpg，2.1.13 的证明。
\begin{proof}
设 $F$ 是 $X$ 上的自由阿贝尔群，$L$ 是 $X$ 上的自由群，$N$ 是 $L$ 中包含所有 $a_ia_ja_i^{-1}a_j^{-1}$ 的最小正规子群。则 $L/N$ 正是由上述生成元和关系决定的群。因为生成元的像满足
\[
  a_ia_jN=a_ja_iN,
\]
所以 $L/N$ 是阿贝尔群。考虑集合映射
\[
  X\longrightarrow L/N,\qquad a_i\longmapsto a_iN.
\]
由自由阿贝尔群的泛性质，存在群同态 $\varphi:F\to L/N$。由于 $L/N$ 由这些陪集生成，$\varphi$ 是满同态。

另一方面，由自由群的泛性质，存在群同态 $\psi:L\to F$，将每个形式生成元 $a_i$ 映为 $F$ 中对应的基元素。其像包含 $F$ 的基，故 $\psi$ 是满同态。由于 $F$ 是阿贝尔群，
\[
  \psi(a_ia_ja_i^{-1}a_j^{-1})=0.
\]
因此 $N\subseteq\ker\psi$。于是 $\psi$ 诱导出满同态
\[
  \overline\psi:L/N\longrightarrow F.
\]
复合 $\overline\psi\circ\varphi$ 在 $X$ 上与恒等映射一致，由泛性质的唯一性，
\[
  \overline\psi\circ\varphi=\operatorname{id}_F.
\]
故 $\varphi$ 为单射；结合其满射性，得到 $F\cong L/N$。
\end{proof}

\begin{purplenote}
\end{purplenote}

\begin{proposition}\label{prop:free-abelian-also-free-group}
一个自由阿贝尔群是自由群，当且仅当它是无限循环群。
\end{proposition}

\begin{proposition}\label{prop:direct-sum-free-abelian}
一族自由阿贝尔群的直和仍是自由阿贝尔群；但其完全直和不一定是自由阿贝尔群。
\end{proposition}

\begin{example}\label{ex:finitely-generated-nonabelian-subgroup}
设 $G$ 是由实矩阵
\[
  a=\begin{pmatrix}2&0\\0&1\end{pmatrix},
  \qquad
  b=\begin{pmatrix}1&1\\0&1\end{pmatrix}
\]
在矩阵乘法下生成的群。令 $H$ 为 $G$ 中主对角元均为 $1$ 的矩阵组成的子集，则 $H$ 是一个不能由有限个元素生成的子群。
\end{example}

\begin{purplenote}
例~\ref{ex:finitely-generated-nonabelian-subgroup} 说明，推论~\ref{cor:subgroups-finitely-generated-abelian} 中的阿贝尔条件不能去掉。
\end{purplenote}

\begin{proposition}\label{prop:finitely-generated-torsion-free-abelian}
设 $G$ 是有限生成的阿贝尔群，其中除单位元外的元素都没有有限阶，则 $G$ 是自由阿贝尔群。
\end{proposition}

% 来源：微信图片_20261114000850_38_26.jpg，2.1.17 的证明。
% 补证（根据讨论）：用 2.1.10 和有限支撑证明线性无关集有限，再用模 2 商群估计其大小；
% 据此选取极大线性无关组，后半段沿用原稿的公倍数与单同态构造。
\begin{proof}
若 $G=0$，则按空基的约定，$G$ 是秩为 $0$ 的自由阿贝尔群。以下设 $G\ne0$，并写作 $G=\langle g_1,\ldots,g_n\rangle$。

先证明：$G$ 的任意线性无关子集 $Y$ 都是有限集，且 $|Y|\leq n$。$Y=\varnothing$ 时显然，以下设 $Y\ne\varnothing$。令 $H=\langle Y\rangle$，由推论~\ref{cor:subgroups-finitely-generated-abelian}，可写成
\[
  H=\langle h_1,\ldots,h_r\rangle,\qquad r\leq n.
\]
每个 $h_j$ 都是 $Y$ 中有限多个元素的整数线性组合。将这些表示中用到的元素合并，得到有限集 $Y_0\subseteq Y$，于是
\[
  H=\langle h_1,\ldots,h_r\rangle
  \subseteq\langle Y_0\rangle\subseteq H,
\]
故 $H=\langle Y_0\rangle$。若有 $y\in Y\setminus Y_0$，则存在整数 $a_z$（$z\in Y_0$），使得
\[
  y-\sum_{z\in Y_0}a_z z=0.
\]
这是 $Y$ 中元素的整数线性关系，且 $y$ 的系数为 $1$，与 $Y$ 线性无关矛盾。因此 $Y=Y_0$，从而 $Y$ 有限。

记 $s=|Y|$。由于 $Y$ 是 $H$ 的基，由定理~\ref{thm:free-abelian-characterizations} 及命题~\ref{prop:multiples-and-direct-sums}，
\[
  H\cong\mathbb Z^s,\qquad
  H/2H\cong(\mathbb Z/2\mathbb Z)^s,
  \qquad |H/2H|=2^s.
\]
另一方面，由 $H=\langle h_1,\ldots,h_r\rangle$，$H/2H$ 中每个元素均可写成
\[
  \sum_{j=1}^{r}\varepsilon_j(h_j+2H),
  \qquad \varepsilon_j\in\{0,1\},
\]
因为整数系数都可模 $2$ 化简。这样的组合至多有 $2^r$ 种，故
\[
  2^s=|H/2H|\leq2^r,
  \qquad |Y|=\operatorname{rank}H=s\leq r\leq n.
\]

由于 $G$ 无非零有限阶元素，任意非零元素组成的单元素集都线性无关。所有线性无关子集的大小又都不超过 $n$，故可选取一个元素个数最大的线性无关子集 $X$；它必然是极大的。记
\[
  X=\{x_1,\ldots,x_m\},\qquad
  F=\langle X\rangle\leq G,\qquad 1\leq m\leq n.
\]
由线性无关性，$F$ 是以 $X$ 为基、秩为 $m$ 的自由阿贝尔群。

由 $X$ 的极大性，对每个 $y\in G$，都存在整数 $k_0,k_1,\ldots,k_m$，其中 $k_0\ne0$，使得
\[
  k_0y=k_1x_1+\cdots+k_mx_m.
\]
这里若 $y\in X$，可取 $k_0=1$；若 $y\notin X$，由极大性，$X\cup\{y\}$ 线性相关，故有一组不全为零的整数系数使上式成立。若 $k_0=0$，则由 $X$ 线性无关又得 $k_1=\cdots=k_m=0$，矛盾。因此 $k_0\ne0$，从而 $|k_0|y\in F$。

特别地，对每个 $g_i$ 选取正整数 $n_i$，使得 $n_ig_i\in F$。令
\[
  T=\operatorname{lcm}(n_1,\ldots,n_n).
\]
由于 $g_1,\ldots,g_n$ 生成 $G$，对任意 $y\in G$ 都有 $Ty\in F$。于是定义
\[
  \psi:G\longrightarrow F,\qquad \psi(y)=Ty.
\]
这里 $T$ 是固定的，故 $\psi$ 是群同态。若 $\psi(y)=0$，则 $Ty=0$；由于 $G$ 无非零有限阶元素且 $T>0$，只能有 $y=0$。所以 $\psi$ 为单射，得到
\[
  G\cong\psi(G)\leq F.
\]
再由定理~\ref{thm:subgroup-finite-rank-free-abelian}，$\psi(G)$ 是自由阿贝尔群，故 $G$ 也是。
\end{proof}

\begin{purplenote}
注意极大线性无关组与生成集没有相互蕴含的关系。由于
\[
  Tx=Ty\quad\Longleftrightarrow\quad T(x-y)=0
  \quad\Longleftrightarrow\quad x=y,
\]
可知 $TG\cong G$，而 $TG\subseteq\langle X\rangle$。因此，原稿中“$G$ 其实也不大于 $F$”应理解为存在到 $F$ 的单同态。

整理注：本节开头按原稿要求基非空；涉及零群时，补用通常的约定：零群是以空集为基、秩为 $0$ 的自由阿贝尔群。命题~\ref{prop:free-abelian-also-free-group} 的陈述按原稿的非零情形理解。
\end{purplenote}

\begin{proposition}\label{prop:rationals-not-free-abelian}
\begin{enumerate}[label=(\arabic*)]
  \item 有理数加法群 $\mathbb Q$ 不是有限生成的；
  \item $\mathbb Q$ 不是自由阿贝尔群；
  \item 因而命题~\ref{prop:finitely-generated-torsion-free-abelian} 中的“有限生成”条件不能去掉。
\end{enumerate}
\end{proposition}

% 排版：将最后一个命题及其说明放在同一页，避免只把末条结论留到末页。
\noindent\begin{minipage}{\linewidth}
\setlength{\parskip}{0.65\baselineskip}
\begin{proposition}\label{prop:free-group-abelianization-rank}
设 $F$ 是 $X$ 上的自由群（不一定是阿贝尔群），$G$ 是 $Y$ 上的自由群。令
\[
  F'=\langle aba^{-1}b^{-1}\mid a,b\in F\rangle,
  \qquad
  G'=\langle aba^{-1}b^{-1}\mid a,b\in G\rangle.
\]
则：
\begin{enumerate}[label=(\arabic*)]
  \item $F'\trianglelefteq F$、$G'\trianglelefteq G$，且 $F/F'$、$G/G'$ 都是阿贝尔群；
  \item $F/F'$ 是秩为 $|X|$ 的自由阿贝尔群，其相应的基为 $\{xF'\mid x\in X\}$；类似地，$G/G'$ 是秩为 $|Y|$ 的自由阿贝尔群；
  \item $F\cong G$ 当且仅当 $|X|=|Y|$。
\end{enumerate}
\end{proposition}

\begin{purplenote}
(1) 的证明可参照命题~\ref{prop:free-abelian-commutator-presentation}。对于 (3)，若 $F\cong G$，则同构将换位子子群 $F'$ 映为 $G'$，从而 $F/F'\cong G/G'$；结合 (2) 和命题~\ref{prop:free-abelian-isomorphism-rank}，得到 $|X|=|Y|$。反向由自由群的泛性质得到。

由命题~\ref{prop:free-group-abelianization-rank} 可知，命题~\ref{prop:free-abelian-isomorphism-rank} 中相应的分类结论对给定的自由群也成立。
\end{purplenote}
\end{minipage}\par


% 新增：第 2.2 节独立维护，整理有限生成阿贝尔群的分解与全系不变量。
% !TeX root = ../main.tex
\subsection{有限生成阿贝尔群}\label{sec:finitely-generated-abelian-groups}

% 来源目录：../../NotePhotos/BaiduSyncdisk/algebra/2.2/（相对于本科目目录）。
% 共 8 张照片，文件名编号 39--46；编号不是原稿页码。
% 覆盖原稿 2.2.1--2.2.15；另页证明移至对应结论之后。
% 黑笔正文用普通环境，原紫笔说明用紫色；整理修正另在注释或“整理注”中注明。

% 来源：微信图片_20261115105904_39_26.jpg，节首与 2.2.1 的陈述。
\begin{purplenote}
本节的目的是寻找有限生成阿贝尔群的全系不变量，从而对其作分类。
\end{purplenote}

\begin{theorem}\label{thm:fg-abelian-invariant-factor-existence}
任意有限生成阿贝尔群 $G$ 都同构于有限个循环群的直和，其中有限循环群（若存在）的阶 $m_1,\ldots,m_t$ 满足
\[
  m_i>1,\qquad m_1\mid m_2\mid\cdots\mid m_t.
\]
换言之，存在非负整数 $s,t$，使得
\[
  G\cong\mathbb Z_{m_1}\oplus\cdots\oplus\mathbb Z_{m_t}
       \oplus\mathbb Z^s.
\]
\end{theorem}

\begin{purplenote}
该定理给出了有限生成阿贝尔群的一种分解。
\end{purplenote}

% 来源：微信图片_20261115105906_40_26.jpg，2.2.1 的证明。
% 整理：自由群的形式基记为 e_i，避免与 G 中的生成元混淆；单列核为零的情形。
\begin{proof}
设 $G=\langle g_1,\ldots,g_n\rangle$。取以形式元素 $e_1,\ldots,e_n$ 为基的自由阿贝尔群 $F$，由泛性质，存在满同态
\[
  \pi:F\longrightarrow G,\qquad e_i\longmapsto g_i.
\]
记 $K=\ker\pi$，由\hyperref[cor:first-isomorphism]{第一同构定理}，$F/K\cong G$。

若 $K=0$，则 $G\cong F\cong\mathbb Z^n$，结论成立。以下设 $K\ne0$。由定理~\ref{thm:subgroup-finite-rank-free-abelian}，存在 $F$ 的一个基 $\{x_1,\ldots,x_n\}$、正整数 $r\leq n$ 及正整数 $d_1,\ldots,d_r$，使得
\[
  d_1\mid d_2\mid\cdots\mid d_r,\qquad
  K=\langle d_1x_1\rangle\oplus\cdots\oplus\langle d_rx_r\rangle.
\]
令 $d_j=0$（$j>r$），则
\[
  F=\bigoplus_{j=1}^{n}\langle x_j\rangle,\qquad
  K=\bigoplus_{j=1}^{n}\langle d_jx_j\rangle.
\]
同构 $\langle x_j\rangle\cong\mathbb Z$ 将 $\langle d_jx_j\rangle$ 对应到 $d_j\mathbb Z$。由推论~\ref{cor:quotient-direct-products}，
\[
  G\cong F/K
  \cong\bigoplus_{j=1}^{n}
       \bigl(\langle x_j\rangle/\langle d_jx_j\rangle\bigr)
  \cong\bigoplus_{j=1}^{n}\mathbb Z/d_j\mathbb Z.
\]
其中
\[
  \mathbb Z/d_j\mathbb Z\cong
  \begin{cases}
    0,&d_j=1,\\
    \mathbb Z_{d_j},&d_j>1,\\
    \mathbb Z,&d_j=0.
  \end{cases}
\]
去掉零直和项，把 $d_j$ 中不为 $0$ 或 $1$ 的项依次记为 $m_1,\ldots,m_t$，把 $d_j=0$ 的项数记为 $s$，即得
\[
  G\cong\mathbb Z_{m_1}\oplus\cdots\oplus\mathbb Z_{m_t}
       \oplus\mathbb Z^s,
  \qquad m_1\mid\cdots\mid m_t.
\]
其中 $\mathbb Z^s$ 是秩为 $s$ 的自由阿贝尔群。
\end{proof}

% 来源：微信图片_20261115105904_39_26.jpg，2.2.2 的陈述。
\begin{lemma}\label{lem:cyclic-prime-power-decomposition}
若正整数 $m$ 的素因数分解为
\[
  m=p_1^{n_1}\cdots p_t^{n_t},
\]
其中 $p_1,\ldots,p_t$ 是互不相同的质数，$n_i>0$，则
\[
  \mathbb Z_m\cong
  \mathbb Z_{p_1^{n_1}}\oplus\cdots\oplus\mathbb Z_{p_t^{n_t}}.
\]
\end{lemma}

% 来源：微信图片_20261115105906_40_26.jpg，2.2.2 的证明。
% 整理：区分不同剩余类群中的元素，保留原稿用嵌入与裴蜀等式构造满同态的证明。
\begin{proof}
只需证明：若 $\gcd(r,m)=1$，则
\[
  \mathbb Z_{rm}\cong\mathbb Z_r\oplus\mathbb Z_m.
\]
在 $\mathbb Z_{rm}$ 中，元素 $\overline m$ 的阶为 $r$，故
$\langle\overline m\rangle\cong\mathbb Z_r$。定义
\[
  \psi_1:\mathbb Z_r\longrightarrow\mathbb Z_{rm},
  \quad \overline k\longmapsto\overline{mk},
  \qquad
  \psi_2:\mathbb Z_m\longrightarrow\mathbb Z_{rm},
  \quad \overline k\longmapsto\overline{rk}.
\]
两者均为单同态。由定理~\ref{thm:weak-direct-sum-universal-property}，存在群同态
\[
  \psi:\mathbb Z_r\oplus\mathbb Z_m\longrightarrow\mathbb Z_{rm},
  \qquad
  \psi(u,v)=\psi_1(u)+\psi_2(v).
\]
由 $\gcd(r,m)=1$，存在 $a,b\in\mathbb Z$ 使 $ra+mb=1$。因此，对任意整数 $k$，
\[
  \overline k=\overline{rak+mbk}
  =\psi_1(\overline{bk})+\psi_2(\overline{ak})
  =\psi(\overline{bk},\overline{ak}),
\]
其中两个输入分别取模 $r$、模 $m$ 的剩余类。故 $\psi$ 为满同态。又因两群均有 $rm$ 个元素，$\psi$ 为同构。对素因数分解逐次使用这一结论即可。
\end{proof}

% 来源：微信图片_20261115105904_39_26.jpg，2.2.3--2.2.5。
\begin{theorem}\label{thm:fg-abelian-elementary-divisor-existence}
任意有限生成阿贝尔群都同构于有限个循环群的直和，其中每个循环群的阶要么是无限的，要么是某个质数的幂。
\end{theorem}

\begin{purpleproof}
由定理~\ref{thm:fg-abelian-invariant-factor-existence}
和引理~\ref{lem:cyclic-prime-power-decomposition} 直接得到。
\end{purpleproof}

\begin{corollary}\label{cor:finite-abelian-subgroups-all-divisors}
若有限阿贝尔群 $G$ 有 $n$ 个元素，则对于每个正整数 $m\mid n$，$G$ 都有一个含 $m$ 个元素的子群。
\end{corollary}

\begin{purplenote}
证明提示：结合定理~\ref{thm:fg-abelian-elementary-divisor-existence}，利用对于质数 $p$ 及 $0\leq i\leq r$，
\[
  p^{r-i}\mathbb Z_{p^r}\cong\mathbb Z_{p^i}.
\]
\end{purplenote}

% 整理：原稿“下列任意一条是子群”实际对应 (1)--(4)；(5)--(7) 是同构性质。
\begin{lemma}\label{lem:abelian-characteristic-subgroups}
设 $G,H$ 及各 $G_i$ 都是阿贝尔群，$m\in\mathbb Z$，$p$ 为质数，则：
\begin{enumerate}[label=(\arabic*)]
  \item $mG=\{mu\mid u\in G\}$ 是 $G$ 的子群。
  \item $G[m]=\{u\in G\mid mu=0\}$ 是 $G$ 的子群。
  \item
  \[
    G(p)=\{u\in G\mid |u|=p^n\text{，对某个整数 }n\geq0\}
  \]
  是 $G$ 的子群。
  \item $G_t=\{u\in G\mid |u|<\infty\}$ 是 $G$ 的子群。
  \item 对 $n\geq1$，有
  \[
    \mathbb Z_{p^n}[p]\cong\mathbb Z_p;
    \qquad
    p^a\mathbb Z_{p^n}\cong\mathbb Z_{p^{n-a}}
    \quad(0\leq a<n).
  \]
  \item 若 $g:G\longrightarrow\sum_{i\in I}G_i$ 是同构，则它在 $mG$ 与 $G[m]$ 上的限制分别给出同构
  \[
    mG\cong\sum_{i\in I}mG_i,\qquad
    G[m]\cong\sum_{i\in I}G_i[m].
  \]
  \item 若 $f:G\longrightarrow H$ 是同构，则其在 $G_t$ 与 $G(p)$ 上的限制也分别为同构，即
  \[
    G_t\cong H_t,\qquad G(p)\cong H(p).
  \]
\end{enumerate}
\end{lemma}

% 来源：微信图片_20261115105904_39_26.jpg，第 (5) 条的紫笔旁注；
% 微信图片_20261115105909_41_26.jpg，承接 2.2.5 的紫笔说明。
\begin{purplenote}
第 (5) 条中，$\mathbb Z_{p^n}[p]$ 的同构类型与 $n\geq1$ 无关，始终为 $\mathbb Z_p$。

第 (1)--(4) 条中的阿贝尔条件不可删去：$S_3$ 可作为第 (1)--(3) 条的反例；第 (4) 条的反例见教材 Chapter~I, Exercise~3.5。

引理中的 $G_t$ 称为 $G$ 的\textbf{挠子群}。若 $G=G_t$，称 $G$ 为\textbf{挠群}；若 $G_t=0$，称 $G$ 为\textbf{无挠群}。
\end{purplenote}

% 来源：微信图片_20261115105909_41_26.jpg，2.2.6 的陈述。
\begin{theorem}\label{thm:fg-abelian-decomposition-uniqueness}
设 $G$ 是有限生成阿贝尔群，则：
\begin{enumerate}[label=(\arabic*)]
  \item 存在唯一的非负整数 $s$，使得将 $G$ 分解成循环群的直和时，任意一种分解中无限循环群的个数都是 $s$。
  \item 要么 $G$ 是自由阿贝尔群；要么存在唯一的一列不一定两两不同的正整数 $m_1,\ldots,m_t$，使得
  \[
    m_1>1,\qquad m_1\mid m_2\mid\cdots\mid m_t,
    \qquad
    G\cong\mathbb Z_{m_1}\oplus\cdots\oplus\mathbb Z_{m_t}\oplus F,
  \]
  其中 $F$ 为自由阿贝尔群。
  \item 要么 $G$ 是自由阿贝尔群；要么存在一列在换序意义下唯一的质数幂 $p_1^{s_1},\ldots,p_k^{s_k}$，使得
  \[
    G\cong
    \mathbb Z_{p_1^{s_1}}\oplus\cdots\oplus\mathbb Z_{p_k^{s_k}}\oplus F,
  \]
  其中 $p_1,\ldots,p_k$ 不一定不同，$s_1,\ldots,s_k$ 不一定为不同的正整数，$F$ 为自由阿贝尔群。
\end{enumerate}
\end{theorem}

% 来源：微信图片_20261115105911_42_26.jpg，2.2.6 的证明；原稿顺序为 (1)、(3)、(2)。
% 整理：比较素数幂指数时明确 a_k<c_k；不变因子分解中的素数指数允许为零。
\begin{proof}
(1) 由定理~\ref{thm:fg-abelian-invariant-factor-existence}，循环群的直和分解存在；略去零直和项后，有限生成性保证这样的分解只有有限项。将有限循环群的直和记为 $H$，无限循环群的直和记为 $F$，则
\[
  G\cong H\oplus F.
\]
这里 $H$ 为有限群，$F$ 为有限秩自由阿贝尔群。由引理~\ref{lem:abelian-characteristic-subgroups} 的第 (7) 条，
\[
  G_t\cong(H\oplus F)_t=H\oplus0.
\]
因此
\[
  G/G_t\cong(H\oplus F)/(H\oplus0)\cong F.
\]
故 $F$ 的同构类型只由 $G$ 决定，与所选分解无关；再由定理~\ref{thm:free-abelian-rank}，其秩 $s$ 唯一。

(3) 存在性由定理~\ref{thm:fg-abelian-elementary-divisor-existence} 得到。设有两种分解
\[
  G\cong\bigoplus_{i=1}^{u}\mathbb Z_{p_i^{s_i}}\oplus F
  \cong\bigoplus_{j=1}^{v}\mathbb Z_{q_j^{r_j}}\oplus F'.
\]
对每个固定质数 $p$，由引理~\ref{lem:abelian-characteristic-subgroups} 的第 (5)、(6) 条，
\[
  G[p]\cong
  \bigoplus_{\{i:p_i=p\}}\mathbb Z_p
  \cong\bigoplus_{\{j:q_j=p\}}\mathbb Z_p.
\]
两边元素个数相同，因此 $p$ 的幂次循环直和项的个数相同，记为 $m_p$。比较 $p$ 分量，可写成
\[
  G(p)\cong\bigoplus_{i=1}^{m_p}\mathbb Z_{p^{a_i}}
       \cong\bigoplus_{i=1}^{m_p}\mathbb Z_{p^{c_i}},
  \qquad
  1\leq a_1\leq\cdots\leq a_{m_p},\quad
  1\leq c_1\leq\cdots\leq c_{m_p}.
\]
若两列指数不同，令 $k$ 为第一个满足 $a_k\ne c_k$ 的指标。交换两种分解后，不妨设 $a_k<c_k$。于是
\[
  p^{a_k}G(p)\cong
  \bigoplus_{i=k+1}^{m_p}\mathbb Z_{p^{a_i-a_k}}.
\]
右侧至多有 $m_p-k$ 个非零循环直和项，因此
\[
  \bigl|(p^{a_k}G(p))[p]\bigr|\leq p^{m_p-k}.
\]
但从另一种分解看，对所有 $i\geq k$ 均有 $c_i\geq c_k>a_k$，从而
\[
  p^{a_k}G(p)\cong
  \bigoplus_{i=1}^{m_p}\mathbb Z_{p^{\max(c_i-a_k,0)}}
\]
至少有 $m_p-k+1$ 个非零循环直和项。因此
\[
  \bigl|(p^{a_k}G(p))[p]\bigr|\geq p^{m_p-k+1},
\]
矛盾。故 $a_i=c_i$ 对所有 $i$ 成立，即所有质数幂直和项连同重数在换序下唯一。

(2) 存在性由定理~\ref{thm:fg-abelian-invariant-factor-existence} 得到。设
\[
  G\cong\mathbb Z_{m_1}\oplus\cdots\oplus\mathbb Z_{m_t}\oplus F
   \cong\mathbb Z_{\ell_1}\oplus\cdots\oplus\mathbb Z_{\ell_d}\oplus F',
\]
其中
\[
  1<m_1\mid\cdots\mid m_t,\qquad
  1<\ell_1\mid\cdots\mid\ell_d.
\]
取所有出现的质数 $p_1,\ldots,p_r$，写出素因数分解
\[
  m_j=\prod_{i=1}^{r}p_i^{a_{ji}}\quad(1\leq j\leq t),
  \qquad
  \ell_j=\prod_{i=1}^{r}p_i^{c_{ji}}\quad(1\leq j\leq d),
\]
其中指数为非负整数，且对每个 $i$，
\[
  a_{1i}\leq\cdots\leq a_{ti},\qquad
  c_{1i}\leq\cdots\leq c_{di}.
\]
由 (3)，每个质数对应的非零指数及其重数均唯一。由于 $m_1>1$，可取一个质数 $p\mid m_1$；整除链使得 $p$ 出现在全部 $t$ 项中，而另一条链也必须有这些项，故 $t\leq d$。对称地 $d\leq t$，所以 $t=d$。对每个 $p_i$，将其唯一的非零指数列在左边补零到长度 $t$，便得到相同的两列指数。因此 $a_{ji}=c_{ji}$，从而 $m_j=\ell_j$，唯一性得证。
\end{proof}

\begin{purplenote}
定理~\ref{thm:fg-abelian-decomposition-uniqueness} 第 (2) 条中的 $m_1,\ldots,m_t$ 称为 $G$ 的\textbf{不变因子组}；第 (3) 条中的质数幂连同其重数称为 $G$ 的\textbf{初等因子组}。
\end{purplenote}

% 来源：微信图片_20261115105909_41_26.jpg，2.2.7--2.2.8。
\begin{corollary}\label{cor:fg-abelian-isomorphism-invariants}
设 $G,H$ 是有限生成阿贝尔群，则 $G\cong H$ 当且仅当 $G/G_t$ 与 $H/H_t$ 的秩相同，且 $G,H$ 的不变因子组相同（等价地，初等因子组相同）。
\end{corollary}

\begin{purplenote}
下面是一些习题结论。
\end{purplenote}

\begin{proposition}[$p$-群]\label{prop:torsion-abelian-primary-decomposition}
若一个群的每个元素的阶都是某个固定质数 $p$ 的幂，则称它为一个 \textbf{$p$-群}；其中单位元的阶 $1=p^0$ 也计入。

设 $G$ 是阿贝尔挠群，则：
\begin{enumerate}[label=(\arabic*)]
  \item $G(p)$ 是 $G$ 中唯一的最大 $p$-子群，即 $G$ 的任意 $p$-子群都包含在 $G(p)$ 中。
  \item
  \[
    G=\sum_{p}G(p),
  \]
  求和遍及使 $G(p)\ne0$ 的所有质数。
  \item 若 $H$ 是另一个阿贝尔挠群，则
  \[
    G\cong H
    \quad\Longleftrightarrow\quad
    G(p)\cong H(p)\text{ 对每个质数 }p\text{ 均成立}.
  \]
\end{enumerate}
\end{proposition}

\begin{purplenote}
第 (2) 条可对各有限生成子群运用定理~\ref{thm:fg-abelian-elementary-divisor-existence}；也可先考察各 $G(p)$ 与其余分量之和的交，再利用裴蜀定理。
\end{purplenote}

% 来源：微信图片_20261115105914_43_26.jpg，2.2.9--2.2.13。
\begin{proposition}\label{prop:cyclic-multiples-and-torsion}
设 $k,m$ 为正整数。若 $\gcd(k,m)=1$，则
\[
  k\mathbb Z_m=\mathbb Z_m,\qquad
  \mathbb Z_m[k]=\{0\}.
\]
若 $k\mid m$，写作 $m=kd$，则
\[
  k\mathbb Z_m\cong\mathbb Z_d,\qquad
  \mathbb Z_m[k]\cong\mathbb Z_k.
\]
\end{proposition}

\begin{theorem}\label{thm:maximal-order-element-direct-summand}
设 $G$ 是有限阿贝尔群，$x\in G$ 是一个阶最大的元素，则 $\langle x\rangle$ 是 $G$ 的一个直和因子，即存在子群 $H\leq G$，使得
\[
  G=\langle x\rangle\oplus H.
\]
\end{theorem}

\begin{purplenote}
证明过程值得细细品味。
\end{purplenote}

% 来源：微信图片_20261115105916_44_26.jpg，2.2.10 的证明及紫笔旁注。
% 整理：原稿的子群交写成空集，统一为零子群；区分商群补因子与其完整原像。
\begin{proof}
先证明一个辅助结论：若 $\langle x\rangle\ne G$，则存在 $G$ 中一个阶为质数的子群 $K$，满足
\[
  K\cap\langle x\rangle=\{0\}.
\]
令 $m=|x|$。由命题~\ref{prop:abelian-lcm-order} 和 $m$ 的最大性，$G$ 中每个元素的阶均整除 $m$。取 $y\notin\langle x\rangle$，设 $|y|=k$，则 $k\mid m$。考虑标准满同态
\[
  \pi:G\longrightarrow G/\langle x\rangle,
  \qquad d=|\pi(y)|.
\]
因为 $y\notin\langle x\rangle$，有 $d>1$，且 $d\mid k$。取质数 $p\mid d$，则
\[
  \overline y=\frac{d}{p}\pi(y)
\]
在商群中的阶为 $p$。又 $dy\in\langle x\rangle$，故存在 $a\in\mathbb Z$，使得 $dy=ax$。于是
\[
  0=\frac{k}{d}\,dy=\frac{ak}{d}\,x,
  \qquad m\mid\frac{ak}{d}.
\]
所以存在 $t\in\mathbb Z$ 使 $ak=mdt$。由于 $k\mid m$，
\[
  a=\frac{m}{k}\,dt,
\]
从而 $p\mid a$，可写 $a=pq$。令
\[
  v=\frac{d}{p}y-qx.
\]
则 $pv=dy-pqx=0$，且
\[
  \pi(v)=\frac{d}{p}\pi(y)=\overline y\ne0.
\]
因此 $v\notin\langle x\rangle$，$|v|=p$。取 $K=\langle v\rangle$，便有 $K\cap\langle x\rangle=\{0\}$，辅助结论得证。

现在对 $|G|$ 作归纳。$|G|=1$ 或 $G=\langle x\rangle$ 时结论显然。否则取上述 $K$，考察标准满同态
\[
  \rho:G\longrightarrow G/K.
\]
由于 $K\cap\langle x\rangle=\{0\}$，元素 $\overline x=\rho(x)$ 的阶仍为 $m$。商群中任意元素的阶整除其原像的阶，故不超过 $m$。因此 $\overline x$ 仍是 $G/K$ 中阶最大的元素。由归纳假设，
\[
  G/K=\langle\overline x\rangle\oplus\overline H
\]
对某个子群 $\overline H$ 成立。令 $H=\rho^{-1}(\overline H)$，则易知 $K\leq H$。

若 $u\in\langle x\rangle\cap H$，则
\[
  \rho(u)\in\langle\overline x\rangle\cap\overline H=\{0\}.
\]
所以 $u\in K\cap\langle x\rangle=\{0\}$。另一方面，任取 $g\in G$，可写
\[
  \rho(g)=a\overline x+\overline h,
  \qquad a\in\mathbb Z,\quad\overline h\in\overline H.
\]
选取 $h\in H$ 使 $\rho(h)=\overline h$，则 $g-ax-h\in K\leq H$，故 $g\in\langle x\rangle+H$。因此 $G=\langle x\rangle\oplus H$。
\end{proof}

\begin{purplenote}
先考虑阶为质数 $p$ 的子群，是因为这种子群与 $\langle x\rangle$ 的交只可能为零子群或它自身。

用好商群的性质，可以把 $\rho(ax)=0$ 转化为 $ax\in K$，再与元素的阶联系起来。
\end{purplenote}

\begin{proposition}\label{prop:lift-order-from-maximal-cyclic-quotient}
设 $G$ 是有限阿贝尔 $p$-群，$x\in G$ 有最大的阶。若
$\overline y\in G/\langle x\rangle$ 的阶为 $p^r$，则存在代表元 $y\in G$，使得 $|y|=p^r$。
\end{proposition}

\begin{purplenote}
若 $|x|=p^t$，则由阶的最大性，$p^tG=0$。取 $\overline y$ 的任一代表元 $y$，有 $r\leq t$，且
\[
  p^ry=ax
\]
对某个 $a\in\mathbb Z$ 成立。因此
\[
  0=p^ty=p^{t-r}(p^ry)=p^{t-r}ax,
\]
由 $|x|=p^t$ 得 $p^r\mid a$。写作 $a=p^rq$，取 $y'=y-qx$，便有 $p^ry'=0$；又 $y'$ 在商群中的像的阶为 $p^r$，而同态像的阶整除原像的阶，故 $|y'|=p^r$。
\end{purplenote}

\begin{proposition}\label{prop:finite-abelian-p-group-filtration}
设 $G$ 是有限阿贝尔 $p$-群，则：
\begin{enumerate}[label=(\arabic*)]
  \item 对所有 $n\geq0$，
  \[
    p^{n+1}G\cap G[p]\leq p^nG\cap G[p].
  \]
  \item 商群
  \[
    (p^nG\cap G[p])/(p^{n+1}G\cap G[p])
  \]
  是若干个 $\mathbb Z_p$ 的直和；记直和项的数量为 $k_n$。
  \item 将 $G$ 分解为阶为质数幂的循环群的有限直积，则 $k_n$ 恰为其中 $\mathbb Z_{p^{n+1}}$ 的个数。
\end{enumerate}
\end{proposition}

\begin{purplenote}
第 (2) 条可直接使用结构定理：先证明该商群中每个元素的阶至多为 $p$。
\end{purplenote}

\begin{proposition}\label{prop:fg-abelian-cancellation}
设 $G,H,K$ 都是有限生成阿贝尔群，则：
\begin{enumerate}[label=(\arabic*)]
  \item 若 $G\oplus G\cong H\oplus H$，则 $G\cong H$。
  \item 若 $G\oplus H\cong G\oplus K$，则 $H\cong K$。
  \item 若 $G_1$ 是秩为 $\aleph_0$ 的自由阿贝尔群，则
  \[
    G_1\oplus\mathbb Z\oplus\mathbb Z\cong G_1\oplus\mathbb Z,
    \qquad
    \mathbb Z\oplus\mathbb Z\not\cong\mathbb Z.
  \]
\end{enumerate}
这里第 (3) 条给出无限生成情形的对照。
\end{proposition}

% 来源：微信图片_20261115105918_45_26.jpg，2.2.14--2.2.15 的陈述。
\begin{theorem}\label{thm:finite-abelian-quotient-embeds}
设 $G$ 是有限阿贝尔群，则对任意子群 $H\leq G$，$G$ 都有一个子群同构于 $G/H$。
\end{theorem}

% 来源：微信图片_20261115105921_46_26.jpg，2.2.14 的证明与紫笔说明。
% 整理修正：|p^kG/p^{k+1}G| 的指数应计数 a_i>=k+1，而非原稿的 a_i>=k；
% 因此末步取 k=b_i-1。原稿“子群交为空集”等记号亦统一为零子群。
\begin{proof}
由于 $G$ 是有限阿贝尔群，由命题~\ref{prop:torsion-abelian-primary-decomposition}，
\[
  G=\bigoplus_{i=1}^{r}G(p_i),\qquad
  H=\bigoplus_{i=1}^{r}H(p_i).
\]
由 $H\leq G$，$H(p_i)=H\cap G(p_i)\leq G(p_i)$。再由推论~\ref{cor:quotient-direct-products}，
\[
  G/H\cong\bigoplus_{i=1}^{r}G(p_i)/H(p_i).
\]
因此只需讨论 $G$ 为有限阿贝尔 $p$-群的情形。

先证明一个辅助结论。设 $A,B$ 为有限群，$\varphi:A\to B$ 为满同态，$B_1\trianglelefteq B$，令 $A_1=\varphi^{-1}(B_1)$，则 $A_1\trianglelefteq A$，且
\[
  |B/B_1|\mid|A/A_1|,\qquad |B/B_1|\leq|A/A_1|.
\]
事实上，对 $a\in A$、$a_1\in A_1$，
\[
  \varphi(aa_1a^{-1})
  =\varphi(a)\varphi(a_1)\varphi(a)^{-1}\in B_1,
\]
故 $A_1\trianglelefteq A$。映射
\[
  \overline\varphi:A/A_1\longrightarrow B/B_1,\qquad
  aA_1\longmapsto\varphi(a)B_1
\]
良定义且为满同态。由\hyperref[cor:first-isomorphism]{第一同构定理}
和\hyperref[cor:lagrange]{Lagrange 定理}，
\[
  |B/B_1|=\frac{|A/A_1|}{|\ker\overline\varphi|},
\]
辅助结论得证。

回到原定理。由结构定理，不妨写成
\[
  G\cong\bigoplus_{i=1}^{t}\mathbb Z_{p^{a_i}},
  \qquad
  G/H\cong\bigoplus_{i=1}^{s}\mathbb Z_{p^{b_i}},
\]
其中
\[
  a_1\geq\cdots\geq a_t\geq1,\qquad
  b_1\geq\cdots\geq b_s\geq1.
\]
只需证明 $s\leq t$，且 $b_i\leq a_i$（$1\leq i\leq s$）。因为此时每个 $\mathbb Z_{p^{b_i}}$ 都能嵌入相应的 $\mathbb Z_{p^{a_i}}$，其直和便给出所需子群。

令 $\pi:G\to G/H$ 为标准满同态。对每个 $k\geq0$，$\pi$ 在 $p^kG$ 上的限制满射到 $p^k(G/H)$。取
\[
  A=p^kG,\quad B=p^k(G/H),\quad B_1=p^{k+1}(G/H),
  \quad A_1=(\pi|_A)^{-1}(B_1).
\]
由于 $p^{k+1}G\leq A_1$，由辅助结论得
\[
  \left|p^k(G/H)/p^{k+1}(G/H)\right|
  \leq |A/A_1|
  \leq \left|p^kG/p^{k+1}G\right|.
\]
另一方面，每个循环直和项的贡献为
\[
  p^k\mathbb Z_{p^a}/p^{k+1}\mathbb Z_{p^a}
  \cong
  \begin{cases}
    \mathbb Z_p,&a\geq k+1,\\
    0,&a\leq k.
  \end{cases}
\]
因此
\[
  \left|p^kG/p^{k+1}G\right|
  =p^{\#\{i:a_i\geq k+1\}},
\]
对 $G/H$ 也有同样的公式。由此得到
\[
  \#\{i:b_i\geq k+1\}
  \leq\#\{i:a_i\geq k+1\}
  \qquad(k\geq0).
\]
取 $k=0$，得 $s\leq t$。对 $1\leq i\leq s$，取 $k=b_i-1$；由于 $b_1\geq\cdots\geq b_i$，左侧至少为 $i$，故右侧至少为 $i$，从而 $a_i\geq b_i$，结论成立。
\end{proof}

\begin{purplenote}
遇到 $p$-群时，应重视向各个素分量化归的技巧。上面的辅助结论值得积累，在学习向量空间后也可用更简单的语言描述。

另一个思路是考察
\[
  \left|p^kG/p^{k+1}G\right|
  =\frac{|p^kG|}{|p^{k+1}G|},
\]
相当于利用乘以 $p^k$ 来过滤；可与命题~\ref{prop:finite-abelian-p-group-filtration} 的第 (1)、(3) 条对照。

整理注：原稿把本证明中的计数写为 $\#\{i:a_i\geq k\}$。这里统一为 $\#\{i:a_i\geq k+1\}$，相应地将最后的取值改为 $k=b_i-1$，以与商群 $p^kG/p^{k+1}G$ 的实际阶一致。
\end{purplenote}

\begin{theorem}\label{thm:fg-abelian-rank-quotient}
设 $G$ 是有限生成阿贝尔群，$G/G_t$ 的秩为 $n$。若 $H\leq G$，且 $H/H_t$ 的秩为 $m$，则
\[
  m\leq n,\qquad
  \operatorname{rank}\bigl((G/H)/(G/H)_t\bigr)=n-m.
\]
\end{theorem}

% 来源：微信图片_20261115105918_45_26.jpg，2.2.15 的紫笔引言与黑笔证明。
\begin{purplenote}
需注意：有限生成群的子群不一定有限生成，但有限生成阿贝尔群的子群一定有限生成，这是由推论~\ref{cor:subgroups-finitely-generated-abelian} 得到的。
\end{purplenote}

% 整理：原稿的 Z_n 按“秩为 n 的自由阿贝尔群”语境写为 Z^n；
% H=F∪H_t 按随后直和结论统一为 H=F+H_t；补出纯化在商群中的对应与 m=0 情形。
\begin{proof}
考察标准满同态
\[
  \pi:G\longrightarrow G/G_t.
\]
注意到 $H_t=H\cap G_t$。定义
\[
  \varphi:H/H_t\longrightarrow(H+G_t)/G_t,\qquad
  h+H_t\longmapsto h+G_t.
\]
若 $h-h'\in H_t$，则 $h-h'\in G_t$，故 $\varphi$ 良定义。它显然为满同态；若 $\varphi(h+H_t)=G_t$，则 $h\in H\cap G_t=H_t$，所以 $\varphi$ 为单射。因此
\[
  H/H_t\cong(H+G_t)/G_t=\pi(H)\leq G/G_t.
\]
由定理~\ref{thm:fg-abelian-decomposition-uniqueness} 第 (1) 条，
$G/G_t$ 是秩为 $n$ 的自由阿贝尔群。再由定理~\ref{thm:subgroup-finite-rank-free-abelian}，$\pi(H)$ 是秩为某个 $m\leq n$ 的自由阿贝尔群（零子群按秩为 $0$ 理解）。

还可进一步说明 $H$ 有限生成。取 $H/H_t$ 的一个基 $\{y_1,\ldots,y_m\}$，为每个 $y_i$ 选取标准投影$H \rightarrow H/H_t$下的原像 $x_i\in H$，令 $F=\langle x_1,\ldots,x_m\rangle$, 易得$F\cap H_t=\{0\}$且$H=F+H_t$。所以
\[
  H=F\oplus H_t.
\]
这里 $F$ 由 $m$ 个元素生成，而 $G_t$ 为有限群，故 $H_t\leq G_t$ 也有限。因此 $H$ 有限生成。

下面证明商群的秩公式。令
\[
  K=\{g\in G\mid kg\in H\text{ 对某个正整数 }k\}.
\]
这是标准满同态 $G\to G/H$ 下 $(G/H)_t$ 的原像，因此 $H\leq K\leq G$，并且映射
\[
  \psi:G/H\longrightarrow G/K,\qquad
  g+H\longmapsto g+K
\]
是满同态，其核为 $(G/H)_t$。由\hyperref[cor:first-isomorphism]{第一同构定理}，
\[
  (G/H)/(G/H)_t\cong G/K.
\]
又 $G_t\leq K$，由\hyperref[cor:third-isomorphism]{第三同构定理}，
\[
  (G/G_t)/(K/G_t)\cong G/K.
\]
记
\[
  \overline G=G/G_t,\qquad
  \overline H=(H+G_t)/G_t,\qquad
  \overline K=K/G_t.
\]
则 $\overline G$ 自由且秩为 $n$，$\overline H$ 的秩为 $m$。此外，
\[
  \overline K=
  \{\overline g\in\overline G\mid
      k\overline g\in\overline H\text{ 对某个正整数 }k\}.
\]
一方面，$kg\in H$ 蕴含 $k\overline g\in\overline H$。反过来，若 $k\overline g\in\overline H$，则 $kg=h+t$，其中 $h\in H$、$t\in G_t$；取正整数 $a$ 使 $at=0$，便得 $akg=ah\in H$，故 $g\in K$。

若 $m=0$，则 $\overline H=0$；由 $\overline G$ 无挠，$\overline K=0$，此时结论直接成立。若 $m>0$，由定理~\ref{thm:subgroup-finite-rank-free-abelian}，可取 $\overline G$ 的一个基 $\{u_1,\ldots,u_n\}$ 和正整数
$d_1\mid\cdots\mid d_m$，使得
\[
  \overline H=\langle d_1u_1,\ldots,d_mu_m\rangle.
\]
若 $\overline u=\sum_{i=1}^{n}c_iu_i\in\overline K$，则存在正整数 $k$ 使 $k\overline u\in\overline H$。比较基的系数，得
\[
  d_i\mid kc_i\quad(1\leq i\leq m),\qquad
  c_i=0\quad(i>m).
\]
故 $\overline K\leq\langle u_1,\ldots,u_m\rangle$。反过来，取 $D=\operatorname{lcm}(d_1,\ldots,d_m)$，则对任意
$\overline u\in\langle u_1,\ldots,u_m\rangle$，有 $D\overline u\in\overline H$，从而 $\overline u\in\overline K$。因此
\[
  \overline K=\langle u_1,\ldots,u_m\rangle,\qquad
  G/K\cong\overline G/\overline K
      \cong\langle u_{m+1},\ldots,u_n\rangle
      \cong\mathbb Z^{n-m}.
\]
结合前面的同构，即得所需秩公式。
\end{proof}

% 来源：微信图片_20261115105921_46_26.jpg，末尾紫笔说明。
% 整理修正：原稿末注写“2.2.5”，但 K 与 bar K 出现在 2.2.15 的证明中，故引用后者。
\begin{purplenote}
定理~\ref{thm:fg-abelian-rank-quotient} 证明中的 $K$ 和 $\overline K$，分别称为 $H$ 在 $G$ 中、$\overline H$ 在 $\overline G$ 中的\textbf{纯化}。

即使 $g$ 在 $G$ 中有无限阶，$g+H$ 在 $G/H$ 中仍可能有有限阶。而$H$在$G$中的纯化$K$使得只要 $g\notin K$，就有 $|g+K|=\infty$。这是因为若某个非零正整数倍 $kg$ 属于 $K$，则存在正整数 $a$ 使 $akg\in H$，从而 $g\in K$。这种对非零整数倍的封闭性体现了纯化的作用。
\end{purplenote}



% 新增：第 2.3 节独立维护，整理 Krull--Schmidt 定理及正规子群链条件。
% !TeX root = ../main.tex
\subsection{Krull--Schmidt 定理}\label{sec:krull-schmidt}

% 来源目录：../../NotePhotos/BaiduSyncdisk/algebra/2.3/（相对于本科目目录）。
% 共 7 张照片，文件名编号 47--53；编号不是原稿页码。
% 覆盖原稿 2.3.1--2.3.17；另页证明移至对应条目之后，末页补充结论不另编号。
% 本节恢复乘法记号。黑笔主体用普通环境，原紫笔说明用紫色；整理性说明另行注明。

% 来源：微信图片_20261116121258_47_26.jpg，节首与 2.3.1--2.3.5。
\begin{purplenote}
本节探索哪些群（不一定是阿贝尔群）可以分解为不可分解群的直积。本节用乘法记号表示群的二元运算。
\end{purplenote}

\begin{definition}[不可分解群]\label{def:indecomposable-group}
群 $G$ 称为\textbf{不可分解的}，当且仅当 $G\ne\{e\}$，且 $G$ 不能表示为两个非平凡子群的内直积。
\end{definition}

% 整理：原稿的 A_n (n≠4) 按非平凡情形注明 n≥3；
% 作为“不可分解但非简单”的例子时，p^a 的指数取 a≥2，S_n 取 n≥3。
\begin{purplenote}
等价地，$G$ 不可分解，当且仅当
\[
  G\ne\{e\},\qquad
  G\cong H\times K\ \Longrightarrow\
  H=\{e\}\text{ 或 }K=\{e\}.
\]
任意非平凡简单群都是不可分解的，例如 $A_n$（$n=3$ 或 $n\geq5$）。反过来，不可分解群不一定是简单群，例如
\[
  \mathbb Z,\qquad
  \mathbb Z_{p^a}\ (p\text{ 为质数},\ a\geq2),\qquad
  S_n\ (n\geq3).
\]
\end{purplenote}

\begin{definition}[升链条件与降链条件]\label{def:group-chain-conditions}
群 $G$ 称为满足在（正规）子群上的\textbf{升链条件}（ACC），当且仅当对任意由 $G$ 的（正规）子群构成的升链
\[
  G_1\leq G_2\leq\cdots,
\]
存在正整数 $n$，使得
\[
  G_i=G_n\qquad(i\geq n).
\]
相应地，$G$ 称为满足在（正规）子群上的\textbf{降链条件}（DCC），当且仅当对任意由 $G$ 的（正规）子群构成的降链
\[
  G_1\geq G_2\geq\cdots,
\]
存在正整数 $n$，使得 $G_i=G_n$ 对所有 $i\geq n$ 成立。
\end{definition}

\begin{example}\label{ex:group-chain-conditions}
任意有限群都满足上述两个链条件。整数加法群 $\mathbb Z$ 满足子群上的升链条件，但不满足降链条件；对任意固定质数 $p$，Prüfer 群 $\mathbb Z(p^\infty)$ 满足子群上的降链条件，但不满足升链条件。
\end{example}

\begin{theorem}\label{thm:indecomposable-decomposition-existence}
若非平凡群 $G$ 满足在正规子群上的升链条件或降链条件，则 $G$ 是有限个不可分解子群的内直积。
\end{theorem}

% 整理：平凡群用空内直积理解，避免与不可分解群必须非平凡的定义冲突。
\begin{purplenote}
该结论包括 $G$ 本身不可分解的情形。

证明思路是不断细分直积因子：若不能经过有限次分解就全部得到不可分解因子，则可继续选取尚能细分的因子，由此构造与升链条件或降链条件矛盾的链。该定理给出分解存在的充分条件；后面的引理和推论主要用于研究分解在什么条件下唯一。
\end{purplenote}

\begin{definition}[正规自同态]\label{def:normal-endomorphism}
设 $G$ 是群。自同态 $f:G\to G$ 称为\textbf{正规自同态}，当且仅当
\[
  f(aba^{-1})=af(b)a^{-1}
  \qquad(a,b\in G).
\]
\end{definition}

% 来源：微信图片_20261116121302_48_26.jpg，2.3.6--2.3.10。
% 第 2.3.6 条中“正规”是原稿在自同态一词上方补写的条件。
\begin{lemma}\label{lem:normal-endomorphism-chain-automorphism}
设 $f$ 是群 $G$ 的正规自同态。
\begin{enumerate}[label=(\arabic*)]
  \item 若 $G$ 满足正规子群上的升链条件，则 $f$ 是自同构，当且仅当 $f$ 是满同态。
  \item 若 $G$ 满足正规子群上的降链条件，则 $f$ 是自同构，当且仅当 $f$ 是单同态。
\end{enumerate}
\end{lemma}

% 来源：微信图片_20261116121306_49_26.jpg，2.3.6 的证明。
\begin{proof}
(1) 只需证明满射蕴含单射。考察正规子群的升链
\[
  \ker f\leq\ker f^2\leq\cdots.
\]
由升链条件，存在正整数 $m$，使得
\[
  \ker f^n=\ker f^m\qquad(n\geq m).
\]
若 $f(a)=e$，由 $f$ 满射，可取 $c\in G$，使 $a=f^m(c)$。于是
\[
  f^{m+1}(c)=f(a)=e,
\]
故 $c\in\ker f^{m+1}=\ker f^m$，从而 $a=f^m(c)=e$。因此 $f$ 单射，进而为自同构。

(2) 只需证明单射蕴含满射。由于 $f$ 正规，其各次复合也正规，故 $\operatorname{im}f^n\trianglelefteq G$。考察降链
\[
  \operatorname{im}f\geq\operatorname{im}f^2\geq\cdots.
\]
由降链条件，存在正整数 $m$，使得
\[
  \operatorname{im}f^n=\operatorname{im}f^m\qquad(n\geq m).
\]
任取 $a\in G$，由 $\operatorname{im}f^m=\operatorname{im}f^{m+1}$，可取 $b\in G$，使
\[
  f^m(a)=f^{m+1}(b)=f^m(f(b)).
\]
由于 $f$ 单射，$f^m$ 也单射，故 $a=f(b)$。因此 $f$ 满射。
\end{proof}

\begin{purplenote}
降链情形用到了 $f$ 的正规性，以保证各 $\operatorname{im}f^n$ 都是 $G$ 的正规子群。
\end{purplenote}

\begin{lemma}[Fitting]\label{lem:fitting-normal-endomorphism}
设群 $G$ 同时满足在正规子群上的升链条件和降链条件，且 $f$ 为 $G$ 的正规自同态。则存在正整数 $n$，使得
\[
  G=\ker f^n\times\operatorname{im}f^n,
\]
这里是子群的内直积。
\end{lemma}

% 原稿只有证明提示，保留其范围；子群交中的空集记号统一为单位子群。
\begin{purplenote}
证明提示：结合链条件和引理~\ref{lem:normal-endomorphism-chain-automorphism}，取足够大的 $n$，使 $\ker f^n$ 与 $\operatorname{im}f^n$ 都不再变化；再证明
\[
  \ker f^n\cap\operatorname{im}f^n=\{e\},
  \qquad
  G=\langle\ker f^n\cup\operatorname{im}f^n\rangle.
\]
\end{purplenote}

\begin{definition}[幂零自同态]\label{def:nilpotent-endomorphism}
设 $f$ 是群 $G$ 的自同态。若存在正整数 $n$，使得
\[
  f^n(g)=e\qquad(g\in G),
\]
则称 $f$ 为\textbf{幂零的}。
\end{definition}

\begin{corollary}\label{cor:indecomposable-normal-endomorphism-dichotomy}
设 $G$ 是不可分解群，且同时满足正规子群上的升链条件和降链条件。若 $f$ 是 $G$ 的正规自同态，则 $f$ 要么幂零，要么是自同构。
\end{corollary}

\begin{purpleproof}
直接应用引理~\ref{lem:fitting-normal-endomorphism}。因为 $G$ 不可分解，分解中的一个因子必为单位子群：若 $\operatorname{im}f^n=\{e\}$，则 $f$ 幂零；若 $\ker f^n=\{e\}$，则 $f$ 单射，再由引理~\ref{lem:normal-endomorphism-chain-automorphism} 得 $f$ 为自同构。
\end{purpleproof}

\begin{purplenote}
若 $f,g$ 是 $G\to G$ 的函数，定义它们的“和”为逐点乘积：
\[
  (f+g)(a)=f(a)g(a)\qquad(a\in G).
\]
在此运算下，所有 $G\to G$ 的函数组成一个群，单位元为常值映射
\[
  0_G:G\longrightarrow G,\qquad 0_G(a)=e.
\]
但即使 $f,g$ 都是自同态，$f+g$ 也未必是自同态。

本节映射的乘积 $fg$ 表示复合 $f\circ g$，$f^n$ 表示 $n$ 次复合；它们与这里的逐点“和”应加以区分。
\end{purplenote}

\begin{corollary}\label{cor:sum-nilpotent-normal-endomorphisms}
设 $G\ne\{e\}$ 是不可分解群，且同时满足正规子群上的升链条件和降链条件。设 $f_1,\ldots,f_n$ 是 $G$ 的正规幂零自同态，并且对任意
\[
  1\leq i_1<\cdots<i_r\leq n,
\]
子和 $f_{i_1}+\cdots+f_{i_r}$ 都是自同态。则 $f_1+\cdots+f_n$ 是幂零的。
\end{corollary}

% 来源：微信图片_20261116121306_49_26.jpg，2.3.10 的证明。
% 整理：展开式按复合字说明，避免在非交换的逐点运算中无说明地使用交换环的二项式公式。
\begin{proof}
由正规性的定义，各子和只要是自同态，就仍为正规自同态。因此对 $n$ 归纳，只需证明 $n=2$ 的情形。

由推论~\ref{cor:indecomposable-normal-endomorphism-dichotomy}，$f_1+f_2$ 要么幂零，要么为自同构。假设后者成立，令
\[
  g=(f_1+f_2)^{-1},\qquad
  g_1=f_1g,\quad g_2=f_2g.
\]
正规自同构的逆仍正规，正规自同态的复合也正规，故 $g_1,g_2$ 均为正规自同态，且
\[
  g_1+g_2=(f_1+f_2)g=1_G.
\]
若某个 $g_i$ 为自同构，则 $f_i=g_ig^{-1}$ 也是自同构，这与 $G$ 非平凡且 $f_i$ 幂零矛盾。再由推论~\ref{cor:indecomposable-normal-endomorphism-dichotomy}，$g_1,g_2$ 都幂零。取正整数 $m$，使得
\[
  g_1^m=g_2^m=0_G.
\]

先说明 $g_1g_2=g_2g_1$。由 $1_G=g_1+g_2$，对任意 $x\in G$，
\[
  x^{-1}=g_1(x^{-1})g_2(x^{-1})
         =g_1(x)^{-1}g_2(x)^{-1},
\]
取逆得 $x=g_2(x)g_1(x)$，因此也有 $1_G=g_2+g_1$。利用复合的分配关系，
\[
  g_1(g_1+g_2)=g_1^2+g_1g_2=g_1,
  \qquad
  (g_1+g_2)g_1=g_1^2+g_2g_1=g_1.
\]
在逐点运算的群中消去 $g_1^2$，得到 $g_1g_2=g_2g_1$。

现在考察
\[
  1_G=(g_1+g_2)^{2m}.
\]
按分配律展开，右侧是由 $g_1,g_2$ 组成的长度为 $2m$ 的复合字之和。每个字中至少有一种映射出现 $m$ 次；由于 $g_1,g_2$ 在复合下交换，可以把这些映射移到一起，故每个字都等于 $0_G$。于是 $1_G=0_G$，与 $G\ne\{e\}$ 矛盾。因此 $f_1+f_2$ 必为幂零自同态，归纳即得结论。
\end{proof}

% 来源：微信图片_20261116121312_50_26.jpg，2.3.11 的陈述。
\begin{theorem}[Krull--Schmidt]\label{thm:krull-schmidt}
设群 $G$ 同时满足在正规子群上的升链条件和降链条件。若有两种内直积分解
\[
  G=G_1\times\cdots\times G_s
   =H_1\times\cdots\times H_t,
\]
其中各 $G_i,H_j$ 都不可分解，则 $s=t$。对 $H_j$ 重新排列指标后，有
\[
  G_i\cong H_i\qquad(1\leq i\leq s),
\]
而且对每个 $0\leq r\leq s$，都有内直积分解
\[
  G=G_1\times\cdots\times G_r
       \times H_{r+1}\times\cdots\times H_s.
\]
\end{theorem}

\begin{purplenote}
Krull--Schmidt 定理说明了定理~\ref{thm:indecomposable-decomposition-existence} 中分解的唯一性。这种唯一性既指对应因子同构，也包含上面逐个替换因子的内直积分解。
\end{purplenote}

% 来源：微信图片_20261116121315_51_26.jpg，2.3.11 的完整证明。
% 整理：投影统一记为 phi_i、psi_j；明确各子和的同态性及换因子映射 theta 的正规性。
% 本证明使用本节后列的 2.3.14--2.3.16，其位置和编号按原稿保留。
\begin{proof}
平凡群的情形按空直积理解，以下设 $G\ne\{e\}$。对
$0\leq r\leq\min(s,t)$ 作归纳。记 $P(r)$ 为如下命题：对 $H_1,\ldots,H_t$ 适当重排指标后，
\[
  G_i\cong H_i\quad(1\leq i\leq r),
  \qquad
  G=G_1\times\cdots\times G_r\times H_{r+1}\times\cdots\times H_t.
\]
$P(0)$ 由已知分解成立。

设 $P(r-1)$ 已成立，因而可将 $H_j$ 排列为
\[
  G=G_1\times\cdots\times G_{r-1}\times H_r\times\cdots\times H_t.
\]
对于原来的 $G_i$ 分解，令 $\pi_i:G\to G_i$ 为标准投影，$\lambda_i:G_i\to G$ 为嵌入，记
\[
  \varphi_i=\lambda_i\pi_i\qquad(1\leq i\leq s).
\]
对于上面的混合分解，记对应的投影后再嵌入的映射为 $\psi_1,\ldots,\psi_t$。由命题~\ref{prop:direct-product-normal-projections}，这些映射都是 $G$ 的正规自同态，并满足
\[
  \begin{gathered}
    \varphi_1+\cdots+\varphi_s=1_G,\qquad
    \varphi_i^2=\varphi_i,\qquad
    \varphi_i\varphi_j=0_G\ (i\ne j),\\
    \psi_1+\cdots+\psi_t=1_G,\qquad
    \psi_i^2=\psi_i,\qquad
    \psi_i\psi_j=0_G\ (i\ne j),
  \end{gathered}
\]
以及
\[
  \operatorname{im}\varphi_i=G_i,\qquad
  \operatorname{im}\psi_i=
  \begin{cases}
    G_i,&i<r,\\
    H_i,&i\geq r.
  \end{cases}
\]
特别地，$\varphi_r\psi_i=0_G$（$i<r$），而
\[
  \varphi_r
  =\varphi_r(\psi_1+\cdots+\psi_t)
  =\varphi_r\psi_r+\cdots+\varphi_r\psi_t.
\]
将此式限制在 $G_r$ 上，得到
\[
  1_{G_r}
  =(\varphi_r\psi_r)|_{G_r}+\cdots+(\varphi_r\psi_t)|_{G_r}.
\]
各项都是 $G_r$ 的正规自同态。
% 根据 2026-11-18 的讨论补充：原稿此处直接应用 2.3.10；以下展开任意子和的同态性。
为应用推论~\ref{cor:sum-nilpotent-normal-endomorphisms}，还需核实上述各项的任意子和都是 $G_r$ 的自同态。任取非空指标集 $J\subseteq\{r,\ldots,t\}$，记
\[
  P_J=\sum_{j\in J}\psi_j,
\]
其中各和按指标递增的顺序理解。由于 $\psi_j$ 来自同一个混合内直积分解，$P_J$ 恰好保留 $J$ 中各因子的分量，并把其余分量变成单位元，即投影到所选因子的直积后再嵌入 $G$。由命题~\ref{prop:direct-product-normal-projections}，$P_J$ 是 $G$ 的自同态。

又因 $\varphi_r$ 是群同态，对任意 $x\in G_r$，有
\[
  \begin{aligned}
    \left(\sum_{j\in J}(\varphi_r\psi_j)|_{G_r}\right)(x)
      &=\prod_{j\in J}\varphi_r(\psi_j(x))\\
      &=\varphi_r\!\left(\prod_{j\in J}\psi_j(x)\right)
       =(\varphi_r\circ P_J)(x).
  \end{aligned}
\]
因此该子和等于 $(\varphi_r\circ P_J)|_{G_r}$。复合映射 $\varphi_r\circ P_J$ 是同态，限制到子群 $G_r$ 后仍是同态；且其像包含于 $\operatorname{im}\varphi_r=G_r$，故该子和确实是 $G_r$ 的自同态。

由定理~\ref{thm:normal-chains-in-direct-factors}，$G_r$ 及各 $H_j$ 都继承所需的两个正规子群链条件。

若右侧各项都幂零，则推论~\ref{cor:sum-nilpotent-normal-endomorphisms} 将推出 $1_{G_r}$ 幂零，矛盾。因此存在 $r\leq j\leq t$，使得
\[
  \alpha=(\varphi_r\psi_j)|_{G_r}
\]
不是幂零的；由推论~\ref{cor:indecomposable-normal-endomorphism-dichotomy}，$\alpha$ 是自同构。

再考察 $H_j$ 上的正规自同态
\[
  \beta=(\psi_j\varphi_r)|_{H_j}.
\]
对所有 $n\geq1$，有
\[
  (\varphi_r\psi_j)^{n+1}|_{G_r}
  =\varphi_r(\psi_j\varphi_r)^n\psi_j|_{G_r}.
\]
若 $\beta$ 幂零，则此式将推出 $\alpha$ 幂零，矛盾。因此 $\beta$ 也非幂零，进而为自同构。由 $\alpha,\beta$ 都是自同构，可知
\[
  \psi_j|_{G_r}:G_r\longrightarrow H_j,\qquad
  \varphi_r|_{H_j}:H_j\longrightarrow G_r
\]
既单射又满射，所以 $G_r\cong H_j$。交换 $H_r$ 与 $H_j$ 后，可设 $j=r$，于是 $G_r\cong H_r$。

还需证明替换后的因子确实给出 $G$ 的内直积分解。令
\[
  B=G_1\cdots G_{r-1}H_{r+1}\cdots H_t.
\]
$\psi_r$ 在 $B$ 上为平凡映射，而 $\psi_r|_{G_r}$ 是单射，故
\[
  G_r\cap B=\{e\}.
\]
这里 $G_r$ 与 $B$ 都是 $G$ 的正规子群，因而相互交换。因此可组成内直积子群
\[
  G^*=G_1\times\cdots\times G_{r-1}\times G_r
          \times H_{r+1}\times\cdots\times H_t\leq G.
\]
利用归纳假设下的混合分解，定义
\[
  \begin{aligned}
    \theta:G&\longrightarrow G,\\
    g_1\cdots g_{r-1}h_rh_{r+1}\cdots h_t
      &\longmapsto
    g_1\cdots g_{r-1}\varphi_r(h_r)h_{r+1}\cdots h_t.
  \end{aligned}
\]
因为 $\varphi_r|_{H_r}:H_r\to G_r$ 是同构，其余因子上的映射是恒等映射，且目标各因子组成内直积，故 $\theta$ 是单同态，$\operatorname{im}\theta=G^*$。又
\[
  \theta=\psi_1+\cdots+\psi_{r-1}
          +\varphi_r\psi_r+\psi_{r+1}+\cdots+\psi_t.
\]
% 根据 2026-11-18 的讨论修订：theta 已知是自同态，此处只需验证正规性，无须再要求任意子和为同态。
上式各项均为正规自同态。结合前面已知 $\theta$ 是自同态，可知 $\theta$ 为正规自同态。由引理~\ref{lem:normal-endomorphism-chain-automorphism} 的降链情形，$\theta$ 是自同构。因此
\[
  G=\operatorname{im}\theta=G^*,
\]
完成 $P(r)$ 的归纳步骤。

最后取 $r=\min(s,t)$。若 $r=s<t$，则
\[
  G=G_1\times\cdots\times G_s
   =G_1\times\cdots\times G_s\times H_{s+1}\times\cdots\times H_t,
\]
迫使剩余因子为平凡群，与其不可分解矛盾。若 $r=t<s$，同样由
$G=G_1\times\cdots\times G_t$ 与原来的 $G_i$ 分解矛盾。因此 $s=t$，定理得证。
\end{proof}

% 来源：微信图片_20261116121312_50_26.jpg，2.3.12--2.3.15。
下面列出本节的一些重要结论。

\begin{proposition}\label{prop:indecomposable-group-examples}
\begin{enumerate}[label=(\arabic*)]
  \item 群 $G$ 不可分解，当且仅当
  \[
    G\ne\{e\},\qquad
    G\cong H\times K\ \Longrightarrow\
    H=\{e\}\text{ 或 }K=\{e\}.
  \]
  \item 对所有 $n\geq2$，对称群 $S_n$ 不可分解。
  \item 整数加法群 $\mathbb Z$、素数阶循环群 $\mathbb Z_p$（$p$ 为质数）及有理数加法群 $\mathbb Q$ 都不可分解。
  \item 不可分解群在非平凡同态下的像不一定不可分解。
\end{enumerate}
\end{proposition}

% 整理：为使原稿中的 Z -> Z_mn 确实给出可分解的像，明确 m,n>1。
\begin{purplenote}
第 (4) 条的反例：取互素的整数 $m,n>1$，标准满同态
\[
  \mathbb Z\longrightarrow\mathbb Z_{mn}
\]
的定义域不可分解，而由引理~\ref{lem:cyclic-prime-power-decomposition} 证明中的互素分解，
\[
  \mathbb Z_{mn}\cong\mathbb Z_m\times\mathbb Z_n
\]
是两个非平凡群的直积。
\end{purplenote}

\begin{proposition}\label{prop:abelian-subgroup-chain-conditions}
\begin{enumerate}[label=(\arabic*)]
  \item $\mathbb Z$ 满足子群上的升链条件，但不满足降链条件。
  \item 任意有限生成阿贝尔群都满足子群上的升链条件。
  \item 对任意固定质数 $p$，$\mathbb Z(p^\infty)$ 满足子群上的降链条件，但不满足升链条件。
\end{enumerate}
\end{proposition}

\begin{theorem}\label{thm:normal-chains-in-direct-factors}
设 $G=H\times K$ 是子群的内直积，则：
\begin{enumerate}[label=(\arabic*)]
  \item 若 $N\trianglelefteq H$，则 $N\trianglelefteq G$。
  \item 若 $G$ 满足正规子群上的升链条件（或降链条件），则 $H,K$ 也分别满足相应的正规子群链条件。
\end{enumerate}
\end{theorem}

\begin{proposition}\label{prop:normal-endomorphism-properties}
设 $f,g$ 是群 $G$ 的正规自同态，则：
\begin{enumerate}[label=(\arabic*)]
  \item $f\circ g$ 是正规自同态。
  \item 若 $H\trianglelefteq G$，则 $f(H)\trianglelefteq G$。
  \item 若 $f+g$ 是自同态，则 $f+g$ 是正规自同态。
\end{enumerate}
\end{proposition}

% 来源：微信图片_20261116121318_52_26.jpg，2.3.16--2.3.17。
\begin{proposition}\label{prop:direct-product-normal-projections}
设 $G=G_1\times\cdots\times G_n$ 是内直积。令 $\lambda_i:G_i\to G$ 为嵌入，$\pi_i:G\to G_i$ 为标准投影，记
\[
  \varphi_i=\lambda_i\circ\pi_i.
\]
则对任意互异的指标 $i_1,\ldots,i_k$（$1\leq k\leq n$），
\[
  \varphi_{i_1}+\cdots+\varphi_{i_k}
\]
都是 $G$ 的正规自同态。
\end{proposition}

\begin{proposition}\label{prop:krull-schmidt-cancellation}
\begin{enumerate}[label=(\arabic*)]
  \item 若 $G\times G\cong H\times H$，且 $G$ 同时满足正规子群上的 ACC 与 DCC，则 $G\cong H$。
  \item 若 $G\cong H\times K$、$G\cong H\times J$，且 $G$ 同时满足正规子群上的 ACC 与 DCC，则 $K\cong J$。
\end{enumerate}
\end{proposition}

% 来源：微信图片_20261116121322_53_26.jpg，无编号的关键补充结论与证明。
% 原稿详细展开 ACC 情形；DCC 情形用同样的比较论证说明，不另增加条目编号。
\par\textbf{补充结论（原稿另页）。}
下面的结论可用于命题~\ref{prop:krull-schmidt-cancellation}
以及命题~\ref{prop:abelian-subgroup-chain-conditions} 第 (2) 条：若 $K,H$ 分别满足正规子群上的 ACC（或分别满足 DCC），则 $G=K\times H$ 也满足相应的正规子群链条件。

\begin{proof}
先讨论 ACC。设
\[
  M_1\leq M_2\leq\cdots
\]
是 $G$ 的一条正规子群升链。取标准投影
\[
  \pi:G\longrightarrow H,\qquad \pi(k,h)=h.
\]
则
\[
  \pi(M_1)\leq\pi(M_2)\leq\cdots,\qquad
  M_1\cap K\leq M_2\cap K\leq\cdots
\]
分别是 $H$ 和 $K$ 的正规子群升链。因此存在正整数 $m$，使对所有 $n\geq m$，
\[
  \pi(M_n)=\pi(M_m),\qquad M_n\cap K=M_m\cap K.
\]
任取 $y\in M_n$。由投影像相同，可选 $x\in M_m$，使 $\pi(x)=\pi(y)$。由于 $M_m\leq M_n$，有
\[
  xy^{-1}\in M_n\cap K=M_m\cap K.
\]
从而 $y=(xy^{-1})^{-1}x\in M_m$。故 $M_n=M_m$，升链稳定。

DCC 情形同理。若 $M_1\geq M_2\geq\cdots$，则取 $m$ 使投影像和与 $K$ 的交在 $n\geq m$ 时均稳定。对 $y\in M_m$，可取 $x\in M_n$ 使 $\pi(x)=\pi(y)$；这时
\[
  xy^{-1}\in M_m\cap K=M_n\cap K,
\]
故 $y\in M_n$。因此 $M_m=M_n$，降链也稳定。
\end{proof}

\begin{purplenote}
这个证明的要点是：在一条包含链中，结合子群的一个投影像与它同投影核的交，就能控制稳定性(但两个都投影就不行)。原稿所说的“一个投影和一个截面”，这里分别对应 $\pi(M_i)$ 与 $M_i\cap K$。
\end{purplenote}



% 新增：第 2.4 节独立维护，整理群作用、轨道与稳定子及陪集作用的应用。
% !TeX root = ../main.tex
\subsection{群作用}\label{sec:group-actions}

% 来源目录：../../NotePhotos/BaiduSyncdisk/algebra/2.4/（相对于本科目目录）。
% 共 4 张照片，完整文件名分别记在对应正文前；文件名编号 55--58 不是原稿页码。
% 覆盖原稿 2.4.1--2.4.15，保留黑笔主体、紫笔说明与原稿的证明／提示范围。

% 来源：微信图片_20261118134403_55_26.jpg，节首、2.4.1--2.4.3 及紫笔说明。
\begin{purplenote}
本节研究的技术将用于后续对非阿贝尔有限群结构的研究。
\end{purplenote}

\begin{definition}[群作用]\label{def:group-action}
设 $G$ 是群，$S$ 是集合。$G$ 在 $S$ 上的一个\textbf{作用}是一个映射
\[
  G\times S\longrightarrow S,\qquad (g,x)\longmapsto gx,
\]
满足对任意 $x\in S$ 及 $g_1,g_2\in G$，都有
\[
  ex=x,\qquad (g_1g_2)x=g_1(g_2x).
\]
给定这样的映射后，称 $G$ 作用于集合 $S$。
\end{definition}

\begin{purplenote}
本节采用左作用，应注意作用与平移的方向；对于一般子群的左陪集，不能直接用右平移定义作用。

对给定的群 $G$ 和集合 $S$，可以有许多不同的作用映射，因此使用 $gx$ 等记号时需要注意上下文。不过，后面研究结构时，许多结论并不依赖作用映射的具体形式。
\end{purplenote}

\begin{example}\label{ex:group-actions-permutations-translations-conjugation}
\begin{enumerate}[label=(\arabic*)]
  \item 对称群 $S_n$ 在集合 $I_n=\{1,\ldots,n\}$ 上有自然作用
  \[
    (\sigma,x)\longmapsto\sigma(x).
  \]

  \item 设 $H\leq G$。由 $G$ 中的乘法定义 $H$ 在集合 $G$ 上的作用
  \[
    (h,x)\longmapsto hx.
  \]
  对固定的 $h\in H$，映射 $x\mapsto hx$ 称为由 $h$ 给出的\textbf{左平移}（left translation）。若 $K\leq G$，$S$ 是 $K$ 在 $G$ 中的所有左陪集组成的集合，则 $H$ 还可在 $S$ 上以左平移作用：
  \[
    (h,xK)\longmapsto hxK.
  \]

  \item 设 $H\leq G$。可以定义 $H$ 在集合 $G$ 上的\textbf{共轭作用}
  \[
    (h,x)\longmapsto hxh^{-1}.
  \]
  由固定的 $h\in H$ 给出的映射 $x\mapsto hxh^{-1}$ 称为由 $h$ 带来的共轭，元素 $hxh^{-1}$ 称为 $x$ 的一个\textbf{共轭元}。

  若 $K\leq G$，则对任意 $h\in H$，都有
  \[
    hKh^{-1}\leq G,\qquad hKh^{-1}\cong K.
  \]
  因此 $H$ 也可以在 $G$ 的所有子群组成的集合上以共轭作用：
  \[
    (h,K)\longmapsto hKh^{-1}.
  \]
  子群 $hKh^{-1}$ 称为 $K$ 的一个\textbf{共轭子群}。
\end{enumerate}
\end{example}

\begin{purplenote}
在上述第 (3) 种作用中，为避免与第 (2) 种左平移的记号 $hx$ 混淆，通常直接写出 $hxh^{-1}$。
\end{purplenote}

% 整理：原稿 2.4.3(2) 的“x∈G”按作用的定义域修正为“x∈S”。
\begin{theorem}\label{thm:action-orbit-equivalence-stabilizer}
设群 $G$ 作用于集合 $S$，则：
\begin{enumerate}[label=(\arabic*)]
  \item 在 $S$ 上定义关系
  \[
    x\sim x'\quad\Longleftrightarrow\quad
    \text{存在 }g\in G\text{，使 }gx=x',
  \]
  则 $\sim$ 是等价关系。
  \item 对任意 $x\in S$，集合
  \[
    G_x=\{g\in G\mid gx=x\}
  \]
  是 $G$ 的一个子群。
\end{enumerate}
\end{theorem}

\begin{purplenote}
定理~\ref{thm:action-orbit-equivalence-stabilizer} 第 (1) 条中的等价类称为 $G$ 在 $S$ 上的\textbf{轨道}，$x\in S$ 所在的轨道记为
\[
  \bar x=\{gx\mid g\in G\}.
\]
$G_x$ 称为 $x$ 的\textbf{稳定子群}，也称\textbf{迷向群}。
\end{purplenote}

% 来源：微信图片_20261118134405_56_26.jpg，2.4.4--2.4.8 及紫笔说明、短证。
\begin{example}\label{ex:conjugacy-centralizer-normalizer}
若 $G$ 通过共轭作用于自身，则 $x\in G$ 的轨道
\[
  \bar x=\{gxg^{-1}\mid g\in G\}
\]
称为 $x$ 的\textbf{共轭类}。

若 $H\leq G$ 通过共轭作用于 $G$，则 $x$ 的稳定子群是
\[
  H_x=\{h\in H\mid hxh^{-1}=x\}
      =\{h\in H\mid hx=xh\}.
\]
它称为 $x$ 在 $H$ 中的\textbf{中心子}，记为 $C_H(x)$。当 $H=G$ 时，$C_G(x)$ 简称为 $x$ 的中心子。

若 $H$ 通过共轭作用于 $G$ 的所有子群组成的集合，则子群 $K\leq G$ 的稳定子群为
\[
  H_K=\{h\in H\mid hKh^{-1}=K\},
\]
称为 $K$ 在 $H$ 中的\textbf{正规子}，记为 $N_H(K)$。当 $H=G$ 时，$N_G(K)$ 简称为 $K$ 的正规子。
\end{example}

\begin{purplenote}
显然有
\[
  K\trianglelefteq N_G(K),\qquad
  \text{且有} K\trianglelefteq G\quad\Longleftrightarrow\quad N_G(K)=G.
\]
\end{purplenote}

\begin{theorem}\label{thm:orbit-stabilizer-cardinality}
设群 $G$ 作用于集合 $S$，则对每个 $x\in S$，其轨道的基数为
\[
  |\bar x|=[G:G_x].
\]
\end{theorem}

\begin{purpleproof}
对任意 $g,h\in G$，
\[
  gx=hx
  \quad\Longleftrightarrow\quad g^{-1}hx=x
  \quad\Longleftrightarrow\quad g^{-1}h\in G_x
  \quad\Longleftrightarrow\quad gG_x=hG_x.
\]
故 $gG_x\mapsto gx$ 给出了左陪集集合与轨道 $\bar x$ 之间的双射。
\end{purpleproof}

% 整理：类方程对全部共轭类求和，明确原稿“不同共轭类”是完整代表系。
\begin{corollary}\label{cor:conjugacy-class-equation}
设 $G$ 是有限群，$K\leq G$，则：
\begin{enumerate}[label=(\arabic*)]
  \item 对任意 $x\in G$，其共轭类中的元素个数为
  \[
    |\bar x|=[G:C_G(x)],\qquad [G:C_G(x)]\mid |G|.
  \]
  \item 若 $\bar x_1,\ldots,\bar x_n$ 是 $G$ 的全部不同共轭类，则
  \[
    |G|=\sum_{i=1}^{n}[G:C_G(x_i)].
  \]
  \item $G$ 中与 $K$ 共轭的子群的数目为
  \[
    [G:N_G(K)],\qquad [G:N_G(K)]\mid |G|.
  \]
\end{enumerate}
\end{corollary}

\begin{purplenote}
推论~\ref{cor:conjugacy-class-equation} 第 (2) 条的等式称为有限群 $G$ 的\textbf{类方程}。
\end{purplenote}

\begin{theorem}\label{thm:action-induced-permutation-homomorphism}
若群 $G$ 作用于集合 $S$，则这个作用诱导出一个群同态
\[
  \tau:G\longrightarrow A(S),
\]
其中 $A(S)$ 是 $S$ 上的所有置换组成的群。
\end{theorem}

% 整理：按原稿给出的 tau_g 定义，写明其逆映射及由作用公理得到的复合等式。
\begin{purpleproof}
对每个 $g\in G$，定义
\[
  \tau_g:S\longrightarrow S,\qquad \tau_g(x)=gx.
\]
由于 $\tau_{g^{-1}}$ 是 $\tau_g$ 的逆映射，$\tau_g\in A(S)$；又由作用公理，
\[
  \tau_{gh}(x)=(gh)x=g(hx)=(\tau_g\circ\tau_h)(x).
\]
故 $g\mapsto\tau_g$ 就是所需的群同态。
\end{purpleproof}

\begin{purplenote}
需注意这里并没有具体指定$G$在$S$上的作用映射，因此所诱导的群同态也不唯一。后续研究中，通常只需要知道存在这样的群同态即可。
\end{purplenote}

\begin{corollary}[Cayley]\label{cor:cayley-permutation-embedding}
对任意群 $G$，存在群单同态
\[
  G\longrightarrow A(G).
\]
% 根据 2026-11-19 的讨论恢复原稿“一些置换组成的群”的含义，明确为 A(G) 的子群。
因此任意群都同构于一个由某些置换组成的群，即 $A(G)$ 的一个子群。特别地，每个有限群 $G$ 都同构于 $S_n$ 的一个子群，其中 $n=|G|$。
\end{corollary}

\begin{purpleproof}
令 $G$ 在自身上以左平移作用，应用定理~\ref{thm:action-induced-permutation-homomorphism}。对任意 $x\in G$，由消去律有 $gx=x$ 当且仅当 $g=e$，故所诱导的群同态是单射。
\end{purpleproof}

% 来源：微信图片_20261118134408_57_26.jpg，2.4.9--2.4.12 及紫笔说明、证明。
\begin{purplenote}
对于群 $G$，$\operatorname{Aut}G$ 表示从 $G$ 到自身的所有群同构，即所有自同构组成的集合。以映射复合作为二元运算，$\operatorname{Aut}G$ 是一个群。
\end{purplenote}

\begin{corollary}\label{cor:inner-automorphisms-center-kernel}
设 $G$ 是群，则：
\begin{enumerate}[label=(\arabic*)]
  \item 对每个 $g\in G$，由 $g$ 诱导的共轭映射
  \[
    \tau_g:G\longrightarrow G,\qquad \tau_g(x)=gxg^{-1},
  \]
  是 $G$ 的自同构。
  \item 存在群同态
  \[
    \tau:G\longrightarrow\operatorname{Aut}G,\qquad g\longmapsto\tau_g,
  \]
  其核为
  \[
    C(G)=\{g\in G\mid gx=xg\text{ 对所有 }x\in G\text{ 成立}\}.
  \]
\end{enumerate}
\end{corollary}

\begin{purplenote}
推论~\ref{cor:inner-automorphisms-center-kernel} 第 (1) 条中的 $\tau_g$ 称为由 $g$ 诱导的\textbf{内同构}，正规子群 $C(G)=\ker\tau$ 称为 $G$ 的\textbf{中心}。

元素 $g\in C(G)$，当且仅当它的共轭类中只有 $g$ 一个元素。因此，当 $G$ 有限时，对 $x\in C(G)$ 有 $[G:C_G(x)]=1$，类方程可改写为
\[
  |G|=|C(G)|+\sum_{i=1}^{m}[G:C_G(x_i)],
\]
其中 $x_1,\ldots,x_m\in G\setminus C(G)$ 代表全部不同的非中心共轭类，各 $[G:C_G(x_i)]>1$。
\end{purplenote}

\begin{proposition}\label{prop:coset-action-kernel-contained-subgroup}
设 $G$ 是群，$H\leq G$。令 $G$ 在 $H$ 的所有左陪集组成的集合 $S$ 上以左平移作用，则所诱导的群同态
\[
  \tau:G\longrightarrow A(S)
\]
的核包含于 $H$。
\end{proposition}

\begin{purpleproof}
若 $g\in\ker\tau$，则 $g$ 固定左陪集 $H$，即 $gH=H$，故 $g\in H$。
\end{purpleproof}

\begin{corollary}\label{cor:core-free-subgroup-permutation-embedding}
设 $H\leq G$，$[G:H]=n$。若 $G$ 的非平凡正规子群均不包含于 $H$，则 $G$ 同构于 $S_n$ 的一个子群。
\end{corollary}

\begin{purpleproof}
直接应用命题~\ref{prop:coset-action-kernel-contained-subgroup}：所得同态的核是包含于 $H$ 的正规子群，由假设只能为 $\{e\}$。
\end{purpleproof}

\begin{corollary}\label{cor:smallest-prime-index-subgroup-normal}
设 $G$ 是有限群，$H\leq G$，且 $[G:H]=p$，其中 $p$ 是整除 $|G|$ 的最小质数，则 $H\trianglelefteq G$。
\end{corollary}

\begin{purpleproof}
令 $G$ 在 $H$ 的所有左陪集组成的集合 $S$ 上以左平移作用，记诱导的群同态为 $\varphi:G\to A(S)$，令 $K=\ker\varphi$。由命题~\ref{prop:coset-action-kernel-contained-subgroup}，$K\leq H$，只需证明 $K=H$。

因为 $|S|=p$，由\hyperref[cor:first-isomorphism]{第一同构定理}及\hyperref[cor:lagrange]{拉格朗日定理}，有
\[
  G/K\cong\operatorname{im}\varphi\leq A(S)\cong S_p,
  \qquad [G:K]\mid p!.
\]
又由\hyperref[thm:index-multiplication]{指数乘法公式}，
\[
  [G:K]=[G:H][H:K]=p[H:K],
  \qquad |G|=p[H:K]|K|.
\]
由 $p$ 的最小性，$[H:K]$ 的每个质因子都不小于 $p$；结合 $[G:K]\mid p!$，知这些质因子只能是 $p$，故 $[H:K]=p^a$，其中 $a\geq0$。但 $p!$ 中只含一个因子 $p$，而
\[
  p^{a+1}=[G:K]\mid p!,
\]
故 $a=0$，从而 $[H:K]=1$，即 $H=K\trianglelefteq G$。
\end{purpleproof}

% 来源：微信图片_20261118134410_58_26.jpg，2.4.13--2.4.15 及紫笔证明提示。
下面是本节的一些习题结论。

\begin{proposition}\label{prop:normalizer-centralizer-subgroup-properties}
设 $H\leq G$，则：
\begin{enumerate}[label=(\arabic*)]
  \item 若 $K\leq G$ 且 $H\trianglelefteq K$，则 $K\leq N_G(H)$。
  \item 定义
  \[
    C_G(H)=\{g\in G\mid gh=hg\text{ 对所有 }h\in H\text{ 成立}\},
  \]
  则 $C_G(H)\trianglelefteq N_G(H)$。
\end{enumerate}
\end{proposition}

% 字迹放大核对：原稿 2.4.14 是“阶为 p”的子群，不是“指数为 p”。
\begin{proposition}\label{prop:prime-order-subgroup-small-cofactor-normal}
若 $|G|=pn$，$p>n$ 且 $p$ 为质数，$H$ 是 $G$ 的一个阶为 $p$ 的子群，则 $H\trianglelefteq G$。
\end{proposition}

\begin{purplenote}
证明思路与推论~\ref{cor:smallest-prime-index-subgroup-normal} 的证明类似：仍考察陪集作用的核 $K$，证明 $H=K$ 即可。
\end{purplenote}

\begin{proposition}\label{prop:unique-normal-subgroup-prime-power-index}
若 $p,q$ 都是质数，$|G|=p^nq$ 且 $p>q$，则 $G$ 包含唯一一个指数为 $q$ 的正规子群。
\end{proposition}

\begin{purplenote}
% 整理引用：第 2.5 节已接入，为原稿“下一节”补充可点击的章节引用。
该结论要用到第~\ref{sec:sylow-theorems} 节的 Sylow 定理和推论~\ref{cor:smallest-prime-index-subgroup-normal}。
\end{purplenote}


% 新增：第 2.5 节独立维护，整理 Cauchy 定理、Sylow 定理及其应用。
% !TeX root = ../main.tex
\subsection{Sylow 定理}\label{sec:sylow-theorems}

% 来源目录：../../NotePhotos/BaiduSyncdisk/algebra/2.5/（相对于本科目目录）。
% 共 5 张照片，文件名编号 61--65 不是原稿页码；覆盖原稿 2.5.1--2.5.17。
% 2.5.8、2.5.11、2.5.12 的另页证明移至对应条目后；无编号补充性质不占用原编号。

% 来源：微信图片_20261119110950_61_26.jpg，节首、2.5.1--2.5.5 及紫笔证明。
% 整理：原稿紫笔补写 G 为有限群，已并入节首的明确假设。
本节的一个动机是：设 $G$ 为有限群，若 $m\mid |G|$，是否一定存在 $G$ 的一个阶为 $m$ 的子群？对于第~\ref{sec:finitely-generated-abelian-groups} 节讨论的有限阿贝尔群，推论~\ref{cor:finite-abelian-subgroups-all-divisors} 给出了肯定回答；对于非阿贝尔群则未必成立。Sylow 定理讨论的是 $m=p^a$，$p$ 为质数的情形。

\begin{lemma}\label{lem:p-group-action-fixed-point-congruence}
设群 $H$ 的阶为 $p^n$，其中 $p$ 为质数。若 $H$ 作用于有限集合 $S$，记
\[
  S_0=\{x\in S\mid hx=x\text{ 对所有 }h\in H\text{ 成立}\},
\]
则
\[
  |S|\equiv |S_0|\pmod p.
\]
\end{lemma}

% 整理：原稿将求和指标简写为 x 不属于 S_0；这里明确每个非单点轨道只取一个代表，避免重复计数。
\begin{purpleproof}
若 $x\in S_0$，则其轨道 $\bar x=\{x\}$。在其余各轨道中分别取一个代表 $x_1,\ldots,x_r$，由定理~\ref{thm:orbit-stabilizer-cardinality}，有
\[
  |S|=|S_0|+\sum_{i=1}^{r}[H:H_{x_i}].
\]
由于 $|H|=p^n$，每个非单点轨道的基数 $[H:H_{x_i}]$ 都是大于 $1$ 的 $p$ 的幂，因而均被 $p$ 整除，结论成立。
\end{purpleproof}

\begin{theorem}[Cauchy]\label{thm:cauchy-prime-order-element}
若 $G$ 是有限群，且质数 $p$ 整除 $|G|$，则 $G$ 中存在一个阶为 $p$ 的元素。
\end{theorem}

% 整理：原稿明确提示须检验循环移位作用良定义；此处写出非交换情形的乘积检验。
\begin{purpleproof}
证明的思路是构造适当的群作用。考虑
\[
  S=\{(a_1,\ldots,a_p)\in G^p\mid a_1\cdots a_p=e\}.
\]
前 $p-1$ 个分量任意选定后，最后一个分量由乘积条件唯一确定，故
\[
  |S|=|G|^{p-1},\qquad p\mid |S|.
\]
令 $\mathbb Z_p$ 通过循环移位作用于 $S$：对 $0\leq k<p$，定义
\[
  \bar k\cdot(a_1,\ldots,a_p)
    =(a_{k+1},\ldots,a_p,a_1,\ldots,a_k).
\]
该作用是良定义的：一次循环移位后，乘积为
\[
  a_2\cdots a_pa_1=a_1^{-1}(a_1\cdots a_p)a_1=e,
\]
所以仍落在 $S$ 中；循环移位 $p$ 次为恒等映射，且移位次数相加对应于映射复合。

由引理~\ref{lem:p-group-action-fixed-point-congruence}，公共不动点集 $S_0$ 满足
\[
  |S_0|\equiv |S|\equiv0\pmod p.
\]
一个元组属于 $S_0$，当且仅当它的所有分量相同。因此
\[
  S_0=\{(a,\ldots,a)\mid a\in G,\ a^p=e\}.
\]
又 $(e,\ldots,e)\in S_0$，故 $|S_0|\geq p>1$。于是存在 $a\ne e$，使 $a^p=e$。因为 $p$ 为质数，$a$ 的阶恰为 $p$。
\end{purpleproof}

\begin{purplenote}
下面再次明确 $p$-群的定义。
\end{purplenote}

\begin{definition}[$p$-群]\label{def:p-group-p-subgroup}
给定质数 $p$，若群 $G$ 中所有元素的阶都是 $p$ 的非负整数次幂，则称 $G$ 为\textbf{$p$-群}。若 $H\leq G$ 且 $H$ 是 $p$-群，则称 $H$ 是 $G$ 的一个\textbf{$p$-子群}。

特别地，平凡子群 $\{e\}$ 对每个质数 $p$ 都是 $G$ 的一个 $p$-子群。
\end{definition}

\begin{corollary}\label{cor:finite-p-group-order-characterization}
有限群 $G$ 是 $p$-群，当且仅当 $|G|$ 是 $p$ 的非负整数次幂。
\end{corollary}

\begin{purpleproof}
直接应用定理~\ref{thm:cauchy-prime-order-element} 即可。
\end{purpleproof}

\begin{corollary}\label{cor:nontrivial-finite-p-group-center}
非平凡有限 $p$-群 $G$ 的中心
\[
  C(G)=\{g\in G\mid gx=xg\text{ 对所有 }x\in G\text{ 成立}\}
\]
至少包含两个元素。
\end{corollary}

\begin{purpleproof}
由\hyperref[cor:conjugacy-class-equation]{类方程}，
\[
  |G|=|C(G)|+\sum_{i=1}^{r}[G:C_G(x_i)],
\]
其中 $x_i$ 代表全部不同的非中心共轭类。每个 $[G:C_G(x_i)]>1$ 都是 $p$ 的幂，故 $p\mid |C(G)|$。又 $e\in C(G)$，所以 $|C(G)|\geq p\geq2$。
\end{purpleproof}

% 来源：微信图片_20261119110952_62_26.jpg，2.5.6--2.5.10 及紫笔说明。
\begin{lemma}\label{lem:p-subgroup-normalizer-index-congruence}
若 $H$ 是有限群 $G$ 的一个 $p$-子群，则
\[
  [N_G(H):H]\equiv [G:H]\pmod p.
\]
\end{lemma}

\begin{purpleproof}
令 $S$ 为 $H$ 在 $G$ 中的所有左陪集组成的集合，并令 $H$ 在 $S$ 上以左平移作用。对 $xH\in S$，有
\[
  \begin{aligned}
    xH\in S_0
      &\Longleftrightarrow hxH=xH\quad\text{对所有 }h\in H\\
      &\Longleftrightarrow x^{-1}Hx\leq H\\
      &\Longleftrightarrow x^{-1}Hx=H\\
      &\Longleftrightarrow x\in N_G(H).
  \end{aligned}
\]
这里由包含关系得到相等，用到了 $H$ 有限以及 $|x^{-1}Hx|=|H|$。故
\[
  |S_0|=[N_G(H):H].
\]
再由引理~\ref{lem:p-group-action-fixed-point-congruence}，即得
\[
  [G:H]=|S|\equiv |S_0|=[N_G(H):H]\pmod p.
\]
\end{purpleproof}

\begin{corollary}\label{cor:p-subgroup-proper-normalizer}
若 $H$ 是有限群 $G$ 的一个 $p$-子群，且 $p\mid[G:H]$，则 $N_G(H)\ne H$。
\end{corollary}

\begin{purpleproof}
由引理~\ref{lem:p-subgroup-normalizer-index-congruence}，有 $p\mid[N_G(H):H]$，故 $[N_G(H):H]>1$。
\end{purpleproof}

\begin{theorem}[第一 Sylow 定理]\label{thm:first-sylow}
设群 $G$ 的阶为 $p^nm$，其中 $p$ 为质数，$n\geq1$ 且 $\gcd(p,m)=1$。则对每个 $1\leq i\leq n$，$G$ 都包含阶为 $p^i$ 的子群。而且，$G$ 中任意一个阶为 $p^i$ 的子群（$1\leq i<n$），都是 $G$ 中某个阶为 $p^{i+1}$ 的子群的正规子群。
\end{theorem}

% 来源：微信图片_20261119110954_63_26.jpg，2.5.8 的另页证明。
% 修订：原稿首句误引“2.5.8”自身；所用的阶为 p 元素的存在性实际是 Cauchy 定理 2.5.2。
\begin{purpleproof}
首先由定理~\ref{thm:cauchy-prime-order-element}，存在 $a\in G$ 的阶为 $p$。于是 $\langle a\rangle$ 是阶为 $p$ 的子群。

现设 $H\leq G$ 的阶为 $p^i$，其中 $1\leq i<n$。由推论~\ref{cor:finite-p-group-order-characterization}，$H$ 是 $p$-子群；又因 $i<n$，有 $p\mid[G:H]$。结合引理~\ref{lem:p-subgroup-normalizer-index-congruence}，得到
\[
  p\mid[N_G(H):H].
\]
由于 $H\trianglelefteq N_G(H)$，商群 $N_G(H)/H$ 的阶被 $p$ 整除。再由~\ref{thm:cauchy-prime-order-element} 定理，其中存在一个阶为 $p$ 的子群 $\overline H_1$。由推论~\ref{cor:quotient-subgroup-correspondence}，它对应于一个子群 $H_1$，满足
\[
  H\leq H_1\leq N_G(H),\qquad H_1/H=\overline H_1.
\]
且有
\[
  |H_1|=|H_1/H|\,|H|=p\cdot p^i=p^{i+1}.
\]
这证明了所述的正规扩张性质；从阶为 $p$ 的子群出发逐次应用，即得到各个阶为 $p^i$ 的子群，完成归纳。
\end{purpleproof}

% 整理：原稿“最大 p-群”按随后的包含关系定义写为“极大 p-子群”。
\begin{definition}[Sylow $p$-子群]\label{def:sylow-p-subgroup}
设 $p$ 为质数，$P\leq G$。若 $P$ 是 $G$ 的一个按包含关系极大的 $p$-子群，即 $P$ 是 $p$-群，且只要 $H$ 是 $G$ 的 $p$-子群并满足
\[
  P\leq H\leq G,
\]
就有 $P=H$，则称 $P$ 为 $G$ 的一个\textbf{Sylow $p$-子群}。
\end{definition}

\begin{purplenote}
Sylow $p$-子群总是存在，尽管可能是平凡子群，且Sylow $p$-子群可能不唯一。每个 $p$-子群都包含在某个 Sylow $p$-子群中；对于无限群，这一存在性需要使用 Zorn 引理来证明。

第一 Sylow 定理进一步说明：对于任意有限群 $G$ 及任意整除 $|G|$ 的质数 $p$，$G$ 都有非平凡的 $p$-子群。
\end{purplenote}

\begin{corollary}\label{cor:sylow-order-conjugates-unique-normal}
设 $|G|=p^nm$，其中 $p$ 为质数，$n\geq1$ 且 $\gcd(p,m)=1$，$H$ 是 $G$ 的一个 $p$-子群，则：
\begin{enumerate}[label=(\arabic*)]
  \item $H$ 是 Sylow $p$-子群，当且仅当 $|H|=p^n$。
  \item 任意一个 Sylow $p$-子群的共轭子群也是 Sylow $p$-子群。
  \item 若 $G$ 有且仅有一个 Sylow $p$-子群 $P$，则 $P\trianglelefteq G$。
\end{enumerate}
\end{corollary}

% 来源：微信图片_20261119110957_64_26.jpg，2.5.11--2.5.17、紫笔说明及无编号补充性质。
\begin{theorem}[第二 Sylow 定理]\label{thm:second-sylow}
若 $H$ 是有限群 $G$ 的一个 $p$-子群，$P$ 是 $G$ 的任意一个 Sylow $p$-子群，则存在 $x\in G$，使得
\[
  H\leq xPx^{-1}.
\]
特别地，$G$ 中任意两个 Sylow $p$-子群都是共轭的。
\end{theorem}

\begin{purplenote}
这里关于 Sylow $p$-子群共轭性的结论，可视为推论~\ref{cor:sylow-order-conjugates-unique-normal} 第 (2) 条的逆向结论。
\end{purplenote}

% 来源：微信图片_20261119110959_65_26.jpg，2.5.11 的另页证明。
\begin{purpleproof}
令 $S$ 为 $P$ 在 $G$ 中的所有左陪集组成的集合，并令 $H$ 在 $S$ 上以左平移作用。注意，对 $x\in G$，
\[
  \begin{aligned}
    H\leq xPx^{-1}
      &\Longleftrightarrow x^{-1}Hx\leq P\\
      &\Longleftrightarrow x^{-1}hx\in P\quad\text{对所有 }h\in H\\
      &\Longleftrightarrow hxP=xP\quad\text{对所有 }h\in H\\
      &\Longleftrightarrow xP\in S_0.
  \end{aligned}
\]
由于 $P$ 是 Sylow $p$-子群，$p\nmid[G:P]=|S|$。再由引理~\ref{lem:p-group-action-fixed-point-congruence}，
\[
  |S_0|\equiv|S|\not\equiv0\pmod p,
\]
故 $S_0\ne\varnothing$。取 $xP\in S_0$，即得 $H\leq xPx^{-1}$。

若 $H$ 本身也是 Sylow $p$-子群，则 $|H|=|P|=|xPx^{-1}|$，故上述包含关系必为相等。
\end{purpleproof}

\begin{theorem}[第三 Sylow 定理]\label{thm:third-sylow}
设 $G$ 是有限群，$p$ 为质数。若 $k$ 是 $G$ 的所有 Sylow $p$-子群的数目，则
\[
  k\mid |G|,\qquad k\equiv1\pmod p.
\]
\end{theorem}

% 来源：微信图片_20261119110959_65_26.jpg，2.5.12 的另页证明。
\begin{purpleproof}
固定 $G$ 的一个 Sylow $p$-子群 $P$，令 $S$ 为 $G$ 的所有 Sylow $p$-子群组成的集合。由定理~\ref{thm:second-sylow}，
\[
  S=\{xPx^{-1}\mid x\in G\}.
\]
因此 $S$ 是共轭作用下 $P$ 的轨道，由推论~\ref{cor:conjugacy-class-equation} 第 (3) 条，
\[
  k=|S|=[G:N_G(P)]\mid |G|.
\]

现令 $P$ 通过共轭作用于 $S$。若 $Q\in S_0$，则对所有 $x\in P$ 都有 $xQx^{-1}=Q$，即
\[
  P\leq N_G(Q).
\]
$P,Q$ 均是 $G$ 的 Sylow $p$-子群，故也是 $N_G(Q)$ 的 Sylow $p$-子群。由第二 Sylow 定理，它们在 $N_G(Q)$ 中共轭；但 $Q\trianglelefteq N_G(Q)$，因此只能有 $P=Q$。反过来，显然 $P\in S_0$，所以
\[
  S_0=\{P\}.
\]
再由引理~\ref{lem:p-group-action-fixed-point-congruence}，得到
\[
  k=|S|\equiv|S_0|=1\pmod p.
\]
\end{purpleproof}

\begin{theorem}\label{thm:sylow-normalizer-self-normalizing}
若 $P$ 是有限群 $G$ 的一个 Sylow $p$-子群，则
\[
  N_G(N_G(P))=N_G(P).
\]
\end{theorem}

\begin{purpleproof}
显然 $N_G(P)\leq N_G(N_G(P))$。反过来，若 $x\in N_G(N_G(P))$，则
\[
  xN_G(P)x^{-1}=N_G(P).
\]
由于 $P\leq N_G(P)$，有 $xPx^{-1}\leq N_G(P)$。$P$ 与 $xPx^{-1}$ 都是 $N_G(P)$ 的 Sylow $p$-子群，由定理~\ref{thm:second-sylow}，存在 $x_0\in N_G(P)$，使得
\[
  xPx^{-1}=x_0Px_0^{-1}.
\]
而 $P\trianglelefteq N_G(P)$，故右边等于 $P$。因此 $xPx^{-1}=P$，即 $x\in N_G(P)$，结论成立。
\end{purpleproof}

\par\textbf{补充性质（原稿无编号）。}
若 $G$ 是有限 $p$-群，$H\trianglelefteq G$ 且 $H\ne\{e\}$，则
\[
  H\cap C(G)\ne\{e\}.
\]

下面列出本节的一些习题性质。

\begin{proposition}\label{prop:p-group-normal-subgroups-all-orders}
设 $p$ 为质数，$|G|=p^n$。则对每个 $0\leq k\leq n$，$G$ 都有一个阶为 $p^k$ 的正规子群。
\end{proposition}

\begin{purplenote}
证明提示：考虑
\[
  S_k=\{H\leq G\mid |H|=p^k\}
\]
及 $G$ 在 $S_k$ 上的共轭作用，并考察相应的公共不动点集 $(S_k)_0$。
\end{purplenote}

\begin{proposition}\label{prop:infinite-p-group-finite-subgroup-orders}
设 $p$ 为质数，$G$ 为无限群且为 $p$-群。则下列两种情形至少有一种成立：
\begin{enumerate}[label=(\arabic*)]
  \item 对每个 $n\geq1$，$G$ 都有一个阶为 $p^n$ 的子群。
  \item 存在 $m\in\mathbb N^*$，使 $G$ 的所有有限子群的阶都不超过 $p^m$。
\end{enumerate}
\end{proposition}

\begin{proposition}\label{prop:normal-sylow-orders-12-28-56-200}
任意阶为 $12,28,56,200$ 的群都有一个非平凡的真正规 Sylow 子群，从而不是简单群。
\end{proposition}

% 整理：不同质数阶子群的交为单位子群，原稿计数时的“完全不交”据此表述。
\begin{purpleproof}
记 Sylow $p$-子群的数目为 $n_p$。由第三 Sylow 定理，$n_p\mid |G|$ 且 $n_p\equiv1\pmod p$，故 $n_p$ 与 $p$ 互素，并整除 $|G|$ 中去掉 $p$ 的幂后的部分。

若 $|G|=28=2^2\cdot7$，则 $n_7\mid4$ 且 $n_7\equiv1\pmod7$，故 $n_7=1$。若 $|G|=200=2^3\cdot5^2$，则 $n_5\mid8$ 且 $n_5\equiv1\pmod5$，故 $n_5=1$。由推论~\ref{cor:sylow-order-conjugates-unique-normal} 第 (3) 条，相应的 Sylow 子群都是正规的。

若 $|G|=12=2^2\cdot3$，则 $n_3=1$ 或 $4$。前者已经给出正规 Sylow 子群；若 $n_3=4$，由于不同的阶为 $3$ 的子群只交于单位元，它们共包含
\[
  4(3-1)=8
\]
个非单位元素，剩下 $12-8=4$ 个元素（包括单位元）。任何 Sylow $2$-子群都有 $4$ 个元素，且其非单位元素不可能属于上述阶为 $3$ 的子群，因而只能恰由这 $4$ 个剩余元素组成。所以 Sylow $2$-子群唯一，进而正规。

若 $|G|=56=2^3\cdot7$，同理 $n_7=1$ 或 $8$。前者给出正规 Sylow 子群；若 $n_7=8$，这些阶为 $7$ 的子群共有 $8(7-1)=48$ 个非单位元素。剩下的 $56-48=8$ 个元素只能组成唯一的 Sylow $2$-子群，故该子群正规。

在上述各情形中，所得正规 Sylow 子群均非平凡且为真子群，因此 $G$ 不是简单群。
\end{purpleproof}

\begin{proposition}\label{prop:finite-group-normal-sylow-direct-product}
若 $G$ 是有限群，且对每个质数 $p$，$G$ 的 Sylow $p$-子群都是正规的，则 $G$ 是其各 Sylow 子群的内直积。
\end{proposition}


\end{document}
